\documentclass[11pt,oneside,openany]{book}

% BEGIN ORIGINAL INPUT: preamble.tex
\usepackage[T1]{fontenc}
\usepackage{lmodern}
\usepackage{amsmath,amssymb,amsthm,mathtools}
\usepackage[a4paper,margin=25mm,headheight=15pt]{geometry}
\usepackage{microtype,enumitem,xcolor,longtable,booktabs,array,needspace}
\usepackage[colorlinks=true,linkcolor=blue!40!black,urlcolor=blue!40!black]{hyperref}
\usepackage[nameinlink,noabbrev]{cleveref}
\hypersetup{pdftitle={Multivariable Analysis and Vector Calculus -- Unified Undergraduate Notes},pdfauthor={MathGenome mathematical preparation series}}
\usepackage{fancyhdr}
\pagestyle{fancy}\fancyhf{}
\fancyhead[L]{\small MathGenome: Multivariable Analysis}
\fancyhead[R]{\small\thepage}
\renewcommand{\headrulewidth}{0.3pt}
\setlength{\parindent}{0pt}
\setlength{\parskip}{5pt}
\setlength{\emergencystretch}{3em}
\allowdisplaybreaks[2]\raggedbottom
\setcounter{tocdepth}{1}
\newtheorem{theorem}{Theorem}[chapter]
\newtheorem{lemma}[theorem]{Lemma}
\newtheorem{proposition}[theorem]{Proposition}
\newtheorem{corollary}[theorem]{Corollary}
\theoremstyle{definition}
\newtheorem{definition}[theorem]{Definition}
\newtheorem{example}[theorem]{Example}
\theoremstyle{remark}
\newtheorem{remark}[theorem]{Remark}
\newcounter{foundation}
\renewcommand{\thefoundation}{F\arabic{foundation}}
\crefname{foundation}{foundation}{foundations}
\Crefname{foundation}{Foundation}{Foundations}
\newcommand{\N}{\mathbb N}
\newcommand{\Z}{\mathbb Z}
\newcommand{\Q}{\mathbb Q}
\newcommand{\R}{\mathbb R}
\newcommand{\C}{\mathbb C}
\newcommand{\eps}{\varepsilon}
\newcommand{\abs}[1]{\lvert#1\rvert}
\newcommand{\norm}[1]{\lVert#1\rVert}
\newcommand{\ip}[2]{\langle#1,#2\rangle}
\newcommand{\dd}{\,\mathrm d}
\newcommand{\D}{\mathrm D}
\DeclareMathOperator{\id}{id}
\DeclareMathOperator{\supp}{supp}
\DeclareMathOperator{\rank}{rank}
\DeclareMathOperator{\tr}{tr}
\DeclareMathOperator{\diag}{diag}
\DeclareMathOperator{\Span}{span}
\DeclareMathOperator{\Hom}{Hom}
\DeclareMathOperator{\Alt}{Alt}
\DeclareMathOperator{\End}{End}
\DeclareMathOperator{\vol}{vol}
\DeclareMathOperator{\sgn}{sgn}
\DeclareMathOperator{\dist}{dist}
\DeclareMathOperator{\diam}{diam}
\DeclareMathOperator{\GL}{GL}
\DeclareMathOperator{\Int}{Int}
\DeclareMathOperator{\cl}{cl}
\newcommand{\family}[1]{}
\newcommand{\source}[1]{\par\begingroup\small\sffamily\color{black!65}Source comparison: #1\par\endgroup}

% END ORIGINAL INPUT: preamble.tex

\begin{document}
\frontmatter

% BEGIN ORIGINAL INPUT: frontmatter.tex
\hypersetup{pageanchor=false}
\begin{titlepage}
\centering
\vspace*{24mm}
{\Large MathGenome Mathematical Preparation Series\par}
\vspace{18mm}
{\Huge\bfseries Multivariable Analysis\par}
\vspace{5mm}
{\LARGE and Vector Calculus\par}
\vspace{14mm}
{\Large Unified Undergraduate Notes\par}
\vspace{8mm}
{\large D04: fifteen connected modules\par}
\vfill
{\large Original mathematical synthesis with explicit proofs,\par
source comparisons and declared foundations\par}
\vspace{10mm}
{\large Manuscript prepared 1 November 2026\par}
\vspace{6mm}
{\small This is a mathematical preparation artifact.\par
It is not a Lean verification certificate or a new MathGenome release.\par}
\end{titlepage}
\hypersetup{pageanchor=true}

\chapter*{Purpose and reading policy}
\addcontentsline{toc}{chapter}{Purpose and reading policy}
These notes assemble the undergraduate theory of finite-dimensional real
multivariable analysis into a single consistent language. The central sequence
is total differentiation, local solvability, multiple integration and change of
variables, followed by curves, surfaces, differential forms and Stokes' theorem.
The Euclidean manifold material is included to close the foundations used by
integration; it is not a substitute for a general course on abstract manifolds.

The scope was organized into fifteen modules and 158 knowledge families before
writing. A family is a coverage unit, not a quota of theorems or future Lean
goals. Several related formulations may share one general result and its proved
specializations. Conversely a single family may need several independent
results, counterexamples and representation bridges. The machine-readable
coverage register associates these families with actual manuscript labels.

This text is an original synthesis. Source references identify comparison
passages, conventions and proof checks; they do not assert that every exercise
or every theorem in ten comparison books belongs to D04. The scalar and
finite-dimensional algebra foundations are declared separately. Every new
multivariable assertion is to be justified here, rather than hidden under a
reference to a more advanced change-of-variables or Stokes theorem.

\section*{What fidelity requires}
The required mathematical content consists of the stated results, their
hypotheses and conclusions, the definitions they concern, and the bridges
needed to identify different presentations of those objects. Equivalent proof
routes are permitted. An intermediate assertion used only in a source's chosen
proof does not automatically become a separate required result. A named result
actually stated in these notes, however, is not excused from proof by that rule.
Examples and counterexamples must satisfy their stated assumptions.

Finite regularity, injectivity, support, integrability, dimension and orientation
are mathematical hypotheses. They are not interchangeable implementation
choices. A vector-valued mean-value estimate is distinguished from a scalar
mean-value equality. Ordinary Jacobians use absolute determinants; oriented
forms use signed determinants. Multiple Riemann integrability does not assert
ordinary integrability of every slice. The relevant qualifications are retained
where the theorems are used.

The five bounded extensions are a finite-regularity Morse normal form, moving
domains given by a regular family of diffeomorphisms, vector potentials,
punctured-domain obstructions, and curvilinear-coordinate operators. Full
measure-theoretic integration, abstract manifold theory, general partial
differential equations, numerical analysis and general Banach-space calculus
are separate directions. No claim of covering all those subjects is made.

\section*{Notation and order of use}
Vectors are columns and $\D f(a)$ is a linear map, represented by the matrix
$Jf(a)_{ij}=\partial_j f_i(a)$. Thus $J(g\circ f)(a)=Jg(f(a))Jf(a)$.
The Euclidean metric identifies $\dd f$ with the gradient; it does not identify
every differential form with a vector without degree and orientation data.
Coordinate alternating forms are normalized by
$\dd x_{i_1}\wedge\cdots\wedge\dd x_{i_k}(e_{i_1},\ldots,e_{i_k})=1$ for increasing
indices. Boundaries use the outward-normal-first orientation.

Forward references to general Stokes in the classical vector-calculus chapter
are deliberate: the general theorem is proved later from local calculations,
partitions and substitution, independently of its classical applications.
Neither the surface chapter nor the vector identities are used to assume
Stokes in that proof. The source register records convention mismatches and
corrections that would otherwise be easy to miss.

The final preparation records distinguish mathematical editorial verification
from future machine verification. Existing Lean declarations listed in the
foundation interface are source-checked reuse candidates in the fixed library.
New notation, domains and types in this manuscript still need explicit Lean
bindings when a formalization project is authorized.

% END ORIGINAL INPUT: frontmatter.tex

\tableofcontents

% BEGIN ORIGINAL INPUT: foundations.tex
\chapter*{Declared foundations and reuse boundary}
\addcontentsline{toc}{chapter}{Declared foundations and reuse boundary}
These notes are self-contained relative to the interfaces below. Re-proving all
scalar analysis and finite-dimensional algebra would obscure the new subject;
silently importing multivariable analytic conclusions would be circular.
The interfaces are mathematical assumptions backed by the earlier directions
and the fixed library, not unproved D04 lemmas disguised as references.

\section*{F1. Scalar universe}\refstepcounter{foundation}\label{base:real}
We use classical sets, functions, finite sums and products, induction and
countable choices of witnesses, a complete ordered field $\R$, and
$\C=\R+i\R$ as a two-dimensional real vector space. Completeness means the
least-upper-bound property. Natural numbers include zero. An empty product is
one and the determinant of the $0\times0$ matrix is one. Topological and analytic
claims about $\R^n$ are proved in the chapters, not included merely by writing
the symbol $\R^n$.

\section*{F2. Finite-dimensional algebra}\refstepcounter{foundation}\label{base:linear}
The algebraic interface consists of vector spaces and linear maps over $\R$ or
$\C$, bases, dimension, rank--nullity, quotient spaces, duals, multilinear and
alternating maps, matrices and their change-of-basis rules. Determinants are
alternating multilinear functions of columns, normalized by $\det I=1$, with
the product, cofactor and transpose identities and Cauchy--Binet. A real inner
product gives Cauchy--Schwarz and Gram--Schmidt. A real symmetric matrix has an
orthonormal eigenbasis; its quadratic form has invariant positive, negative and
zero indices (Sylvester inertia). In particular a positive-definite Gram matrix
has positive determinant. All these are algebraic finite-dimensional facts.

The actual public library modules are
\nolinkurl{MathGenome.LinearAlgebra.Determinants},
\nolinkurl{MathGenome.LinearAlgebra.InnerProducts},
\nolinkurl{MathGenome.LinearAlgebra.BilinearForms},
\nolinkurl{MathGenome.LinearAlgebra.SpectralTheory} and
\nolinkurl{MathGenome.LinearAlgebra.MultilinearAlgebra}.
The accompanying foundation-interface file records inspected declaration names
and types. Geometric volume transformation, differentiability of inverses,
surface area, a moving orthonormal frame, and differential forms with variable
coefficients are not consequences admitted at this interface: they need the
proofs below. In particular an exterior algebra alone does not prove the
naturality or regularity of the exterior derivative.

\section*{F3. One-variable analysis}\refstepcounter{foundation}\label{base:scalar}
The scalar interface includes sequences and completeness in $\R$ and $\C$;
compactness of a closed bounded real interval; scalar continuity, uniform
continuity on such intervals and the intermediate value theorem; ordinary
one-variable differentiation with product and chain rules, Rolle's theorem,
the real mean-value theorem, Taylor's formula with the stated finite
regularity, and the fundamental theorem of calculus. It also includes the
Riemann integral on a compact interval, its Darboux and tagged-sum criteria,
additivity, order and norm bounds, scalar substitution by a $C^1$ monotone
diffeomorphism, and uniform limits and uniform integrable tail arguments.
Exponentials and trigonometric functions and their ordinary derivatives are
available. Finite vector-valued uses are obtained component by component;
they do not admit a common scalar mean-value point for a vector map.

These are the established D03 direction. Relevant public imports include:
\begin{itemize}[leftmargin=*,itemsep=0pt]
\item \nolinkurl{MathGenome.SingleVariableAnalysis.MeanValues}
\item \nolinkurl{MathGenome.SingleVariableAnalysis.Taylor}
\item \nolinkurl{MathGenome.SingleVariableAnalysis.RiemannIntegral}
\item \nolinkurl{MathGenome.SingleVariableAnalysis.IntegrationFormulas}
\end{itemize}
The sequence, series and improper-integral modules are used when indicated. Multivariable
compactness, total differentiation, multivariable Taylor symmetry, multiple
Riemann Fubini, substitution and manifold integration are not admitted here.

\section*{F4. Measure comparison only}\refstepcounter{foundation}\label{base:measure}
For the explicit comparison with Lebesgue integration we use the standard
Lebesgue measure on $\R^n$, with the volume of an axis-parallel box equal to the
product of side lengths. Its elementary interface is countable additivity,
completeness, monotonicity, approximation by countable box covers, measurable
sets and measurable functions, integral order and linearity, and monotone and
dominated convergence. Bounded measurable functions supported on a bounded
box are integrable. These belong to the fixed Mathlib measure-theory base;
their construction is the separate measure-theory direction, not an unstated
proof obligation of classical D04. A null set may be covered by countably many
open boxes of arbitrarily small total volume.

This interface is used to compare the independently constructed Riemann
integral with the measure integral. It does \emph{not} include that comparison
as an assumption. It also does not include a Jacobian substitution theorem,
an area or coarea formula, a measure-theoretic Fubini theorem used in place of
pointwise Riemann Fubini, a multivariable differentiation theorem, partitions
of unity, Poincare's lemma or Stokes' theorem. The classical integration proofs
below do not depend on assuming those conclusions.

\section*{Version and evidence policy}
The inspected dependency base is MathGenome v0.3.0, Lean 4.35.0, and Mathlib
commit \texttt{56a005855ea17e0b0958f73ec801ae140b0a1679}. Existing declarations
were previously accepted in that fixed environment. The present work reads
their source and index; it does not modify that release or claim a new Lean
check for this manuscript. Precise future reuse must retain both the actual
type and the relevant coordinate, integrability and domain bridges.

% END ORIGINAL INPUT: foundations.tex

\mainmatter

% BEGIN ORIGINAL INPUT: chapters/01.tex
\chapter{Euclidean limits, continuity, and local geometry}
\label{ch:d04-01}

\section{Coordinates and neighborhoods}
\family{K04.01.01}\family{K04.01.02}\family{K04.01.03}\family{K04.01.10}

\begin{definition}[The ambient spaces]\label{d04:01:euclidean}
For an integer $n\geq0$, $\R^n$ is the space of functions on
$\{1,\ldots,n\}$, written as columns. Its standard basis is $(e_j)_{j=1}^n$,
its inner product is $\ip{x}{y}=\sum_jx_jy_j$, and
$\norm{x}=\sqrt{\ip{x}{x}}$, $\norm{x}_\infty=\max_j\abs{x_j}$.
When $n=0$, there is one vector, denoted $0$, and both norms are zero;
the maximum in this particular definition is assigned the value zero.
An empty matrix represents the unique linear map of the indicated dimensions.
Complex coordinates, when used, mean pairs of real coordinates unless a
complex-linear assertion is expressly made.
\end{definition}

\begin{proposition}[Coordinate estimates]\label{d04:01:coordinate-estimates}
For $n\geq1$,
\[
 \norm{x}_\infty\leq\norm{x}\leq\sqrt n\norm{x}_\infty.
\]
The Euclidean norm satisfies the triangle inequality, and
$\abs{\norm{x}-\norm{y}}\leq\norm{x-y}$. A sequence in $\R^n$
converges, or is Cauchy, exactly when each coordinate does. In particular
$\R^n$ is complete, and every bounded sequence has a convergent subsequence.
\end{proposition}
\begin{proof}
The displayed inequalities follow by comparing each squared coordinate with
the sum and with the largest squared coordinate. Cauchy--Schwarz from
\cref{base:linear} gives
$\norm{x+y}^2\leq(\norm{x}+\norm{y})^2$; applying the triangle
inequality twice gives the reverse triangle inequality. The coordinate
estimates prove the two sequential equivalences. Coordinatewise scalar
completeness supplies a limit for a Cauchy sequence. Repeatedly take a
convergent subsequence of the first, second, and subsequent coordinates of a
bounded sequence, using scalar Bolzano--Weierstrass from \cref{base:scalar}.
There are only finitely many coordinates, so the final subsequence converges
in norm. The assertions for $n=0$ are immediate.
\end{proof}

\begin{definition}[Ambient and relative topology]\label{d04:01:topology}
Write $B(a,\rho)=\{x:\norm{x-a}<\rho\}$ for $\rho>0$.
A set $U$ is open if each of its points has such a ball in $U$; $F$ is closed
if its complement is open. For $A\subset\R^n$, $\Int A$ consists of the
points with a ball in $A$, $\cl A$ of the points every ball about which meets
$A$, and $\partial A=\cl A\setminus\Int A$. A relative neighborhood of
$a\in A$ contains $A\cap B(a,\rho)$ for some $\rho>0$; a subset of $A$
is relatively open if it is the intersection of $A$ with an ambient open set.
An accumulation point of $A$ is a point whose every punctured ball meets $A$.
This ambient boundary is distinct from the manifold boundary introduced in
\cref{ch:d04-14}.
\end{definition}

\begin{proposition}[Basic topological tests]\label{d04:01:closure-tests}
Arbitrary unions and finite intersections of open sets are open. A set is
closed exactly when it contains the limits of its convergent sequences.
Moreover, $a\in\cl A$ exactly when some sequence in $A$ converges to $a$;
$a$ is an accumulation point exactly when that sequence can be chosen in
$A\setminus\{a\}$. Thus $\cl A$ is closed, $\Int A$ is open, and
$\partial A=\cl A\cap\cl(\R^n\setminus A)$ is closed.
\end{proposition}
\begin{proof}
At a point in a union choose one of its constituent sets. For a finite
intersection take the least of the finitely many ball radii. If every ball
about $a$ meets $A$, choose $a_k\in A\cap B(a,1/k)$; the converse follows
from convergence. The same construction with punctured balls proves the
accumulation statement. If $F$ is closed and a sequence in $F$ converges to
$a\notin F$, a ball about $a$ in its open complement gives a contradiction.
Conversely, if a point outside $F$ has no ball disjoint from $F$, the sequence
construction contradicts sequential closedness. If $a\notin\cl A$, some ball
about $a$ misses $A$, and a smaller ball about every sufficiently close point
also misses $A$; hence $\cl A$ is closed. The union of all balls contained
in $A$ is $\Int A$. Finally a point fails to be interior to $A$ precisely
when every ball meets its complement, proving the boundary formula.
\end{proof}

\section{Limits and continuity on specified domains}
\family{K04.01.04}\family{K04.01.05}

\begin{definition}[Relative limits and continuity]\label{d04:01:limits}
For $f:A\to\R^m$ and an accumulation point $a$ of $A$, the assertion
$\lim_{x\to a,\ x\in A}f(x)=b$ means that for every $\eps>0$ there is
$\delta>0$ such that $x\in A$ and $0<\norm{x-a}<\delta$ imply
$\norm{f(x)-b}<\eps$. The map is continuous at $a\in A$ when the same
condition, with $\norm{x-a}<\delta$ and $b=f(a)$, holds. It is uniformly
continuous on $A$ when $\delta$ can be chosen independently of $a\in A$.
At an isolated domain point continuity is automatic. A punctured limit there
would be vacuous, which is why accumulation was required in its definition.
\end{definition}

\begin{theorem}[Sequential and coordinate criteria]\label{d04:01:sequential-continuity}
A relative limit holds exactly when $f(x_k)\to b$ for every sequence
$x_k\in A\setminus\{a\}$ with $x_k\to a$. A map is continuous at $a$
exactly when it preserves every sequence in $A$ converging to $a$. Limits are
unique. A limit or continuity assertion for a finite-dimensional vector target
is equivalent to the corresponding assertion for every coordinate.
\end{theorem}
\begin{proof}
The epsilon--delta condition implies the sequential condition by applying its
$\delta$ to a tail of the sequence. If it fails, choose a fixed
$\eps_0>0$ and, for each $k$, a violating point at distance less than $1/k$.
These points give a sequence violating the sequential condition. This also
proves the continuity version, without puncturing. If $b\ne c$ were two
limits, choose $\eps=\norm{b-c}/3$ and a domain point in the smaller of
the two punctured balls; the triangle inequality is impossible. Existence of
that point uses accumulation. The coordinate assertions follow from
\cref{d04:01:coordinate-estimates}, choosing the minimum of finitely many
deltas in the reverse implication.
\end{proof}

\begin{proposition}[Continuous operations]\label{d04:01:continuous-operations}
Continuous maps compose, and a map is continuous exactly when preimages of
open sets are relatively open. Sums, scalar multiples, inner products and
finite matrix products of continuous maps are continuous. A scalar quotient
is continuous wherever its denominator is nonzero.
\end{proposition}
\begin{proof}
For composition first choose a target tolerance for the outer map, then a
source tolerance for the inner map. Equivalently use the sequential criterion.
If $f$ is continuous and $f(a)\in O$ with $O$ open, choose a ball about
$f(a)$ in $O$ and pull it back using continuity at $a$. Conversely apply
the preimage condition to each ball about $f(a)$. Coordinate sums and products
are continuous by scalar limit algebra; inner and matrix products are finite
sums of these. If $q(a)\ne0$, continuity gives
$\abs{q(x)}\geq\abs{q(a)}/2$ nearby, and
$\abs{q(x)^{-1}-q(a)^{-1}}\leq2\abs{q(x)-q(a)}/\abs{q(a)}^2$.
\end{proof}

\begin{example}[Approach sets matter]\label{d04:01:path-tests}
On $\R^2\setminus\{0\}$ set
$f(x,y)=x^2y/(x^4+y^2)$. Along a line $y=mx$, its value tends to zero
as $(x,y)\to0$ (and on the vertical axis it is zero). Along $y=x^2$
it equals $1/2$. Consequently all straight-line tests do not establish a
two-dimensional limit. On the relative domain $\{(x,x^2):x>0\}$ the limit
is $1/2$; on the horizontal axis it is zero. The same formula thus has
different limits on different approach domains, without any contradiction.
\end{example}

\section{Compactness and its quantitative consequences}
\family{K04.01.06}\family{K04.01.07}

\begin{definition}[Compactness]\label{d04:01:compact-definition}
A subset $K$ of a Euclidean space is compact if every cover of it by ambient
open sets has a finite subcover. It is sequentially compact if every sequence
in $K$ has a subsequence converging to a point of $K$.
\end{definition}

\begin{theorem}[Finite-dimensional compactness]\label{d04:01:compactness}
For a Euclidean subset $K$, compactness, sequential compactness, and being
closed and bounded are equivalent. The empty set and the one-point space
$\R^0$ obey these assertions.
\end{theorem}
\begin{proof}
A compact set is bounded because the balls $B(0,j)$, $j\geq1$, cover it.
It is closed: for $a\notin K$, the balls
$B(x,\norm{x-a}/3)$, $x\in K$, cover $K$; finitely many suffice, and
the minimum of their positive radii is a radius of a ball about $a$ missing
all those balls. A closed bounded set is sequentially compact by
\cref{d04:01:coordinate-estimates,d04:01:closure-tests}.

Suppose now $K$ is sequentially compact. For each $\eta>0$ it has a finite
$\eta$-net: otherwise recursively choose points whose distances from all
previously chosen points exceed or equal $\eta$, contradicting the existence
of a convergent subsequence. Let $\mathcal U$ be an open cover. There is a
$\delta>0$ such that each $K\cap B(x,\delta)$, $x\in K$, is contained
in one member of $\mathcal U$. If no such $\delta$ existed, choose $x_j\in K$
for which the assertion fails at $1/j$. A convergent subsequence tends to
$x\in K$; some cover member contains $B(x,2\rho)$. For large indices
$\norm{x_j-x}<\rho$ and $1/j<\rho$, making the purported failure
impossible. Choose a finite $\delta/2$-net with centers in $K$, and for each
center choose the cover member containing its $\delta$-ball in $K$. These
finitely many members cover $K$. The empty cases require no choices.
\end{proof}

\begin{lemma}[Lebesgue number]\label{d04:01:lebesgue-number}
If a compact set $K$ is covered by relatively open sets $U_\alpha$, there
is a $\delta>0$ such that each nonempty subset $S\subset K$ with
$\diam S<\delta$ lies in one $U_\alpha$. In particular each sufficiently
small ball in $K$ lies in one cover member.
\end{lemma}
\begin{proof}
The sequential argument in \cref{d04:01:compactness}, applied to ambient open
sets representing the relative opens, supplies $\eta>0$ for which
$K\cap B(x,\eta)$ lies in a cover member for every $x\in K$. Take
$\delta=\eta$. Choosing $x\in S$ gives $S\subset K\cap B(x,\eta)$.
For the final assertion take balls of radius less than $\eta/2$.
\end{proof}

\begin{theorem}[Continuous maps on compact sets]\label{d04:01:compact-continuity}
If $f:K\to\R^m$ is continuous and $K$ is compact, then $f(K)$ is compact,
$f$ is bounded, and $f$ is uniformly continuous. If $m=1$ and $K\ne\varnothing$,
$f$ attains its minimum and maximum.
\end{theorem}
\begin{proof}
Pull back any open cover of $f(K)$ and take a finite subcover of $K$.
Boundedness follows from \cref{d04:01:compactness}. If uniform continuity
failed, there would be $x_j,y_j\in K$ with
$\norm{x_j-y_j}<1/j$ and $\norm{f(x_j)-f(y_j)}\geq\eps_0>0$.
Take a subsequence $x_j\to x\in K$. Then $y_j\to x$, contradicting
continuity. For a nonempty scalar image, boundedness gives its supremum and
infimum by \cref{base:real}; sequences approaching them have their limits
in the closed image. Hence both are image values.
\end{proof}

\begin{proposition}[Distance and separation]\label{d04:01:distance}
For nonempty $A\subset\R^n$ define
$\dist(x,A)=\inf_{a\in A}\norm{x-a}$. This function is $1$-Lipschitz,
and it vanishes exactly on $\cl A$. If $K$ is nonempty compact and $F$ is
nonempty closed with $K\cap F=\varnothing$, then $\dist(K,F)>0$.
For every nonempty compact $K\subset U$ with $U$ open there is a $\rho>0$ such that
$\{x:\dist(x,K)\leq\rho\}\subset U$; this neighborhood is compact.
\end{proposition}
\begin{proof}
The triangle inequality gives
$\dist(x,A)\leq\norm{x-y}+\dist(y,A)$ by taking the infimum over
$A$; interchange $x,y$. The zero-set assertion follows directly from the
ball definition of closure. On $K$ the positive continuous function
$\dist(\cdot,F)$ attains a positive minimum by
\cref{d04:01:compact-continuity}. Apply this to $F=\R^n\setminus U$
and choose $\rho$ less than that minimum. The indicated neighborhood is
closed and bounded and therefore compact. If $U=\R^n$, take any $\rho>0$;
no distance to an empty complement is needed. Empty compact sets can be
handled separately by taking their thickening to be empty; the notation
$\dist(x,K)$ above is used only for nonempty $K$.
\end{proof}

\begin{theorem}[Equivalence of norms]\label{d04:01:norm-equivalence}
If $N$ is any norm on $\R^n$, there are constants $c,C>0$ with
$c\norm{x}\leq N(x)\leq C\norm{x}$ for all $x$. Thus all norms on
a fixed finite-dimensional real space give the same limits, topology,
compact sets, and differentiability remainders.
\end{theorem}
\begin{proof}
For $n>0$, the triangle inequality gives
$N(x)\leq\sum_j\abs{x_j}N(e_j)\leq C\norm{x}$ with
$C=\sum_jN(e_j)>0$. Therefore
$\abs{N(x)-N(y)}\leq C\norm{x-y}$, so $N$ is Euclidean-continuous.
The Euclidean unit sphere is compact. Its positive $N$-values have a positive
minimum $c$. Homogeneity proves the lower bound, and the upper bound was
already obtained. Each ball for either norm contains a suitable ball for the
other; the stated consequences follow. In dimension zero take $c=C=1$.
No infinite-dimensional norm-equivalence assertion is being made.
\end{proof}

\section{Connected domains and linear estimates}
\family{K04.01.08}\family{K04.01.09}

\begin{definition}[Connectedness and paths]\label{d04:01:connected-definition}
A space is connected if it is not the union of two disjoint, nonempty
relatively open sets. A path in $A$ is a continuous map $[0,1]\to A$;
the space is path connected if every pair of its points can be joined by a
path. A polygonal path is a finite concatenation of line segments. A set is
convex if it contains every segment with endpoints in it. A path connected
space is simply connected if each loop can be continuously contracted to a
point within the space while keeping its basepoint fixed. Only the first
two notions are needed for the following theorem.
\end{definition}

\begin{theorem}[Connectedness in Euclidean domains]\label{d04:01:connectedness}
Continuous images of connected spaces are connected. Intervals in $\R$ are
connected. A path connected space is connected. An open connected subset of
$\R^n$ is polygonally connected. In particular convex sets are path connected.
A continuous real-valued map on a connected space has the intermediate value
property.
\end{theorem}
\begin{proof}
A separation of an image pulls back to a separation. If an interval had a
separation and $a<b$ lay in the two pieces, take the supremum of the points
of the first piece in $[a,b]$; relative openness of whichever piece contains
that supremum contradicts its defining property. A path joining points in
different pieces of a separation would separate $[0,1]$ by preimages, so
path connected implies connected. For open $U$, fix $a\in U$ and let $P$
be the points polygonally reachable from $a$. A small ball about a point of
$P$ lies in $P$ because balls are convex. A small ball about a point outside
$P$ also stays outside $P$, since a segment from a reachable point of that
ball would make its center reachable. Thus $P$ and its complement are
relatively open; connectedness gives $P=U$. Convexity supplies a single
segment. Finally a subset of $\R$ omitting a value between two of its points
is separated by the two open half-lines at that value.
\end{proof}

\begin{corollary}[Connected components of an open set]\label{d04:01:components}
For a point $a$ in a space $U$, let $C_a$ be the union of all connected
subsets of $U$ containing $a$. This is the unique maximal connected subset
containing $a$, called its connected component. The components partition
$U$. If $U\subset\R^n$ is open, every component is open and polygonally
connected.
\end{corollary}
\begin{proof}
The union of connected sets having a common point is connected: in a
purported separation of the union, each constituent lies wholly in the
piece containing that common point, leaving the other piece empty. Thus
$C_a$ is connected and is maximal by its definition. If two components
meet, their union is connected by the same argument with an intersection
point; maximality makes them equal. Since every singleton is connected,
the components cover $U$ and form a partition. Now suppose $U$ is open
and $x\in C_a$. A sufficiently small ball about $x$ lies in $U$ and is
connected by convexity. Its union with $C_a$ is connected, so maximality
puts the whole ball in $C_a$. Thus $C_a$ is open, and
\cref{d04:01:connectedness} makes it polygonally connected. Empty $U$
has no components; the proof also covers the one-point space $\R^0$.
\end{proof}

\begin{definition}[Operator norm]\label{d04:01:operator-norm}
For a linear map $A:\R^n\to\R^m$, set
$\norm{A}=\sup_{\norm{v}\leq1}\norm{Av}$. If $n=0$ this is zero.
For a $k$-linear map $T$ use the supremum of $\norm{T(v_1,\ldots,v_k)}$
over $\norm{v_i}\leq1$; for $k=0$ this means the norm of its value.
\end{definition}

\begin{proposition}[Linear and multilinear coordinate bridges]\label{d04:01:operator-estimates}
All finite-dimensional linear and multilinear maps are continuous and have
finite operator norm. For linear maps
\[
 \norm{Av}\leq\norm{A}\norm{v},\qquad
 \norm{BA}\leq\norm{B}\norm{A}.
\]
If $A=(a_{ij})$ with $m,n>0$, then
\[
 \max_{i,j}\abs{a_{ij}}\leq\norm{A}
 \leq\Big(\sum_{i,j}a_{ij}^2\Big)^{1/2}
 \leq\sqrt{mn}\max_{i,j}\abs{a_{ij}}.
\]
Consequently convergence and continuity in operator norm are equivalent to
coordinate convergence and continuity. Analogous coordinate statements hold
for each fixed multilinear order.
\end{proposition}
\begin{proof}
Expanding in basis vectors expresses each output coordinate of a multilinear
map as a finite sum of products of input coordinates; it is continuous by
\cref{d04:01:continuous-operations}. Compactness of a product of finitely many
unit balls (it is closed and bounded in its coordinate space) makes its norm
bounded and its supremum attained. Scaling each nonzero input proves the
multilinear estimate; a zero input gives zero if the order is positive.
The composition estimate follows by applying the linear estimate twice.
For matrices, $\abs{a_{ij}}=\abs{(Ae_j)_i}\leq\norm{A}$;
Cauchy--Schwarz applied separately to every row gives
$\norm{Av}^2\leq\sum_{i,j}a_{ij}^2\norm{v}^2$. The final bound is a
finite-sum comparison. For a multilinear map the same expansion bounds its
norm by the sum of absolute coordinate coefficients, while each coefficient
is bounded by its operator norm. These two inequalities prove the last
assertions, including the zero-dimensional cases with empty coefficients.
\end{proof}

\source{Original finite-dimensional proofs; comparison passages checked:
HUB 1.6, printed pp. 104--106 (PDF pp. 121--123), and Appendix A.3,
printed pp. 713--714 (PDF pp. 730--731), including its finite-cover proof.
HUB initially defines nonempty compact sets as closed and bounded; here
the open-cover definition is used and the equivalence, including the empty
case, is proved. FLE Chapters 1--2 were checked at contents level only.
No general metric-space compactness theorem is admitted here.}

% END ORIGINAL INPUT: chapters/01.tex


% BEGIN ORIGINAL INPUT: chapters/02.tex
\chapter{Total derivatives and the chain rule}
\label{ch:d04-02}

\section{One derivative object}
\family{K04.02.01}\family{K04.02.02}\family{K04.02.03}\family{K04.02.04}

Throughout this chapter $U\subset\R^n$ is open, $n,m\geq0$, and
$f:U\to\R^m$. The openness requirement ensures that all sufficiently small
increments, in every direction, are available.

\begin{definition}[Directional, partial, and total derivatives]\label{d04:02:derivative}
The directional derivative in the vector $v$ at $a\in U$ is
\[
 \partial_vf(a)=\lim_{t\to0}\frac{f(a+tv)-f(a)}t
\]
when it exists. Write $\partial_jf=\partial_{e_j}f$. The map is totally
differentiable at $a$ if there is a linear map $A:\R^n\to\R^m$ such that
\[
 f(a+h)=f(a)+Ah+r(h),\qquad
 \lim_{h\to0,\ h\ne0}\frac{\norm{r(h)}}{\norm{h}}=0.
\]
We write $Df(a)=A$ and $r(h)=o(\norm{h})$, setting $r(0)=0$.
For a general subset the same remainder can be required only for available
increments, but it need not determine $A$ uniquely. We will use that relative
condition only when explicitly stated, and will not identify it with this
open-domain derivative.
\end{definition}

\begin{proposition}[Uniqueness, continuity, and coordinates]\label{d04:02:derivative-bridge}
The total derivative is unique. Differentiability implies continuity and
$\partial_vf(a)=Df(a)v$ for every $v$. If $f=(f_1,\ldots,f_m)$,
then $f$ is differentiable exactly when all $f_i$ are, and the unique matrix
of $Df(a)$ in the standard column bases is
\[
 Jf(a)_{ij}=\partial_jf_i(a).
\]
Conversely a proposed matrix $J$ is the Jacobian of the total derivative
exactly when the map $h\mapsto Jh$ has the stated small remainder.
For scalar $f$, its gradient is the column $\nabla f(a)$ characterized by
$Df(a)h=\ip{\nabla f(a)}{h}$.
\end{proposition}
\begin{proof}
If $A,B$ both work, putting $h=tv$ gives
$\norm{(A-B)v}=o(1)$ for each fixed nonzero $v$ as $t\to0$.
Thus $A=B$. Also
$\norm{f(a+h)-f(a)}\leq(\norm{A}+1)\norm{h}$ for sufficiently
small $h$, proving continuity. Dividing the remainder at $tv$ by $t$ proves
the directional formula; for $v=0$ it is immediate. Coordinate projection
of a small vector remainder is small, and conversely finitely many small
coordinate remainders give a small vector remainder by
\cref{d04:01:coordinate-estimates}. The columns of a linear map are its
values on $e_j$, proving the matrix statement and its converse. The gradient
identity follows by expanding the scalar linear map in its basis coefficients.
For $n=0$ the only linear map has zero domain; for $m=0$ every assertion is
the unique zero assertion. There is no uniqueness problem in these cases.
\end{proof}

\begin{example}[Why the domain is part of the definition]\label{d04:02:relative-derivative}
On $A=\{(x,0):x\in\R\}$ the function $f(x,0)=x$ satisfies the relative
linear-approximation condition at zero both with $A_1(h)=h_1$ and with
$A_2(h)=h_1+h_2$. On an isolated domain point every linear map satisfies
the relative condition. Open-domain uniqueness cannot be exported to these
situations. Derivatives along regular submanifolds will instead be defined
on their tangent spaces through charts.
\end{example}

\section{Algebra and composition}
\family{K04.02.05}\family{K04.02.06}\family{K04.02.07}

\begin{theorem}[Chain rule]\label{d04:02:chain-rule}
Let $V\subset\R^m$ be open, $f:U\to V$ differentiable at $a$, and
$g:V\to\R^p$ differentiable at $b=f(a)$. Then
\[
 D(g\circ f)(a)=Dg(f(a))\circ Df(a),\qquad
 J(g\circ f)(a)=Jg(f(a))Jf(a).
\]
The assertion requires no continuity of either derivative near its basepoint.
\end{theorem}
\begin{proof}
Write $f(a+h)=b+Ah+r(h)$ and $g(b+k)=g(b)+Bk+s(k)$ with small
remainders, both set to zero at zero. Put $k=Ah+r(h)$. For small $h$,
$\norm{k}\leq(\norm{A}+1)\norm{h}$, so $k\to0$. Substitution gives
\[
 g(f(a+h))=g(b)+BAh+Br(h)+s(Ah+r(h)).
\]
The first error is $o(\norm{h})$. Given $\eps>0$, choose a neighborhood
where $\norm{s(k)}\leq\eps\norm{k}$, including $k=0$. Its contribution
is at most $\eps(\norm{A}+1)\norm{h}$ for sufficiently small $h$.
Thus the second error is small also, including increments for which the
inner increment vanishes. Openness of $V$ and continuity of $f$ ensure that
all expressions are defined nearby. The matrix identity follows from
\cref{d04:02:derivative-bridge} and ordinary composition of column matrices.
\end{proof}
\source{HUB Appendix A.4, printed pp. 715--716, physical PDF pp. 732--733:
the two-remainder proof and the zero-inner-increment issue were checked.
The proof above is an original formulation of that argument.}

\begin{proposition}[Linear and bilinear rules]\label{d04:02:bilinear-rule}
Linear maps have themselves as derivatives. Derivatives commute with finite
sums and constant scalar multiplication. If $B:E\times F\to G$ is a fixed
bilinear map between finite-dimensional real spaces and $u,v$ are
differentiable, then
\[
 D(B(u,v))(a)h=B(Du(a)h,v(a))+B(u(a),Dv(a)h).
\]
This includes multiplication of scalars or matrices, the inner product,
and the cross product in $\R^3$ with their declared orders.
\end{proposition}
\begin{proof}
For a linear map the remainder is identically zero; finite sums of small
remainders are small. Set $\Delta u=u(a+h)-u(a)$ and similarly for $v$.
Bilinearity gives
\[
 B(u+\Delta u,v+\Delta v)-B(u,v)
 =B(\Delta u,v)+B(u,\Delta v)+B(\Delta u,\Delta v).
\]
The first two terms have the stated linear parts and small remainders by
\cref{d04:01:operator-estimates}. Differentiability bounds both increments
by a constant times $\norm{h}$, so the final term is $O(\norm{h}^2)$,
hence $o(\norm{h})$. Matrix order is preserved because $B$ need not be
commutative.
\end{proof}

\begin{proposition}[Scalar reciprocal and quotient]\label{d04:02:quotient-rule}
If scalar $q$ is differentiable at $a$ and $q(a)\ne0$, then $1/q$ is defined
on a neighborhood and
\[
 D(q^{-1})(a)h=-\frac{Dq(a)h}{q(a)^2}.
\]
If $u$ is vector-valued and differentiable, then
$D(u/q)(a)h=Du(a)h/q(a)-u(a)Dq(a)h/q(a)^2$.
There is no entrywise-matrix reciprocal in this assertion.
\end{proposition}
\begin{proof}
Continuity of $q$ keeps $\abs{q}$ bounded below nearby. For $c\ne0$,
\[
 \frac1{c+t}=\frac1c-\frac t{c^2}+
 \frac{t^2}{c^2(c+t)}.
\]
The last term is $O(t^2)$. Substitute $t=q(a+h)-q(a)=O(\norm{h})$
and use the chain rule, then the scalar-vector bilinear rule for the quotient.
Matrix inversion and its derivative will be proved in
\cref{d04:04:inverse-stability,d04:04:matrix-inversion}.
\end{proof}

\section{Mean-value bounds and continuous derivatives}
\family{K04.02.08}\family{K04.02.09}\family{K04.02.10}\family{K04.02.11}

\begin{theorem}[Segment mean-value estimate]\label{d04:02:mean-value}
Suppose $[a,b]\subset U$, $f$ is continuous on this segment and
differentiable at all its interior points (in an ambient open neighborhood).
If $\norm{Df(x)}\leq M$ there, then
\[
 \norm{f(b)-f(a)}\leq M\norm{b-a}.
\]
For scalar $f$ there is, when $a\ne b$, a $t\in(0,1)$ such that
$f(b)-f(a)=Df(a+t(b-a))(b-a)$. For vector $f$ the inequality holds,
but such a common exact intermediate point generally does not exist.
\end{theorem}
\begin{proof}
The scalar assertion is the one-variable mean-value theorem applied to
$t\mapsto f(a+t(b-a))$, whose derivative is supplied by
\cref{d04:02:chain-rule}. If the vector difference is nonzero, put
$w=(f(b)-f(a))/\norm{f(b)-f(a)}$ and apply that scalar assertion to
$\ip{w}{f(a+t(b-a))}$. Its endpoint difference equals the norm of the
vector difference, and its derivative has absolute value at most
$M\norm{b-a}$. A zero difference and a zero-length segment are immediate.
The curve $f(t)=(t^2,t^3)$ on $[0,1]$ disproves the vector equality:
$(1,1)=(2c,3c^2)$ would require simultaneously $c=1/2$ and $c^2=1/3$.
\end{proof}

\begin{corollary}[Constancy and Lipschitz bounds]\label{d04:02:zero-derivative}
On a convex open set a bound $\norm{Df}\leq M$ makes $f$ $M$-Lipschitz.
On a connected open set, $Df=0$ implies that $f$ is constant. A $C^1$ map
is locally Lipschitz.
\end{corollary}
\begin{proof}
Apply \cref{d04:02:mean-value} on each segment. When $Df=0$, concatenate
the finitely many segments supplied by \cref{d04:01:connectedness}.
For a $C^1$ map, continuity of $Df$ bounds it on a sufficiently small closed
ball about any given point; the open ball is convex.
\end{proof}

\begin{theorem}[Continuous partial derivatives]\label{d04:02:continuous-partials}
Suppose every $\partial_jf$ exists throughout a neighborhood of $a$ and is
continuous at $a$. Then $f$ is differentiable at $a$, with columns
$\partial_jf(a)$. In particular, if all first partials are continuous on
$U$, then $f$ has a continuous total derivative on $U$.
\end{theorem}
\begin{proof}
It suffices to prove the scalar assertion and then use
\cref{d04:02:derivative-bridge}. Choose a coordinate box about $a$ inside
the neighborhood. For $h$ small set
$a^{(j)}=a+\sum_{i\leq j}h_ie_i$. Telescoping and the scalar mean-value
theorem on each coordinate segment give
\[
 f(a+h)-f(a)=\sum_{j=1}^nh_j\partial_jf(c_j),
\]
where $c_j$ lies on $[a^{(j-1)},a^{(j)}]$; if $h_j=0$ its summand is
zero and choose either endpoint. Each $\norm{c_j-a}\leq\norm{h}$.
After subtracting $\sum_jh_j\partial_jf(a)$, the absolute error is at most
\[
 \sum_j\abs{h_j}\abs{\partial_jf(c_j)-\partial_jf(a)}
 \leq\sqrt n\norm{h}\max_j\abs{\partial_jf(c_j)-\partial_jf(a)}.
\]
The maximum tends to zero by finitely many continuity assertions. This proves
the remainder condition. Continuity of the resulting derivative follows from
the coordinate/operator-norm bridge in \cref{d04:01:operator-estimates}.
For $n=0$ the assertion needs no summation.
\end{proof}

\begin{definition}[Finite differentiability classes]\label{d04:02:cr}
On an open set, $C^0$ means continuous. Recursively $f$ is $C^{r+1}$ if
it is differentiable everywhere and $x\mapsto Df(x)$, as a map to the
finite-dimensional space $\Hom(\R^n,\R^m)$, is $C^r$.
Here $r$ is a nonnegative integer. The class $C^\infty$ means $C^r$ for
every finite $r$. Existing derivatives without continuity are not the same
condition. A map is real analytic if it has a convergent power series locally;
no analyticity conclusion follows merely from membership in every $C^r$.
\end{definition}

\begin{proposition}[Coordinate criterion and closure of $C^r$]\label{d04:02:cr-calculus}
For finite $r\geq1$, $f$ is $C^r$ exactly when every iterated coordinate
partial of order at most $r$ exists and is continuous. The order of iteration
is retained in this criterion until symmetry is proved in
\cref{d04:03:symmetry}. Sums, fixed multilinear operations, composition,
and scalar quotients with nonvanishing denominator preserve $C^r$,
including $r=0$ and $r=\infty$.
\end{proposition}
\begin{proof}
For $r=1$ combine \cref{d04:02:continuous-partials} with the coordinate
representation of $Df$. Inductively, the entries of $Df$ are the first
partials, and taking their iterated derivatives is exactly taking iterated
partials of $f$ of one additional order. Coordinate identification is a
linear isomorphism, so \cref{d04:02:bilinear-rule,d04:02:chain-rule} show
that this identification preserves differentiability; continuity was proved
in \cref{d04:01:operator-estimates}. This proves the criterion by induction.
For closure, the case $r=0$ is
\cref{d04:01:continuous-operations}. For the step from $r$ to $r+1$, the
first derivative formulas already proved express the derivative of each
operation as sums or products of the original maps and their first
derivatives. Apply the order-$r$ closure statement. In the composition case
the formula is $(Dg\circ f)Df$, with both factors of class $C^r$.
The quotient uses $-q^{-2}Dq$ and a neighborhood on which $q\ne0$.
A fixed higher multilinear operation can be treated one slot at a time by
its finite expansion, or by the same increment expansion as for a bilinear
map. Applying every finite-order conclusion proves the smooth case.
\end{proof}

\begin{proposition}[Integration along a segment]\label{d04:02:segment-ftc}
For $f\in C^1(U,\R^m)$ and $[a,a+h]\subset U$,
\[
 f(a+h)-f(a)=\int_0^1Df(a+th)h\,\dd t.
\]
The integral of a continuous vector function means its finitely many scalar
Riemann integrals. It satisfies
$\norm{\int v}\leq\int\norm{v}$.
\end{proposition}
\begin{proof}
The chain rule and the scalar fundamental theorem of calculus apply to each
coordinate of $t\mapsto f(a+th)$. For the norm inequality let $z=\int v$.
If $z\ne0$, pair the integral with $z/\norm{z}$; scalar linearity and
Cauchy--Schwarz give $\norm{z}\leq\int\norm{v}$. If $z=0$ it is immediate.
Thus only one-variable integration from \cref{base:scalar} has been used.
\end{proof}

\section{Separating the hypotheses}
\family{K04.02.12}

\begin{example}[Partials, directions, and differentiability]\label{d04:02:counterexamples}
Define all the following functions to be zero at the origin.
First, $u(x,y)=xy/(x^2+y^2)$ away from the origin has both coordinate
partials equal to zero there but is discontinuous: on $y=x$ it equals $1/2$.
Second, $v(x,y)=x^3/(x^2+y^2)$ is continuous because
$\abs{v(x,y)}\leq\norm{(x,y)}$. Every directional derivative exists,
but in direction $(a,b)\ne0$ its value is $a^3/(a^2+b^2)$, which is
not linear in $(a,b)$ (use $(1,0),(0,1),(1,1)$). Thus $v$ is not totally
differentiable. Third,
\[
 w(x,y)=\frac{x^3y}{x^4+y^2}
\]
is continuous since $2\abs{x^2y}\leq x^4+y^2$ gives
$\abs{w}\leq\abs{x}/2$. In a fixed direction $(a,b)$ with $b\ne0$,
\[
 \frac{w(ta,tb)}t=\frac{t a^3b}{t^2a^4+b^2}\longrightarrow0;
\]
if $b=0$ it is zero. Thus even all directional derivatives assembling to
the zero linear map do not imply total differentiability. Indeed on
$(x,x^2)$ its value is $x/2$, whose ratio to $\norm{(x,x^2)}$ does not
tend to zero. Finally the scalar function
$z(x)=x^2\sin(1/x)$ for $x\ne0$, $z(0)=0$, is differentiable at zero
with derivative zero, but $z'(x)=2x\sin(1/x)-\cos(1/x)$ does not tend
to zero. Total differentiability does not imply continuous partials.
All formulas follow by direct substitution and one-variable differentiation.
\end{example}

\source{Original exposition and proofs. Comparison passages actually checked:
HUB Appendix A.4; SPI Theorem 2-10 and Theorem 2-11, printed p. 35
(PDF p. 48), for coordinate mean-value estimates and the total-derivative
interface. Stronger local solvability conclusions are proved next rather
than assumed from those references.}

% END ORIGINAL INPUT: chapters/02.tex


% BEGIN ORIGINAL INPUT: chapters/03.tex
\chapter{Higher derivatives and Taylor formulas}
\label{ch:d04-03}

\section{Iterated linear maps and symmetry}
\family{K04.03.01}\family{K04.03.02}\family{K04.03.03}\family{K04.03.04}

\begin{definition}[Higher differentials]\label{d04:03:higher-derivative}
The space of $k$-linear maps $(\R^n)^k\to\R^m$ has the operator norm of
\cref{d04:01:operator-norm}. Put $D^0f=f$. Recursively the derivative of
the map $x\mapsto D^kf(x)$ is identified with a $(k+1)$-linear map by
\[
 D^{k+1}f(a)[v_0,v_1,\ldots,v_k]
   =\big(D(D^kf)(a)v_0\big)[v_1,\ldots,v_k].
\]
This is an actual linear identification: expand on the finitely many standard
basis tuples, as in \cref{d04:01:operator-estimates}. Say $f$ is $k$ times
differentiable at $a$ if the successive derivative maps through order $k-1$
are defined on a neighborhood of $a$ and the last of them is differentiable
at $a$ (with the recursive interpretation at the lower orders).
This condition imposes no continuity of the top derivative on a neighborhood.
Write $T[h^k]=T[h,\ldots,h]$ and $T[h^0]=T$.
\end{definition}

\begin{theorem}[Pointwise crossed partials]\label{d04:03:crossed-partials}
Suppose all first partials of $f$ exist in a neighborhood of $a$ and the
maps $\partial_if$ are totally differentiable at $a$. Then
$\partial_j\partial_if(a)=\partial_i\partial_jf(a)$ for every $i,j$.
In particular the conclusion holds for a twice totally differentiable map
at $a$; it does not require continuous second derivatives on a neighborhood.
\end{theorem}
\begin{proof}
Work in one scalar output coordinate, since finitely many coordinates suffice.
For small nonzero $t$ consider the square increment
\[
 Q_t=f(a+te_i+te_j)-f(a+te_i)-f(a+te_j)+f(a).
\]
Apply the one-variable mean-value theorem in the $i$ direction to
$s\mapsto f(a+se_i+te_j)-f(a+se_i)$ between $0$ and $t$.
It gives a $c_t$ between these endpoints with
\[
 Q_t=t\big(\partial_if(a+c_te_i+te_j)-\partial_if(a+c_te_i)\big).
\]
Differentiability of $\partial_if$ at $a$ expands the bracket as
$t\partial_j\partial_if(a)+o(\abs{t})$: its two remainders are small
at vectors of norm at most $\sqrt2\abs{t}$, and the common $c_te_i$
linear term cancels. Consequently $Q_t/t^2\to\partial_j\partial_if(a)$.
The expression $Q_t$ is unchanged on interchanging $i,j$, proving equality.
If $f$ is twice totally differentiable, its partials are coordinate
evaluations of $Df$, so they are differentiable at $a$ by the chain rule.
\end{proof}
\source{HUB Theorem 3.3.8 and Appendix A.9, printed pp. 732--733
(PDF pp. 749--750): the pointwise differentiability hypothesis and square
increment proof were checked.}

\begin{theorem}[Symmetry and coordinate representation]\label{d04:03:symmetry}
If $f$ is $k$ times totally differentiable at $a$, $D^kf(a)$ is symmetric.
Its values on basis tuples are the corresponding ordered $k$th coordinate
partials. Thus for $f\in C^k$ these partials are independent of their order.
For scalar $f$ the Hessian matrix $H_f(a)$ has entries
$\partial_i\partial_jf(a)$ and
\[
 D^2f(a)[u,v]=u^TH_f(a)v,\qquad
 D^2f(a)[h,h]=h^TH_f(a)h.
\]
The quadratic Taylor term is one half of this last expression.
\end{theorem}
\begin{proof}
Coordinate evaluation is linear, so differentiating it commutes with taking
a total derivative. Induction on $k$ proves the basis-tuple assertion.
The cases $k=0,1$ have no nontrivial permutation. For $k=2$ use
\cref{d04:03:crossed-partials} on basis vectors and extend by bilinearity.
Inductively, $D^{k-1}f(x)$ is symmetric for $x$ near $a$ because those
derivatives exist there. Differentiating the equality for each permutation
shows that $D^kf(a)$ is symmetric in its last $k-1$ slots. On the other
hand, apply the $k=2$ result to the finite-dimensional valued map
$D^{k-2}f$ at $a$. It interchanges the first two slots of $D^kf(a)$.
These interchanges generate all adjacent transpositions, hence every
permutation. The Hessian identities are the basis expansions of the bilinear
map, and the factor $1/2$ will also follow from Taylor's formula.
\end{proof}

\begin{lemma}[Derivatives of line restrictions and homogeneous terms]\label{d04:03:line-restriction}
When the indicated derivatives exist on a segment,
$F(t)=f(a+th)$ satisfies $F^{(j)}(t)=D^jf(a+th)[h^j]$.
If $T$ is a symmetric $j$-linear map, the polynomial $p(x)=T[x^j]/j!$
has derivative $Dp(x)v=T[v,x^{j-1}]/(j-1)!$ for $j\geq1$.
\end{lemma}
\begin{proof}
Apply the chain rule repeatedly to the affine line and to the linear
evaluation map on multilinear tensors. For the polynomial, expand
$T[(x+v)^j]$ into the $2^j$ choices of $x$ or $v$ in each slot.
The $j$ terms containing exactly one $v$ are equal by symmetry; the terms
containing two or more have norm $O(\norm{v}^2)$ for $x$ fixed. Dividing
by $j!$ proves the derivative formula.
\end{proof}

\section{Pointwise and segment Taylor formulas}
\family{K04.03.05}\family{K04.03.06}\family{K04.03.10}

\begin{theorem}[Taylor--Peano formula]\label{d04:03:taylor-peano}
For an integer $k\geq1$, if $f$ is $k$ times totally differentiable at
$a$ in the sense of \cref{d04:03:higher-derivative}, then
\[
 f(a+h)=\sum_{j=0}^k\frac1{j!}D^jf(a)[h^j]+o(\norm{h}^k).
\]
In particular $C^k$ on a neighborhood suffices. For $k=0$ the assertion
is $f(a+h)=f(a)+o(1)$ and requires just continuity at $a$.
\end{theorem}
\begin{proof}
The order-one assertion is the definition. Assume the assertion of order
$k-1$ for all finite-dimensional targets. Apply it to $Df$, with target the
space of linear maps. Symmetry in \cref{d04:03:symmetry} and
\cref{d04:03:line-restriction} show that the derivative of
\[
 P(h)=\sum_{j=0}^k D^jf(a)[h^j]/j!
\]
is exactly the order-$(k-1)$ Taylor polynomial of $h\mapsto Df(a+h)$.
Thus for $q(x)=f(x)-P(x-a)$,
$\norm{Dq(a+x)}=o(\norm{x}^{k-1})$ and $q(a)=0$.
Given $\eps>0$, for all sufficiently small $x$ that norm is at most
$\eps\norm{x}^{k-1}$; at $x=0$ it is zero. On the segment from $a$ to
$a+h$, its supremum is at most $\eps\norm{h}^{k-1}$.
The segment mean-value inequality \cref{d04:02:mean-value} gives
$\norm{q(a+h)}\leq\eps\norm{h}^k$. This proves the induction without
assuming the continuity or integrability of the top derivative. The order-zero
case is the definition of continuity.
\end{proof}

\begin{theorem}[Integral remainder]\label{d04:03:taylor-integral}
Let $k\geq1$, $f\in C^k(U,\R^m)$, and $[a,a+h]\subset U$. Then
\[
 f(a+h)=\sum_{j=0}^{k-1}\frac{D^jf(a)[h^j]}{j!}
 +\frac1{(k-1)!}\int_0^1(1-t)^{k-1}D^kf(a+th)[h^k]\,\dd t.
\]
Consequently the remainder after the degree-$k$ Taylor polynomial is
\[
 R_k(a,h)=\frac1{(k-1)!}\int_0^1(1-t)^{k-1}
   \big(D^kf(a+th)-D^kf(a)\big)[h^k]\,\dd t.
\]
If $f\in C^{k+1}$ the alternative degree-$k$ remainder is
\[
 R_k(a,h)=\frac1{k!}\int_0^1(1-t)^kD^{k+1}f(a+th)[h^{k+1}]\,\dd t.
\]
For degree zero this last formula with $k=0$ is the segment FTC.
\end{theorem}
\begin{proof}
Apply the ordinary scalar integral Taylor formula from \cref{base:scalar}
to each coordinate of $t\mapsto f(a+th)$ on $[0,1]$, and use
\cref{d04:03:line-restriction}. Equivalently it follows by repeated scalar
integration by parts, whose endpoint terms are exactly the displayed
coefficients. The integral of the constant $D^kf(a)[h^k]$ against the
weight is $D^kf(a)[h^k]/k!$; subtracting it gives the second formula.
Replace $k$ by $k+1$ in the first formula to obtain the third.
All these are finite coordinate integrals, so no multiple-integration
theorem is being used.
\end{proof}
\source{HUB Appendix A.12, printed pp. 740--744 (PDF pp. 757--761),
including Theorem A12.1 and the final explicit Taylor bound, was compared
at PDF pp. 757, 758 and 761. The manuscript's integral formula is
vector-valued and separates its finite regularity hypotheses.}

\begin{corollary}[Remainder bounds and scalar Lagrange form]\label{d04:03:taylor-bounds}
Assume $k\geq0$, $f\in C^{k+1}(U,\R^m)$, and $[a,a+h]\subset U$.
If $\norm{D^{k+1}f(x)}\leq M$ on this segment, then
$\norm{R_k(a,h)}\leq M\norm{h}^{k+1}/(k+1)!$.
For scalar $f\in C^{k+1}$ and $h\ne0$ there is a $t\in(0,1)$ with
\[
 R_k(a,h)=\frac{D^{k+1}f(a+th)[h^{k+1}]}{(k+1)!}.
\]
If instead every coordinate partial of scalar $f$ of order $k+1$ has
absolute value at most $C$ on the segment, then
\[
 \abs{R_k(a,h)}\leq
 \frac{C(\sum_i\abs{h_i})^{k+1}}{(k+1)!}
 \leq\frac{Cn^{(k+1)/2}\norm{h}^{k+1}}{(k+1)!}.
\]
\end{corollary}
\begin{proof}
Use the integral norm inequality from \cref{d04:02:segment-ftc} and
$\int_0^1(1-t)^k\dd t=1/(k+1)$. The scalar Lagrange assertion is the
ordinary one-variable Taylor theorem on the line, not a vector theorem.
Expansion of the multilinear tensor in basis coordinates bounds
$\abs{D^{k+1}f(x)[h^{k+1}]}$ by
$C\sum_{i_1,\ldots,i_{k+1}}\abs{h_{i_1}\cdots h_{i_{k+1}}}
=C(\sum_i\abs{h_i})^{k+1}$. Finally apply Cauchy--Schwarz.
When $n=0$ there is no nonzero increment and the remainder is zero.
\end{proof}

\section{Uniformity and computation}
\family{K04.03.07}\family{K04.03.08}

\begin{theorem}[Compact-uniform Taylor estimate]\label{d04:03:uniform-taylor}
Let $K\subset U$ be compact and $f\in C^k(U,\R^m)$, $k\geq1$.
There are a $\rho>0$ and a function $\omega:[0,\rho]\to[0,\infty)$
with $\omega(s)\to0$ as $s\downarrow0$ such that for $a\in K$ and
$\norm{h}\leq\rho$,
\[
 \norm{R_k(a,h)}\leq\frac{\omega(\norm{h})}{k!}\norm{h}^k.
\]
The whole segment lies in $U$. If $f\in C^{k+1}$ there is also a uniform
bound $\norm{R_k(a,h)}\leq M\norm{h}^{k+1}/(k+1)!$ on those segments.
For $k=0$ the uniform conclusion is uniform continuity.
\end{theorem}
\begin{proof}
If $K$ is empty the assertions are vacuous; choose any positive $\rho$
and $\omega=0$. Otherwise choose $\rho>0$ so the closed
$\rho$-neighborhood $K_\rho$ lies in $U$
and is compact, using \cref{d04:01:distance}. The map $D^kf$ is uniformly
continuous there by \cref{d04:01:compact-continuity}. Define $\omega(s)$
as the supremum of $\norm{D^kf(x)-D^kf(y)}$ for $x,y\in K_\rho$,
$\norm{x-y}\leq s$. It is finite, and uniform continuity implies its
limit is zero. The second formula of \cref{d04:03:taylor-integral} now
gives the estimate. If the next derivative is continuous its norm has a
finite maximum on $K_\rho$, and \cref{d04:03:taylor-bounds} applies.
\end{proof}

\begin{definition}[Multi-indices and Taylor polynomials]\label{d04:03:multi-index}
A multi-index is $\alpha\in\N^n$, where here $\N$ includes zero.
Set $\abs{\alpha}=\sum_i\alpha_i$, $\alpha!=\prod_i\alpha_i!$,
$h^\alpha=\prod_ih_i^{\alpha_i}$, and let $\partial^\alpha$ denote
the coordinate partial with those multiplicities when symmetry applies.
The order-$k$ Taylor polynomial at $a$, expressed in the increment, is
$P_{f,a}^k(h)=\sum_{j=0}^kD^jf(a)[h^j]/j!$.
For $n=0$ there is a single empty multi-index, of degree zero and factorial
one; its monomial is one. This agrees with the constant polynomial on
the one-point domain.
\end{definition}

\begin{proposition}[The factorial bridge]\label{d04:03:factorial-bridge}
For a symmetric differential of order $j$,
\[
 \frac1{j!}D^jf(a)[h^j]
  =\sum_{\abs{\alpha}=j}\frac{\partial^\alpha f(a)}{\alpha!}h^\alpha.
\]
Thus the multilinear and multi-index Taylor polynomials are the same object.
\end{proposition}
\begin{proof}
Multilinear expansion sums over ordered $j$-tuples of indices. For a fixed
multiplicity vector $\alpha$ there are $j!/\prod_i\alpha_i!$ such
tuples: choose positions for the first index, then the second, and so on.
Their coefficients agree by symmetry. Grouping these terms and dividing by
$j!$ proves the identity. The empty tuple proves the order-zero case.
\end{proof}

\begin{proposition}[Finite Taylor arithmetic]\label{d04:03:taylor-arithmetic}
The Taylor polynomial of a sum is the sum of its Taylor polynomials.
For a fixed bilinear operation it is the product of the two polynomials,
with all monomials of total degree greater than $k$ discarded.
If $f(a)=b$, $f,g$ are $C^k$ near their basepoints, and $k\geq1$, then
\[
 P_{g\circ f,a}^k(h)=
 \big[P_{g,b}^k\big(P_{f,a}^k(h)-b\big)\big]_{\leq k}.
\]
Brackets mean deletion of terms of degree greater than $k$. Degree zero
just records the value of the composition. Scalar reciprocal series are
computed by the same truncated arithmetic with a nonzero constant term.
\end{proposition}
\begin{proof}
In dimension zero every polynomial function is determined by its value at
the sole increment $0$; the asserted arithmetic identities are simply
evaluation identities. Assume henceforth that $n>0$.
First, a polynomial $Q$ of degree at most $k$ with
$Q(h)=o(\norm{h}^k)$ is zero. For each fixed $v$, the polynomial $Q(tv)$
in $t$ has no nonzero coefficient: its lowest nonzero power would contradict
that remainder condition. Hence each homogeneous part vanishes at every
$v$, which means the polynomial function is zero (coordinate polynomial
coefficients can also be recovered by differentiating at zero).

Substitute the Peano expansions. Their remainders are $o(\norm{h}^k)$,
their polynomial parts stay bounded near zero, and the increment
$f(a+h)-b$ is $O(\norm{h})$ by differentiability. A bilinear product error
has a bounded factor times a small remainder. In the composition, the outer
remainder is $o(\norm{f(a+h)-b}^k)=o(\norm{h}^k)$, including zero
inner increments. Replacing the inner increment in each polynomial term by
$P_{f,a}^k(h)-b$ changes it by a sum of products containing a small
remainder, hence by $o(\norm{h}^k)$. Every discarded monomial has degree
at least $k+1$, hence is also small. Uniqueness just proved identifies the
resulting polynomial with the Taylor polynomial. Finally, for $c\ne0$,
the finite identity
$(c+u)^{-1}=\sum_{j=0}^k(-u)^j/c^{j+1}+O(\abs{u}^{k+1})$
follows by multiplying the finite geometric sum by $c+u$. Substitute
$u=q(a+h)-c=O(\norm{h})$ to obtain the reciprocal rule.
\end{proof}

\section{What regularity supplies}
\family{K04.03.09}

\begin{example}[Noncommuting partials]\label{d04:03:mixed-counterexample}
Set $f(0,0)=0$ and
$f(x,y)=xy(x^2-y^2)/(x^2+y^2)$ elsewhere. It is $C^1$: away from zero
it is rational with nonzero denominator, its first derivatives are homogeneous
of degree one and bounded on the unit circle, and at zero
$\abs{f(x,y)}\leq\abs{xy}\leq\norm{(x,y)}^2/2$ gives derivative zero.
The homogeneity bound makes those first derivatives continuous at zero.
Direct difference quotients along the axes give
$\partial_xf(0,y)=-y$ and $\partial_yf(x,0)=x$. Thus
$\partial_y\partial_xf(0)=-1$ while
$\partial_x\partial_yf(0)=1$. Existence of these iterated partials does
not give the total differentiability hypothesis of
\cref{d04:03:crossed-partials}, nor the hypotheses of a second-order Taylor
formula. The boundedness claim on the unit circle follows from continuity
and compactness, so no hidden estimate at zero was used.
\end{example}

\begin{example}[Smoothness is not analyticity]\label{d04:03:flat-example}
The function $f(x)=\exp(-1/x^2)$ for $x\ne0$ and $f(0)=0$ is smooth
with every derivative at zero equal to zero, but is positive away from zero.
Indeed each derivative away from zero is a polynomial in $1/x$ times the
same exponential. For every integer $N$, $\abs{x}^{-N}\exp(-1/x^2)\to0$:
put $t=1/x^2$ and bound $e^t$ below by a sufficiently high term of its
scalar power series. Induction, using the difference quotient at zero and
that estimate with one extra power, proves existence and continuity of every
derivative with value zero there. Every finite Taylor polynomial vanishes,
but the function does not coincide locally with its Taylor series.
\end{example}

\begin{remark}[Finite order and scalar specialization]\label{d04:03:scope}
The pointwise Peano theorem uses iterated total differentiability; the segment
integral theorem uses continuous derivatives on the whole segment; uniform
estimates use compact neighborhoods. These are separate assertions. With
$n=m=1$ all formulas reduce to the admitted one-variable Taylor interface.
The order-zero case is continuity, not a differentiability assumption. An
exact vector Lagrange remainder would already imply the false vector
mean-value equality exhibited in \cref{d04:02:mean-value}.
\end{remark}

% END ORIGINAL INPUT: chapters/03.tex


% BEGIN ORIGINAL INPUT: chapters/04.tex
\chapter{Local inverse, implicit, and constant-rank theorems}
\label{ch:d04-04}

\section{Stable inverses and contractions}
\family{K04.04.01}\family{K04.04.02}

\begin{theorem}[Quantitative stability of matrix inversion]\label{d04:04:inverse-stability}
Let $A:\R^n\to\R^n$ be invertible. If
$q=\norm{A^{-1}(B-A)}<1$, then $B$ is invertible and
\[
 B^{-1}=\sum_{j=0}^\infty\big(-A^{-1}(B-A)\big)^jA^{-1},\qquad
 \norm{B^{-1}}\leq\frac{\norm{A^{-1}}}{1-q},
\]
\[
 \norm{B^{-1}-A^{-1}}\leq
 \frac{\norm{A^{-1}}^2\norm{B-A}}{1-q}.
\]
Thus $\GL(n,\R)$ is open in the matrix space and inversion is continuous.
For $n=0$, its sole element is invertible with determinant one, and the
claims are interpreted in the one-point matrix space.
\end{theorem}
\begin{proof}
Set $E=A^{-1}(B-A)$. The partial sums of $\sum(-E)^j$ are Cauchy in
operator norm because their tails are bounded by the tails of $\sum q^j$.
Finite-dimensional completeness from \cref{d04:01:coordinate-estimates,d04:01:operator-estimates}
gives a limit $S$. Multiplication is continuous,
and the finite geometric identities on both sides show
$(I+E)S=S(I+E)=I$. As $B=A(I+E)$, its inverse is $SA^{-1}$.
Summing the norm bounds gives the first estimate. Summing only the terms
with $j\geq1$ gives
$\norm{B^{-1}-A^{-1}}\leq q\norm{A^{-1}}/(1-q)$ and hence the stated
second estimate. If $\norm{B-A}<1/(2\norm{A^{-1}})$ in positive dimension,
then $q<1/2$; this is a neighborhood of invertible matrices and proves
continuity at $A$. In dimension zero all maps and estimates are the unique
zero ones, while the identity of the zero space is also that unique map.
\end{proof}

\begin{proposition}[Smooth matrix inversion and its differential]\label{d04:04:matrix-inversion}
The map $A\mapsto A^{-1}$ on $\GL(n,\R)$ is $C^\infty$, and
\[
 D(\mathrm{inv})(A)H=-A^{-1}HA^{-1}.
\]
\end{proposition}
\begin{proof}
For $n>0$, the adjugate identity from \cref{base:linear} expresses every
entry of $A^{-1}$ as a signed cofactor, with the row and column transposed,
divided by $\det A$. Determinants and cofactors are finite polynomials,
and repeated use of \cref{d04:02:cr-calculus} makes the quotient smooth
on the open nonzero-determinant set. To identify its derivative without
any cofactor convention ambiguity, write $B=A+H$ and use
\[
 B^{-1}=A^{-1}-A^{-1}HA^{-1}
                 +A^{-1}HA^{-1}HB^{-1}.
\]
The identity follows twice from $B^{-1}=A^{-1}-A^{-1}HB^{-1}$.
The last term is $O(\norm{H}^2)$ by
\cref{d04:04:inverse-stability}. The displayed linear term is therefore
the total derivative. Dimension zero is immediate.
\end{proof}

\begin{theorem}[Contraction principle with error bound]\label{d04:04:contraction}
Let $S$ be a nonempty complete subset of a Euclidean space with its induced
metric, and $T:S\to S$ satisfy
$\norm{T(x)-T(y)}\leq q\norm{x-y}$ for a fixed $0\leq q<1$.
Then $T$ has exactly one fixed point $x_*$. Starting from any $x_0\in S$,
$x_{j+1}=T(x_j)$ converges to it and
\[
 \norm{x_j-x_*}\leq\frac{q^j}{1-q}\norm{x_1-x_0},\qquad
 \norm{x-x_*}\leq\frac{\norm{x-T(x)}}{1-q}\quad(x\in S).
\]
Closed subsets of $\R^n$, in particular closed balls, are complete.
\end{theorem}
\begin{proof}
Induction bounds successive increments by $q^j\norm{x_1-x_0}$.
The triangle inequality and the geometric sum then bound every tail, so
$(x_j)$ is Cauchy and has a limit in $S$. The contraction inequality implies
continuity, giving $T(x_*)=x_*$. Two fixed points would satisfy
$\norm{x_*-y_*}\leq q\norm{x_*-y_*}$ and hence coincide. Passing to
the limit in the tail bound gives the first estimate. For any $x$,
\[
 \norm{x-x_*}\leq\norm{x-T(x)}+q\norm{x-x_*},
\]
which proves the second. A Cauchy sequence in a closed subset converges
in $\R^n$ by completeness and remains in the subset by closedness.
The convention $q^0=1$ includes the case $q=0$.
\end{proof}

\section{The inverse function theorem}
\family{K04.04.03}\family{K04.04.04}

\begin{theorem}[Local inverse with finite regularity]\label{d04:04:inverse-function}
Let $U\subset\R^n$ be open, $f\in C^r(U,\R^n)$ with
$r\in\{1,2,\ldots\}\cup\{\infty\}$, and $a\in U$. If $Df(a)$
is invertible, there are open neighborhoods $X$ of $a$ and $Y$ of $f(a)$
such that $f:X\to Y$ is a bijection and its inverse $g:Y\to X$ is $C^r$.
Every $Df(x)$ on $X$ is invertible and
\[
 Dg(y)=\big(Df(g(y))\big)^{-1}.
\]
The inverse is locally unique with values in $X$ and can be chosen Lipschitz.
More quantitatively, if on a closed ball $\overline B(a,\rho)\subset U$
\[
 \norm{I-A^{-1}Df(x)}\leq q<1,\qquad A=Df(a),
\]
then every $y$ with $\norm{A^{-1}(y-f(a))}<(1-q)\rho$ has exactly
one preimage in that ball, and it lies in the open ball. For two such values
whose solutions lie there,
\[
 \norm{g(y)-g(z)}\leq\frac{\norm{A^{-1}(y-z)}}{1-q}.
\]
\end{theorem}
\begin{proof}
Translations and multiplication in the target by $A^{-1}$ reduce the
construction to $a=0$, $f(0)=0$, $Df(0)=I$. By continuity of $Df$,
choose $\rho>0$ so the closed ball lies in $U$ and
$\norm{Df(x)-I}\leq q<1$ there. The segment mean-value estimate for
$f-\id$ gives
\[
 \norm{(f(x)-x)-(f(z)-z)}\leq q\norm{x-z}.
\]
Consequently $f$ is injective on the ball and
\[
 (1-q)\norm{x-z}\leq\norm{f(x)-f(z)}.
\]
For $y$ with $\norm{y}<(1-q)\rho$, set $T_y(x)=x-f(x)+y$.
It is a contraction with ratio $q$ on the closed ball, and
$\norm{T_y(x)}\leq q\rho+\norm{y}\leq\rho$.
By \cref{d04:04:contraction} it has a unique fixed point, which solves
$f(x)=y$. Its norm is at most $\norm{y}/(1-q)<\rho$, so it is interior.
The preceding lower bound proves the inverse Lipschitz estimate. Choose a
smaller open ball $Y=B(0,\sigma)$ with $0<\sigma<(1-q)\rho$, and set
$X=B(0,\rho)\cap f^{-1}(Y)$. It is open, contains zero, and $f:X\to Y$
is bijective. Undoing the normalizations proves the asserted construction
and the quantitative estimates. Notice that the invariant closed ball was
used for existence, while the final inverse has open domain and range.

By \cref{d04:04:inverse-stability}, $Df(x)$ is invertible throughout $X$.
Fix $y=f(x)\in Y$ and let $k\to0$ with $y+k\in Y$.
Put $h=g(y+k)-g(y)$. Lipschitz continuity gives $h=O(\norm{k})$.
Differentiability of $f$ at $x$ says
$k=Df(x)h+r(h)$, $r(h)=o(\norm{h})$. Solving gives
\[
 h=Df(x)^{-1}k-Df(x)^{-1}r(h),
\]
whose error is $o(\norm{k})$, including $h=0$. This proves the inverse
derivative formula. Continuity of $Df$, of $g$, and of matrix inversion
makes $Dg$ continuous, so $g$ is $C^1$.

If $f\in C^r$ with $r\geq2$, bootstrap the formula
$Dg=\mathrm{inv}\circ Df\circ g$. If $g$ is $C^j$ and $j\leq r-1$,
then $Df$ is at least $C^j$, so \cref{d04:02:cr-calculus} and
\cref{d04:04:matrix-inversion} imply $Dg\in C^j$ and hence
$g\in C^{j+1}$. Induction reaches $r$; applying every finite step proves
the smooth case. Any other solution with values in $X$ agrees with $g$
by injectivity. In dimension zero the unique map between the singleton
neighborhoods is its own inverse, and all derivative assertions hold.
\end{proof}
\source{Comparison: SPI Theorem 2-11, printed pp. 35--39 (PDF pp. 48--52),
and Addenda item 1, printed p. 145 (PDF p. 158, also inspected visually),
were checked. The proof here uses a contraction rather than Spivak's
minimum-distance argument. HUB Appendix A.7, printed pp. 724--728
(PDF pp. 741--745), was checked for the separate continuity,
differentiability and openness issues; its quantitative Lipschitz-derivative
version is not an assumption here.}

\begin{corollary}[A global inverse when injectivity is supplied]\label{d04:04:injective-local-diffeomorphism}
If $f\in C^r(U,\R^n)$ is injective and $Df(x)$ is invertible for every
$x\in U$, then $f(U)$ is open and $f:U\to f(U)$ has a $C^r$ inverse.
\end{corollary}
\begin{proof}
The local inverse theorem at each point gives an open neighborhood in the
image. Their union is $f(U)$. Global injectivity makes their local inverses
agree wherever their targets overlap. The resulting inverse is $C^r$ since
that property is local, and is the set-theoretic inverse by construction.
\end{proof}

\section{Implicit equations and several constraints}
\family{K04.04.05}\family{K04.04.06}\family{K04.04.07}

\begin{theorem}[Implicit function theorem]\label{d04:04:implicit-function}
Let $G:W\subset\R^p\times\R^q\to\R^q$ be $C^r$, $W$ open,
$r\geq1$ finite or infinite. Suppose $G(a,b)=c$ and
$D_yG(a,b):\R^q\to\R^q$ is invertible. There are open neighborhoods
$A$ of $a$ and $B$ of $b$ and a unique $C^r$ map $\varphi:A\to B$
such that, for $(x,y)\in A\times B$,
\[
 G(x,y)=c\quad\Longleftrightarrow\quad y=\varphi(x).
\]
It has $\varphi(a)=b$ and
\[
 D\varphi(x)=-\big(D_yG(x,\varphi(x))\big)^{-1}
                         D_xG(x,\varphi(x)).
\]
In matrix notation these are the solved-variable and parameter blocks of
the same total derivative, not separately defined Jacobians.
\end{theorem}
\begin{proof}
Define $F(x,y)=(x,G(x,y))$. Its derivative has block matrix
\[
 \begin{pmatrix}I&0\\D_xG&D_yG\end{pmatrix},
\]
which is invertible at $(a,b)$: solving for its input first gives the
parameter component and then uses $D_yG^{-1}$ for the remaining component.
By \cref{d04:04:inverse-function}, $F$ has a $C^r$ inverse on an open
neighborhood $V$ of $(a,c)$, taking values in an open neighborhood $O$
of $(a,b)$. The first coordinate of that inverse must be $x$, so write
it as $(x,H(x,z))$. Choose a product $A_0\times B\subset O$ about
$(a,b)$. Shrink $A\subset A_0$ so $(x,c)\in V$ and $H(x,c)\in B$
for every $x\in A$; continuity and openness permit this. Set
$\varphi(x)=H(x,c)$. Its graph solves the equation. Conversely any solution
in $A\times B\subset O$ maps to $(x,c)$ and must be that inverse value,
which proves uniqueness on the entire chosen product, not merely on the
graph. Shrinking $O$ initially keeps $D_yG$ invertible throughout by
\cref{d04:04:inverse-stability}. Differentiate
$G(x,\varphi(x))=c$ with the chain rule to get
$D_xG+D_yG D\varphi=0$ and solve the block equation. Regularity follows
from that of the inverse. If $q=0$ this is the unique map into $\R^0$;
if $p=0$ it asserts the local uniqueness of the solution point. The same
construction covers both cases with the empty blocks.
\end{proof}
\source{SPI Theorem 2-12 and its derivative calculation, printed pp. 41--42
(PDF pp. 54--55), and Theorem 2-13, printed p. 43 (PDF p. 56), were
checked. The explicit shrinking above supplies product neighborhoods and
local uniqueness; finite $C^r$ follows from the inverse bootstrap.}

\section{Rank normal form and regular geometry}
\family{K04.04.08}\family{K04.04.09}\family{K04.04.10}\family{K04.04.11}

\begin{theorem}[Constant-rank normal form]\label{d04:04:constant-rank}
Let $f:U\subset\R^n\to\R^m$ be $C^r$, $r\geq1$, and suppose
$\rank Df(x)=k$ on a neighborhood of $a$. Then $0\leq k\leq\min(n,m)$,
and there are local $C^r$ coordinate diffeomorphisms $\Phi$ in the source
and $\Psi$ in the target, taking $a$ and $f(a)$ to zero, for which
\[
 (\Psi\circ f\circ\Phi^{-1})(u_1,\ldots,u_n)
       =(u_1,\ldots,u_k,0,\ldots,0).
\]
The rank hypothesis is required on a neighborhood, not merely at $a$.
\end{theorem}
\begin{proof}
Translate both basepoints to zero. If $k=0$, the derivative vanishes on a
small ball, so $f$ is constant there by \cref{d04:02:zero-derivative};
translations give the asserted zero form. Otherwise choose a nonzero
$k$-minor at zero and permute input and output coordinates to make it the
first $k$ rows and columns. Write $x=(x',x'')$ and $f=(f',f'')$.
The map $\Phi(x)=(f'(x),x'')$ has invertible derivative at zero. Its local
$C^r$ inverse exists by \cref{d04:04:inverse-function}. In the new source
coordinates, on a sufficiently small product box,
\[
 f\circ\Phi^{-1}(u',u'')=(u',H(u',u'')).
\]
Its derivative is
$\left(\begin{smallmatrix}I&0\\D_{u'}H&D_{u''}H\end{smallmatrix}\right)$.
Invertible coordinate changes preserve rank by the chain rule and linear
algebra. Since this matrix has rank exactly $k$, its last block
$D_{u''}H$ is zero: subtract the first $k$ column contributions from the
remaining ones, or use the displayed block row reduction. For fixed $u'$,
the second-variable box is convex, so \cref{d04:02:zero-derivative} gives
$H(u',u'')=h(u')=H(u',0)$. This $h$ is $C^r$ and $h(0)=0$.
The target transformation
$\Psi(y',y'')=(y',y''-h(y'))$ is a $C^r$ diffeomorphism on a target
neighborhood with $y'$ in that box; its inverse is
$(v',v'')\mapsto(v',v''+h(v'))$. It yields the asserted normal form.
The cases $k=n$ or $k=m$ simply omit the appropriate empty block.
This proof uses constancy of rank exactly in forcing $D_{u''}H=0$.
\end{proof}

\begin{theorem}[Regular levels, graph charts, and tangent kernels]\label{d04:04:regular-level}
Let $g:U\subset\R^n\to\R^q$ be $C^r$, $r\geq1$, and let
$M=g^{-1}(c)$. If $Dg(a)$ is surjective at $a\in M$, then $q\leq n$
and, after permuting coordinates, there is a neighborhood of $a$ in which
\[
 M=\{(u,\varphi(u)):u\in A\},\qquad A\subset\R^{n-q}\text{ open},
\]
where $\varphi$ is $C^r$. The graph map
$\theta(u)=(u,\varphi(u))$ is a one-to-one regular parametrization,
homeomorphic onto that relative neighborhood, and
\[
 D\theta(u)\R^{n-q}=\ker Dg(\theta(u)).
\]
The latter space is exactly the set of velocities at $\theta(u)$ of $C^1$
curves in $M$ defined on an open interval about zero. If every point of $M$
is regular, these graph neighborhoods make $M$ an embedded $C^r$ manifold
of dimension $n-q$ in the Euclidean sense used in \cref{ch:d04-14}.
\end{theorem}
\begin{proof}
Surjectivity supplies $q$ independent columns of $Dg(a)$ by
\cref{base:linear}. Permute them into the solved variables and apply
\cref{d04:04:implicit-function}. Projection onto $u$ is a continuous inverse
of the graph map; $D\theta(u)v=(v,D\varphi(u)v)$ is injective.
Differentiating $g\circ\theta=c$ places its image in the kernel.
Conversely, for $(v,w)\in\ker Dg$, the block equation
$D_ug\,v+D_yg\,w=0$ forces $w=D\varphi(u)v$. This proves equality
and also the dimension count by rank-nullity. Every $C^1$ curve $\gamma$
in the level has $Dg\,\gamma'(0)=0$ by the chain rule. Conversely each
kernel vector $D\theta(u)v$ is realized by
$t\mapsto\theta(u+tv)$ for $t$ sufficiently small of both signs.
Transition maps between overlapping graph coordinates are projections
composed with graph maps, hence $C^r$; their inverses are the reverse
transition maps. This is precisely the required embedded-chart assertion,
without assuming a manifold theorem. When $q=0$, the level is $U$;
when $q=n$, it is locally one point with zero tangent space. Empty levels
have no pointwise assertion to check.
\end{proof}

\begin{corollary}[Submersions and immersions]\label{d04:04:submersion-immersion}
A $C^r$ map is a submersion at $a$ if $Df(a)$ is surjective; equivalently
its scalar component differentials are linearly independent. Near $a$ it
has the projection normal form $(u,v)\mapsto u$. If it is a submersion
everywhere, it is an open map, but need not map onto the whole target.
A $C^r$ map is an immersion at $a$ if $Df(a)$ is injective. Near $a$ it
has the inclusion normal form $u\mapsto(u,0)$ and is an embedding on a
neighborhood of $a$ (a homeomorphism onto its image with the relative topology).
An immersion need not be globally injective or globally an embedding.
\end{corollary}
\begin{proof}
A full-rank minor stays nonzero nearby by continuity of its determinant;
surjectivity or injectivity is therefore a neighborhood condition. Its rank
is then the maximal possible rank, so \cref{d04:04:constant-rank} applies.
Row independence is exactly independence of the scalar covectors in the
column-matrix representation. Projection sends an open neighborhood to an
open neighborhood: for each point a small product ball fits inside the
domain, and its projection contains an open ball. Coordinate diffeomorphisms
preserve openness, proving the open-map assertion. The inclusion map is
one-to-one with continuous inverse the coordinate projection; composing
with the coordinate diffeomorphisms proves the embedding assertion.
\end{proof}

\section{Local conclusions and their limits}
\family{K04.04.12}

\begin{example}[Several distinct failures]\label{d04:04:counterexamples}
The map $f(x,y)=(e^x\cos y,e^x\sin y)$ has determinant $e^{2x}>0$
everywhere, but is $2\pi$-periodic in $y$. It has local inverses and no
global inverse on $\R^2$. Its image omits zero. Already $x\mapsto e^x$
is an everywhere submersion with proper open image $(0,\infty)$.
The immersion $t\mapsto(\cos t,\sin t)$ is locally an embedding but is
periodic. These statements follow from the displayed formulas and their
derivatives.

The function $x\mapsto x^3$ is bijective although its derivative vanishes
at zero; its inverse cannot be differentiable there, since the chain rule
for two differentiable inverses would give $1=0\cdot Dg(0)$. Rank at a
single point is not a normal-form hypothesis: $x\mapsto x^2$ has rank
zero at zero but is not constant on any neighborhood, whereas the rank-zero
normal form would make it constant.

Finally define $f(0)=0$, $f(x)=x/2+x^2\sin(1/x)$ for $x\ne0$.
It is differentiable at zero with derivative $1/2$, yet is not one-to-one
on any neighborhood of zero. For $x\ne0$ its derivative is
$1/2+2x\sin(1/x)-\cos(1/x)$; at arbitrarily small positive points with
$1/x$ an even multiple of $\pi$ it is $-1/2$, and at odd multiples it
is $3/2$. A continuous injective function on an interval must be strictly
monotone: otherwise three ordered points giving a reversal yield, by the
intermediate value theorem on the two subintervals, a repeated value.
The positive derivative at zero precludes decreasing behavior, and the
negative derivatives arbitrarily nearby preclude increasing behavior.
Thus the derivative continuity assumption in the inverse theorem cannot
simply be deleted.
\end{example}
\source{SPI Problems 2-38(b) and 2-39, printed p. 40 (PDF p. 53), were
checked for the exponential-polar and derivative-discontinuity examples.
The constant-rank proof and regular-level tangent bridge here are supplied
in full rather than attributed to an unchecked source theorem.}

% END ORIGINAL INPUT: chapters/04.tex


% BEGIN ORIGINAL INPUT: chapters/05.tex
\chapter{Extrema, constraints, and a bounded Morse lemma}
\label{ch:d04-05}

\section{Existence and interior tests}
\family{K04.05.01}\family{K04.05.02}\family{K04.05.03}
\family{K04.05.04}\family{K04.05.05}\family{K04.05.06}

\begin{definition}[Relative extrema and stationary points]\label{d04:05:extrema}
Let $f:S\to\R$ and $a\in S$. It has a local minimum relative to $S$
at $a$ if $f(x)\geq f(a)$ for $x\in S$ in some ambient ball about $a$;
the minimum is strict if the inequality is strict for $x\ne a$ in that
ball. A global minimum uses all $x\in S$. Maxima reverse the inequalities.
A stationary or critical point of a differentiable function on an open set
is a point where $Df(a)=0$. A saddle point here means that every
neighborhood of $a$ contains both a point with value greater than $f(a)$
and a point with value smaller than $f(a)$. These notions always retain
the indicated feasible set; in particular a boundary extremum need not be
an ambient stationary point.
\end{definition}

\begin{theorem}[Existence and first-order necessity]\label{d04:05:existence-fermat}
A continuous real function on a nonempty compact feasible set attains its
global minimum and maximum. If $f$ is differentiable on an open neighborhood
of $a$ and has an unconstrained local extremum there, then $Df(a)=0$.
More generally, at a relative local extremum its derivative vanishes on
every velocity of a two-sided differentiable feasible curve through $a$.
For a feasible curve defined only for $t\geq0$, a minimum gives
$Df(a)\gamma'(0)\geq0$, with the opposite sign for a maximum.
\end{theorem}
\begin{proof}
Existence is \cref{d04:01:compact-continuity}. Restrict to each line
$a+te_j$ for the unconstrained assertion. The scalar Fermat theorem gives
$\partial_jf(a)=0$ for all $j$ and hence $Df(a)=0$ by
\cref{d04:02:derivative-bridge}. For a two-sided feasible curve apply scalar
Fermat to $f\circ\gamma$ and the chain rule. For a one-sided curve at
a minimum, $[f(\gamma(t))-f(a)]/t\geq0$ for $t>0$; its limit is the
asserted inequality. The maximum assertion follows on replacing $f$ by $-f$.
\end{proof}

\begin{theorem}[Second-order interior tests]\label{d04:05:second-order}
Suppose $f$ is twice totally differentiable at an interior point $a$.
At a local minimum, $Df(a)=0$ and $D^2f(a)[v,v]\geq0$ for all $v$;
at a local maximum the quadratic form is nonpositive. Conversely, if
$Df(a)=0$ and the Hessian is positive definite, $a$ is a strict local
minimum. In positive dimension there are $c,\rho>0$ with
\[
 f(a+h)\geq f(a)+c\norm{h}^2\quad(\norm{h}<\rho).
\]
Negative definiteness gives the corresponding strict maximum. If the
Hessian is indefinite, $a$ is a saddle point. Semidefiniteness alone is
not a sufficient test.
\end{theorem}
\begin{proof}
The first-order conclusion is \cref{d04:05:existence-fermat}. Taylor--Peano
on $tv$ then gives
$f(a+tv)-f(a)=t^2D^2f(a)[v,v]/2+o(t^2)$.
At a minimum divide by $t^2>0$ and pass to the limit. For sufficiency,
the positive continuous form $v\mapsto D^2f(a)[v,v]$ has a positive
minimum $\eta$ on the unit sphere by compactness. Consequently
$D^2f(a)[h,h]\geq\eta\norm{h}^2$. The Peano error is smaller than
$\eta\norm{h}^2/4$ for small $h$, yielding $c=\eta/4$.
Replacing $f$ by $-f$ gives the negative case. In the indefinite case,
if $Df(a)=0$, choose fixed vectors $u,v$ with positive and negative
quadratic values; the line expansions give both signs of $f-f(a)$
arbitrarily near $a$. If $Df(a)\ne0$, choose $v$ with $Df(a)v>0$;
the first-order expansions along $tv$ and $-tv$ for $t>0$ already give
both signs, irrespective of the Hessian.
For failure in the semidefinite case, $x^4$, $-x^4$, and $x^3$ have the
same zero Hessian at zero and respectively a strict minimum, strict maximum,
and saddle. More visibly a nonzero semidefinite Hessian occurs at zero for
$x^2+y^4$, $x^2-y^4$, and $x^2$; these have respectively a strict minimum,
a saddle, and a non-strict minimum. In dimension zero the sole point is
both a strict relative minimum and maximum by the empty comparison set;
quadratic-growth assertions have no nonzero increment to check.
\end{proof}
\source{HUB 3.6 and the proof on printed p. 347 (PDF p. 364, inspected
visually) were checked, including both directions of the saddle argument.
The manuscript uses twice total differentiability and its proved Peano
formula rather than silently requiring all functions to be smooth.}

\section{Regular equality constraints}
\family{K04.05.07}\family{K04.05.08}

\begin{theorem}[Lagrange multipliers]\label{d04:05:lagrange}
Let $f:U\to\R$ and $g:U\to\R^q$ be $C^1$, $U\subset\R^n$ open,
and let $a$ be a relative local extremum of $f$ on $M=g^{-1}(c)$.
If $Dg(a)$ is surjective, there is a unique covector
$\lambda\in(\R^q)^*$, represented by scalars $\lambda_1,\ldots,\lambda_q$,
such that
\[
 Df(a)=\lambda\circ Dg(a),\qquad
 \nabla f(a)=\sum_{i=1}^q\lambda_i\nabla g_i(a).
\]
The fixed sign convention is
$L_\lambda(x)=f(x)-\sum_i\lambda_i g_i(x)$, so $DL_\lambda(a)=0$.
These are necessary equations, not an existence or sufficiency theorem
for extrema.
\end{theorem}
\begin{proof}
By \cref{d04:04:regular-level}, every vector in $\ker Dg(a)$ is the
velocity of a two-sided $C^1$ feasible curve. Thus
\cref{d04:05:existence-fermat} says that $Df(a)$ annihilates this kernel.
Choose a linear right inverse $R$ of the surjective linear map $Dg(a)$,
using finite-dimensional linear algebra. Put $\lambda=Df(a)R$.
For every $v$, $v-RDg(a)v$ belongs to the kernel, so
$Df(a)v=\lambda Dg(a)v$. If two such covectors existed, their difference
would vanish on the image of $Dg(a)$, which is all of $\R^q$.
The gradient formula is the transpose of this same row-map equation under
\cref{d04:02:derivative-bridge}. For $q=0$ it reduces to Fermat; for
$q=n$ the feasible point is locally isolated, and the unique multiplier
still exists because $Dg(a)$ is invertible.
\end{proof}

\begin{theorem}[Tangent second variation]\label{d04:05:constrained-second-variation}
Let $f,g$ be $C^2$ and $Dg(a)$ surjective at $a\in g^{-1}(c)$.
Suppose $Df(a)=\lambda Dg(a)$ and use $L_\lambda=f-\lambda g$.
At a constrained local minimum,
\[
 D^2L_\lambda(a)[v,v]\geq0\qquad(v\in\ker Dg(a)),
\]
and a maximum gives the opposite inequality. If this restricted quadratic
form is positive definite, $a$ is a strict constrained local minimum, with
quadratic growth relative to $M$. Negative definiteness gives a strict
maximum. If it takes both signs on the tangent kernel, there are feasible
points of both higher and lower value in every neighborhood of $a$.
\end{theorem}
\begin{proof}
Take the $C^2$ graph chart $\theta$ of \cref{d04:04:regular-level} with
$\theta(u_0)=a$, and set $F=f\circ\theta$. Its derivative at $u_0$
is zero because $D\theta$ maps onto $\ker Dg(a)$. Differentiating the
chain rule by the bilinear rule gives, for $w$ in the coordinate space and
$v=D\theta(u_0)w$,
\[
 D^2F(u_0)[w,w]=D^2f(a)[v,v]+Df(a)D^2\theta(u_0)[w,w].
\]
Applying the same calculation to $g\circ\theta=c$ gives
$D^2g(a)[v,v]+Dg(a)D^2\theta(u_0)[w,w]=0$.
Multiply by $\lambda$ and subtract, obtaining
\[
 D^2F(u_0)[w,w]=D^2L_\lambda(a)[v,v].
\]
The map $D\theta(u_0)$ is an isomorphism onto the tangent kernel.
The necessary, definite, and indefinite assertions therefore follow from
\cref{d04:05:second-order} applied to $F$. For ambient quadratic growth,
the graph map is locally Lipschitz by \cref{d04:02:zero-derivative}, so
$\norm{\theta(u)-a}\leq C\norm{u-u_0}$; growth in $u$ yields growth
in ambient distance. If the coordinate dimension is zero, the level is
locally a single point, and both strict relative extrema hold vacuously.
\end{proof}

\begin{remark}[The multiplier is fixed during variation]\label{d04:05:fixed-multiplier}
In the second variation $\lambda$ is the value obtained at $a$; it is held
constant while differentiating in $x$. The relevant Hessian is
$H_f(a)-\sum_i\lambda_iH_{g_i}(a)$, restricted to $\ker Dg(a)$.
Using $H_f$ alone ignores the acceleration of feasible curves. Testing
this Lagrangian Hessian on the whole ambient space imposes an unnecessary
condition in the normal directions.
\end{remark}

\section{Boundary constraints and active inequalities}
\family{K04.05.10}

\begin{theorem}[Active-constraint first-order rule]\label{d04:05:active-constraints}
Let $f$, $g:U\to\R^q$, and $h:U\to\R^s$ be $C^1$, and define
\[
 S=\{x\in U:g(x)=c,\ h_i(x)\geq0\ (1\leq i\leq s)\}.
\]
At $a\in S$ put $I=\{i:h_i(a)=0\}$. Assume the differentials of all
$g_j$ and of all $h_i$, $i\in I$, are linearly independent at $a$.
If $a$ is a local minimum, there are unique multipliers $\lambda_j$ and
$\mu_i$ for the active constraints such that
\[
 Df(a)=\lambda Dg(a)+\sum_{i\in I}\mu_iDh_i(a),\qquad \mu_i\geq0.
\]
Set $\mu_i=0$ for inactive constraints; then $\mu_ih_i(a)=0$ for every
$i$. The Lagrangian is $L=f-\lambda g-\sum_i\mu_i h_i$.
For a maximum, apply this assertion to $-f$.
\end{theorem}
\begin{proof}
The combined map $x\mapsto(g(x), (h_i(x))_{i\in I})$ is a submersion.
Complete its independent rows with coordinate rows to obtain a local $C^1$
coordinate diffeomorphism
\[
 x\longmapsto(u,v,w),\qquad
 u=g(x)-c,\quad v=(h_i(x))_{i\in I},
\]
by the inverse theorem, exactly as in the proof of the constant-rank theorem.
Shrink its target to a product neighborhood about $(0,0,w_0)$.
All inactive inequalities stay strictly positive by continuity.
The feasible set in these coordinates is precisely $u=0$, $v_i\geq0$.
Write $F$ for $f$ in these coordinates. At a minimum its derivative in
each free $w$ direction is zero, while its derivative in each increasing
$v_i$ direction is nonnegative by the one-sided assertion of
\cref{d04:05:existence-fermat}. The derivatives in the $u$ directions are
unrestricted. Thus
$DF=\lambda\,\dd u+\sum_i\mu_i\,\dd v_i$ with $\mu_i\geq0$.
Pulling this covector equation back gives the result. Independence of the
constraint covectors proves uniqueness. The case of no active inequalities
is precisely the equality rule.
\end{proof}

\begin{theorem}[Critical-cone second-order tests]\label{d04:05:critical-cone}
Under the active-independence hypothesis above, suppose $f,g,h$ are $C^2$
and the first-order multiplier equations hold. Define
\[
 C=\{v:Dg(a)v=0,\ Dh_i(a)v\geq0\ (i\in I),\ Df(a)v=0\}.
\]
At a constrained minimum, $D^2L(a)[v,v]\geq0$ for every $v\in C$.
If $D^2L(a)[v,v]>0$ for every nonzero $v\in C$, then $a$ is a strict
local minimum on $S$, and there exist $c_0,\rho>0$ such that
$f(x)\geq f(a)+c_0\norm{x-a}^2$ for $x\in S\cap B(a,\rho)$.
This includes the case $C=\{0\}$.
\end{theorem}
\begin{proof}
Use the $C^2$ constraint coordinates from the preceding proof. For
$v\in C$, let its coordinate velocity be $(0,\beta,\zeta)$, with
$\beta_i=Dh_i(a)v\geq0$. The coordinate line
$(0,t\beta,w_0+t\zeta)$ for small $t\geq0$ pulls back to a feasible
$C^2$ curve $\gamma$ with $\gamma'(0)=v$. Along it the equalities are
constant and each active $h_i$ is exactly $t\beta_i$, so their second
derivatives are zero. As in the equality second-variation calculation,
\[
 (f\circ\gamma)''(0)=D^2L(a)[v,v].
\]
The first derivative is $Df(a)v=0$. A one-sided minimum with zero first
derivative has nonnegative second derivative by its scalar order-two
Taylor formula. This proves necessity.

For sufficiency suppose the claimed quadratic growth fails. Choose feasible
$x_k\ne a$ with $t_k=\norm{x_k-a}<1/k$ and
$f(x_k)-f(a)<t_k^2/k$. A subsequence of
$v_k=(x_k-a)/t_k$ converges to a unit vector $v$. First-order expansions
of the equalities and active inequalities show
$Dg(a)v=0$ and $Dh_i(a)v\geq0$. The expansion of $f$, and the preceding
upper bound divided by $t_k$, show $Df(a)v\leq0$; the multiplier formula
and nonnegative $\mu_i$ give the opposite inequality. Thus $v\in C$.
Since $DL(a)=0$, Taylor's formula gives
\[
 f(x_k)-f(a)
 =\tfrac12D^2L(a)[x_k-a,x_k-a]+o(t_k^2)
                  +\sum_{i\in I}\mu_i h_i(x_k)
 \geq\tfrac12t_k^2D^2L(a)[v_k,v_k]+o(t_k^2).
\]
Here the equality constraints cancel between $x_k$ and $a$, inactive
multipliers vanish, and active $h_i(a)=0$. Divide by $t_k^2$ and pass
to the limit. The upper bound tends to zero while the lower limit is
$D^2L(a)[v,v]/2>0$, a contradiction. If $C=\{0\}$ the earlier production
of a unit $v\in C$ is already impossible. This proves quadratic growth
and hence strictness.
\end{proof}

\begin{example}[Compact, singular, and active boundary cases]\label{d04:05:worked-constraints}
For $n\geq1$ and $v\ne0$, minimize and maximize $f(x)=\ip{v}{x}$ on
the unit sphere. The feasible set is nonempty compact. For
$g(x)=\norm{x}^2$, its derivative $2\ip{x}{\cdot}$ is surjective on
the sphere, so $v=2\lambda x$. Together with $\norm{x}=1$ this gives
$x=\pm v/\norm{v}$ and $\lambda=\pm\norm{v}/2$.
Cauchy--Schwarz gives $-\norm{v}\leq f(x)\leq\norm{v}$, attained
exactly at these two points. On $x^\perp$ the Lagrangian Hessian is
$-2\lambda I$, giving the same classification and its quadratic growth
when tangent directions exist. For $n=1$ the sphere consists of two
isolated points; local strictness is relative and immediate. For $v=0$
every feasible point is a global minimum and maximum.

Singular constraints require direct examination. The level
$g(x,y)=x^2+y^2=0$ consists only of the origin. Every function has both
relative extrema there, but for $f(x,y)=x$ the multiplier equation would
say $(1,0)=\lambda(0,0)$, which is impossible. Even a regular-looking
feasible set can have a singular defining equation: $g(x,y)=y^2=0$
is the horizontal axis, and $f(x,y)=y$ is constant on it, while its
nonzero gradient cannot be a multiple of $Dg=0$ on that axis.

On the closed unit disk minimize $f(x,y)=x$. The minimum is $(-1,0)$,
the maximum $(1,0)$, by $-1\leq x\leq1$ and the feasibility equation
at equality. Neither is an ambient stationary point. With
$h=1-x^2-y^2\geq0$, the minimum satisfies
$Df=\mu Dh$ with $\mu=1/2\geq0$; the maximum obeys the same rule
for $-f$. At a corner of $[0,1]^2$, the minimum of $x+y$ at $(0,0)$
has active constraints $h_1=x$, $h_2=y$ and multipliers $\mu_1=\mu_2=1$.
Its critical cone is zero, and the critical-cone theorem gives quadratic
growth locally, consistent with $x+y\geq\sqrt{x^2+y^2}$ in that corner.
These examples show why solving only interior stationary equations cannot
be a complete search for extrema on a compact region. One must include
every regular boundary stratum and examine singular feasible points directly.
\end{example}

\section{Convexity and global optimality}
\family{K04.05.09}

\begin{theorem}[Differential tests for convexity]\label{d04:05:convexity}
Let $U\subset\R^n$ be open and convex. A differentiable scalar function
$f$ is convex if and only if
\[
 f(y)\geq f(x)+Df(x)(y-x)\quad(x,y\in U).
\]
For $f\in C^2(U)$ this is equivalent to its Hessian being positive
semidefinite everywhere. A positive definite Hessian everywhere implies
strict convexity. A stationary point of a differentiable convex function
is a global minimum; of a strictly convex function it is the only minimum.
For a convex feasible subset $S\subset U$, a point $a\in S$ minimizes
a differentiable convex $f$ on $S$ exactly when
$Df(a)(x-a)\geq0$ for every $x\in S$.
\end{theorem}
\begin{proof}
Convexity gives
$f(x+t(y-x))\leq(1-t)f(x)+tf(y)$ for $0<t\leq1$.
Subtract $f(x)$, divide by $t$, and let $t\downarrow0$ to obtain the
supporting inequality. Conversely put $z=(1-t)x+ty$ and apply that
inequality from $z$ to each endpoint. Multiply by $1-t$ and $t$ and add;
the linear terms cancel, yielding convexity.

If the Hessian is nonnegative, on any segment the scalar restriction
$\phi(t)=f(x+t(y-x))$ has $\phi''\geq0$. The scalar mean-value theorem
makes $\phi'$ nondecreasing. Integrating, or applying scalar mean-value
again, gives $\phi(1)-\phi(0)\geq\phi'(0)$, which is the supporting
inequality. Conversely the convex midpoint inequality gives
$f(x+tv)+f(x-tv)-2f(x)\geq0$. Divide the order-two Taylor expansion
by $t^2$ and let $t\to0$ to get $D^2f(x)[v,v]\geq0$.
Positive definiteness makes $\phi''>0$ on nonconstant segments, so
$\phi'$ is strictly increasing and its secant averages give the strict
convexity inequality at $0<t<1$. The supporting inequality proves the
stationary-point assertion. Two distinct minimizers of a strictly convex
function would have a midpoint of smaller value, proving uniqueness.
Finally the stated derivative condition suffices by the supporting
inequality. If $a$ is a minimum on convex $S$, the feasible segment
$a+t(x-a)$ and its one-sided difference quotients prove necessity.
No closedness or existence claim for $S$ is implicit in this last test.
\end{proof}

\begin{remark}[Scope of optimization]\label{d04:05:optimization-scope}
The preceding tests cover smooth local extrema with regular equality or
independent active inequality constraints, compact existence, and the stated
convex global criterion. They do not assert a multiplier rule without a
constraint qualification. General convex duality, nonsmooth subgradients,
other constraint qualifications and numerical optimization belong to D20.
At singular constraints the feasible geometry, rather than an unjustified
multiplier equation, must supply the argument, as in
\cref{d04:05:worked-constraints}.
\end{remark}

\section{A finite-regularity Morse normal form}
\family{K04.05.11}

\begin{theorem}[Bounded Morse lemma]\label{d04:05:morse}
Let $r\geq1$ be a finite integer, let $f$ be $C^{r+2}$ on an open
neighborhood of $a\in\R^n$, and suppose $Df(a)=0$ and $D^2f(a)$ is
nondegenerate. If its inertia has $p$ positive and $q$ negative eigenvalues,
$p+q=n$, then there is a local $C^r$ coordinate diffeomorphism
$y=\chi(x)$ with $\chi(a)=0$ such that
\[
 f(\chi^{-1}(y))=f(a)+y_1^2+\cdots+y_p^2
                                  -y_{p+1}^2-\cdots-y_n^2.
\]
If $f$ is smooth, $\chi$ and its inverse are smooth. This is the bounded
finite-regularity version proved here; it makes no assertion that arbitrary
$C^2$ data admit coordinates of a particular positive differentiability
class, nor does it assert an optimal loss of derivatives.
\end{theorem}
\begin{proof}
Translate $a$ to zero and take a small ball on which every radial segment
stays in the domain. The order-two integral Taylor formula gives the exact
coefficient representation
\[
 f(x)-f(0)=\sum_{i,j}a_{ij}(x)x_ix_j,\qquad
 a_{ij}(x)=\int_0^1(1-t)\partial_i\partial_j f(tx)\,\dd t.
\]
The matrix $A(x)=(a_{ij}(x))$ is symmetric, $A(0)=H_f(0)/2$, and its
entries are $C^r$. Here is the required parameter-integral justification.
For a continuous integrand with continuous first $x$ derivative on a product
of a small closed ball and $[0,1]$, apply the segment FTC in $x$ to the
integrand. After subtracting its proposed integral derivative, the error is
bounded by $\norm{h}$ times the supremum of the change of that derivative
between $x$ and $x+sh$, uniformly in $s,t\in[0,1]$. Uniform continuity on
the compact product makes this supremum tend to zero. Thus the integral
has the proposed total derivative. Its derivative is continuous by the
same uniform estimate. Repeat for each derivative through order $r$.
For the present integrand each such derivative is continuous, by the chain
rule and $f\in C^{r+2}$; factors of $t$ cause no difficulty. This proves
the claimed $C^r$ property using only ordinary one-variable integrals,
without invoking a later parameter-integration theorem.

By the spectral theorem and a constant invertible linear change of input
coordinates, assume $A(0)$ is diagonal with nonzero entries $d_1^0,\ldots,d_n^0$.
We now complete squares in the coefficient matrix, keeping its argument $x$
fixed during the algebra. Its first pivot $d_1(x)=a_{11}(x)$ stays nonzero
near zero, and
\[
 \sum_{i,j}a_{ij}(x)x_ix_j
 =d_1(x)\left(x_1+\sum_{j>1}\frac{a_{1j}(x)}{d_1(x)}x_j\right)^2
  +\sum_{i,j>1}\left(a_{ij}(x)
                -\frac{a_{1i}(x)a_{1j}(x)}{d_1(x)}\right)x_ix_j.
\]
The residual matrix is symmetric $C^r$, and at zero it is diagonal with
entries $d_2^0,\ldots,d_n^0$. Repeating on the residual matrices gives,
after shrinking the neighborhood finitely many times,
\[
 f(x)-f(0)=\sum_{i=1}^n d_i(x)
             \left(x_i+\sum_{j>i}b_{ij}(x)x_j\right)^2,
\]
where all $d_i,b_{ij}$ are $C^r$, each $d_i(0)=d_i^0\ne0$, and
$b_{ij}(0)=0$. This recursive construction is finite: at step $i$ divide
the $i$th residual row by its nonzero diagonal pivot, subtract its outer
product divided by that pivot, and proceed to the next lower-right block.
Consequently it preserves $C^r$ without any differentiation or repeated
regularity loss. The signs $\eps_i=\sgn(d_i(x))=\sgn(d_i^0)$ are constant
near zero. Define
\[
 \chi_i(x)=\sqrt{\abs{d_i(x)}}
                  \left(x_i+\sum_{j>i}b_{ij}(x)x_j\right).
\]
The scalar square root is smooth on positive numbers by ordinary scalar
calculus. Hence $\chi$ is $C^r$, $\chi(0)=0$, and
$D\chi(0)=\diag(\sqrt{\abs{d_1^0}},\ldots,\sqrt{\abs{d_n^0}})$:
all terms differentiating the coefficients multiply zero coordinates, and
the off-diagonal coefficients at zero vanish. This derivative is invertible,
so \cref{d04:04:inverse-function} gives a local $C^r$ inverse. The exact
square identity becomes $f-f(0)=\sum_i\eps_i\chi_i^2$. Permuting these
coordinates puts the positive signs first. Inertia of $A(0)=H_f(0)/2$
identifies their numbers as $p,q$.

For smooth $f$, the integral coefficients are $C^r$ for every finite $r$,
and the same algebraic construction is smooth; the smooth inverse theorem
then supplies a smooth inverse. For $n=0$ the empty quadratic form is
nondegenerate on the zero space and the unique coordinate map suffices.
\end{proof}

\begin{corollary}[The index is independent of the chosen coordinates]\label{d04:05:morse-index}
In any local $C^1$ coordinate change giving a signed-square normal form at
a nondegenerate critical point of a $C^2$ function, the number of negative
squares equals the negative index of its Hessian.
\end{corollary}
\begin{proof}
If $x=\psi(y)$, $\psi(0)=a$, and $B=D\psi(0)$ is invertible, then
$\psi(y)-a=By+o(\norm{y})$. The second-order Peano formula for $f$,
using $Df(a)=0$, gives
$f(\psi(y))-f(a)=D^2f(a)[By,By]/2+o(\norm{y}^2)$.
If the left side is a quadratic form exactly, uniqueness of its degree-two
Taylor polynomial identifies it with $D^2f(a)[By,By]/2$.
Inertia is invariant under invertible linear congruence by
\cref{base:linear}. This argument needs only $C^1$ for the coordinate
change, not an unproved second derivative of it.
\end{proof}

\source{HUB 3.6, printed p. 347 (PDF p. 364), mentions Morse's lemma
and points to Milnor; the passage was read and visually checked. It does
not supply a finite-regularity proof there. The present coefficient-integral
and iterative square-completion proof establishes exactly the displayed
$C^{r+2}$-data/$C^r$-coordinates version. A stronger regularity assertion
from another source is not being imported or claimed proved.}

% END ORIGINAL INPUT: chapters/05.tex


% BEGIN ORIGINAL INPUT: chapters/06.tex
\chapter{Multiple Riemann integrals and iterated integration}
\label{ch:d04-06}
\source{SPI, Chapter 3, Theorems 3-1--3-10, and Addenda 2 (printed pp. 145--146); HUB, Theorem 4.5.10 and Appendix A.17.}

Throughout this chapter dimensions are positive integers unless explicitly stated otherwise. The integrals constructed here are Riemann integrals. A null set will mean a set admitting arbitrarily small countable box covers; this definition does not presuppose a measure integral.

\section{Boxes, Darboux sums, and tags}
\family{K04.06.01}\family{K04.06.02}
\begin{definition}\label{d04:06:boxes}
A box is $Q=\prod_{i=1}^n[a_i,b_i]$, $a_i\leq b_i$, with content $|Q|=\prod_i(b_i-a_i)$. A grid partition is obtained by finitely partitioning each coordinate interval; its cells are the resulting closed boxes and its mesh is the largest cell diameter. Different cells meet on faces. A tagged sum for a bounded $f:Q\to\R$ is
\[
 S(f,P,\xi)=\sum_{C\in P}f(\xi_C)|C|,\qquad \xi_C\in C.
\]
Write $m_C=\inf_C f$, $M_C=\sup_C f$ and
$L(f,P)=\sum_C m_C|C|$, $U(f,P)=\sum_C M_C|C|$.
The lower and upper integrals are $\underline{\int}_Q f=\sup_P L(f,P)$ and $\overline{\int}_Qf=\inf_P U(f,P)$. Equality defines integrability and the common value $\int_Q f$.
On a degenerate box all sums and integrals of bounded functions are zero. In dimension zero the sole box is the point $\R^0$, its content is the empty product $1$, and integration is evaluation; subsequent positive-dimensional null-set assertions are not assertions about this point.
\end{definition}

\begin{theorem}[Equivalent integral criteria]\label{d04:06:criteria}
For bounded $f$ on a nondegenerate box, the following are equivalent: Darboux integrability; for every $\eps>0$ some grid has $U-L<\eps$; and existence of $I$ such that all tagged sums of mesh sufficiently small differ from $I$ by less than $\eps$. In the last condition $I=\int_Q f$.
\end{theorem}
\begin{proof}
Refinement increases lower sums and decreases upper sums: split each old cell and use $m_C\leq m_{C'}\leq M_{C'}\leq M_C$ and additivity of products of interval lengths. Any two grids have a common refinement, so every lower sum is at most every upper sum. Completeness of $\R$ gives the lower and upper integrals, and a common refinement proves their equality equivalent to arbitrarily small gaps.

Fix a grid $P_0$ with gap less than $\eps$ and put $|f|\leq M$. For any other grid of mesh $\delta$, call a cell bad if it crosses an internal coordinate face of $P_0$. The union of bad cells is contained in the $\delta$ neighborhood of the finite union of those faces. Its total cell content tends to zero with $\delta$: a face $x_i=c$ inside $Q$ contributes at most $2\delta\prod_{j\ne i}(b_j-a_j)$, with harmless enlargement at the exterior faces. All other cells lie in a $P_0$ cell. Refine by $P_0$; the change of a tagged sum under this refinement has absolute value at most $2M$ times the total bad content. The refined sum lies between $L(f,P_0)$ and $U(f,P_0)$. These also contain the Darboux integral, proving the mesh assertion.

Conversely choose a grid of mesh below the asserted threshold. Since there are finitely many cells, tags can be chosen so that its tagged sum is arbitrarily close to its upper sum, and other tags so that it is arbitrarily close to its lower sum (suprema need not be attained). Thus $U-L\leq2\eps$. The gap criterion proves integrability, and the inequalities $L\leq S\leq U$ identify the limit.
\end{proof}

\family{K04.06.03}\family{K04.06.04}
\begin{proposition}[Integral calculus on a box]\label{d04:06:algebra}
Integrable functions form a real algebra; the integral is linear, preserves order, and satisfies
\[
 \left|\int_Q f\right|\leq\int_Q|f|\leq |Q|\sup_Q|f|.
\]
A bounded function is integrable on $Q$ if and only if its restrictions to the cells of any fixed grid are integrable; then its integral is the sum of their integrals. Uniform limits of integrable functions are integrable, with convergence of integrals.
For $f:Q\to\R^k$, coordinatewise integration is equivalent to convergence of vector tagged sums, and $\norm{\int_Q f}\leq\int_Q\norm f$.
\end{proposition}
\begin{proof}
For a cell let $\operatorname{osc}_C f=M_C-m_C$. The estimates
\[
 \operatorname{osc}_C(f+g)\leq\operatorname{osc}_C f+\operatorname{osc}_C g,
 \quad\operatorname{osc}_C(fg)\leq\norm f_\infty\operatorname{osc}_C g+\norm g_\infty\operatorname{osc}_C f
\]
and $\operatorname{osc}_C|f|\leq\operatorname{osc}_C f$ prove all integrability assertions by common refinements. Linearity follows by adding tagged sums and taking limits. Positivity follows from nonnegative tagged sums; applying it to $-|f|\leq f\leq|f|$ gives the displayed bounds.
If $f$ is integrable, refine any good grid by the fixed grid; the nonnegative Darboux gap on each fixed cell is at most the total gap. Conversely take good grids on the finitely many cells and refine their coordinate cuts to a grid of $Q$. Lower and upper sums add on the fixed cells, proving integrability and additivity. Faces are duplicated only as possible tags; they have zero cell content and the same refinement argument already handles their values. For uniform limits use
$U(f,P)-L(f,P)\leq U(g,P)-L(g,P)+2|Q|\norm{f-g}_\infty$ and then the integral bound. For vectors finite coordinate convergence is equivalent to norm convergence. The triangle inequality for vector sums gives the norm inequality after taking limits; continuity of the norm and the oscillation estimate make $\norm f$ integrable.
\end{proof}

\section{Null sets, oscillation, and Jordan regions}
\family{K04.06.05}\family{K04.06.06}
\begin{definition}\label{d04:06:null}
A subset of $\R^n$ is null if for every $\eps>0$ it is covered by countably many boxes with sum of contents less than $\eps$. It has content zero if finite covers suffice. Closed or open boxes give the same definitions: enlarge the $j$th box to an open box with increase less than $\eps2^{-j}$; a closed box contained in an open box can conversely be enlarged slightly within it. The oscillation of bounded $f:Q\to\R$ at $x\in Q$ is
\[
 \omega_f(x)=\inf_{r>0}\sup\{|f(y)-f(z)|:y,z\in Q\cap B(x,r)\}.
\]
\end{definition}

\begin{lemma}[Elementary cover facts]\label{d04:06:cover-facts}
Countable unions of null sets are null. Compact null sets have content zero. Finite unions of coordinate faces in a bounded box have content zero. A nondegenerate box is not null. A set of content zero has closure of content zero; bounded null sets need not have content zero.
\end{lemma}
\begin{proof}
Cover the $j$th null set with total content less than $\eps2^{-j}$ and enumerate all boxes. An open cover of a compact null set has a finite subcover, with the same upper bound for total content. A face is covered by a slab of arbitrarily small thickness; finitely many such slabs suffice.
For a finite cover of a box by boxes, insert all endpoints in one grid, clipping to the covered box. Each open grid cell is in at least one covering box; sum cell contents to obtain $|Q|\leq\sum|Q_j|$. A countable open cover of a compact box has a finite subcover, so a countable cover cannot have sum less than $|Q|$ either (first enlarge with arbitrarily small additional content). Finite closed box covers also cover the closure. Finally $\Q\cap[0,1]$ is countable and null, but any finite closed interval cover covers its closure $[0,1]$ and has total length at least $1$. This proves the last distinction.
\end{proof}

\begin{theorem}[Riemann's discontinuity criterion]\label{d04:06:riemann-criterion}
A bounded $f:Q\to\R$ is integrable if and only if its relative discontinuity set is null.
\end{theorem}
\begin{proof}
The set $D_\eta=\{x\in Q:\omega_f(x)\geq\eta\}$ is closed in $Q$: an oscillation bound on a ball persists on smaller balls about nearby points. It is compact, and the discontinuity set is $\bigcup_{k\geq1}D_{1/k}$.

Suppose $f$ is integrable. Choose a grid with gap less than $\eta\eps$. Away from its faces, any point of $D_\eta$ is in the interior of a cell with oscillation at least $\eta$. Such cells have total content at most $(U-L)/\eta<\eps$. Cover all grid faces by boxes of total content less than $\eps$. These two finite collections cover $D_\eta$, proving it has content zero. Notice that a cell merely touching $D_\eta$ at its boundary need not have oscillation at least $\eta$; the separate face cover is essential. This is the repair in Spivak's Addenda.

Conversely suppose the discontinuity set is null and $|f|\leq M$. Fix $\eta,\delta>0$. Cover the compact set $D_\eta$ by finitely many open boxes of total content less than $\delta$. At every point outside $D_\eta$ choose a relative neighborhood on which the oscillation of $f$ is less than $\eta$. These neighborhoods and the boxes cover $Q$. A finite subcover has a Lebesgue number: if none existed, a sequence of subsets with diameters tending to zero and lying in no member would have points converging in $Q$, contradicting openness of a member containing that limit. A sufficiently fine grid therefore has each cell either inside one covering box or in one small-oscillation neighborhood. Assign ambiguous cells in either way. By the finite grid cover inequality in \cref{d04:06:cover-facts}, the total content of cells assigned to the boxes is at most their summed contents, less than $\delta$. Hence
\[
 U(f,P)-L(f,P)\leq 2M\delta+\eta|Q|.
\]
Letting both parameters tend to zero proves integrability.
\end{proof}

\begin{corollary}[Continuity and changes on exceptional sets]\label{d04:06:exceptions}
Continuous functions on a box are integrable. If two bounded integrable functions agree off a null set, their integrals agree. If a bounded function vanishes off a set of content zero, it is integrable with integral zero. Replacing values on a null set alone does not guarantee preservation of integrability.
\end{corollary}
\begin{proof}
The first assertion follows from the criterion. For the second let $h=f-g$. The complement of a null set is dense in every nonempty open box, since such a box is not null. Thus every cell has $\inf h\leq0\leq\sup h$. Integrability forces $\int h=0$. In the third assertion the closure of the exceptional set has content zero; cover it by small open boxes and take a sufficiently fine grid whose cells meeting the closure lie in that cover. Upper sums of $|h|$ tend to zero. For the last assertion change the zero function on $\Q\cap[0,1]$ to $1$ there: its indicator has oscillation $1$ everywhere.
\end{proof}

\family{K04.06.07}\family{K04.06.08}
\begin{definition}\label{d04:06:jordan}
A bounded set $E\subset\R^n$ is Jordan measurable if its indicator, extended by zero to a box containing $\cl E$ in its interior, is Riemann integrable. Its volume is $\vol(E)=\int\mathbf1_E$. For bounded $f:E\to\R$, $\int_Ef$ means the integral of its zero extension, when this extension is integrable. The same definition applies to any function on an open set having bounded support, provided its zero extension is integrable; the open set itself need not be Jordan measurable.
\end{definition}

\begin{proposition}[Boundary and extension tests]\label{d04:06:region-test}
A bounded set $E$ is Jordan measurable exactly when $\partial E$ has content zero. This integral is independent of the enclosing box. Finite unions, intersections, and differences of Jordan sets are Jordan sets, and integrals are finitely additive on disjoint such sets whenever the zero extensions are integrable. If $E$ is Jordan and bounded $f:E\to\R$ is relatively continuous except on a null subset of $E$, its zero extension is integrable. In particular continuous functions on $\cl E$ are integrable over $E$.
\end{proposition}
\begin{proof}
The discontinuities of the indicator, in a box whose interior contains $\cl E$, are exactly $\partial E$: every neighborhood of a boundary point meets $E$ and its complement. This boundary is compact, so the criterion and the compact-null lemma prove the first equivalence. Add a larger enclosing box and subdivide along the old box faces; the integral over added cells is zero, proving independence. Products, sums, and differences of indicators prove the Boolean and additivity assertions using \cref{d04:06:algebra}. A discontinuity of the zero extension lies either in $\partial E$ or at an interior discontinuity of $f$. Their union is null, proving the extension assertion. This is a sufficient condition; integrability of an arbitrary function on an arbitrary subset is always checked on its actual zero extension.
\end{proof}

\section{Fubini with lower and upper slice integrals}
\family{K04.06.09}\family{K04.06.10}
\begin{theorem}[Darboux Fubini]\label{d04:06:fubini}
Let $A\subset\R^p$, $B\subset\R^q$ be compact boxes and let $f$ be Riemann integrable on $A\times B$. Define, for every $x\in A$,
\[
 \ell(x)=\underline{\int}_B f(x,y)\,\dd y,
 \qquad u(x)=\overline{\int}_B f(x,y)\,\dd y.
\]
Both functions are Riemann integrable on $A$ and
\[
 \int_A\ell=\int_{A\times B}f=\int_Au.
\]
The analogous assertions with the two coordinate blocks interchanged hold. In particular, if every inner slice is integrable, its ordinary integral defines an integrable outer function and gives the product integral.
\end{theorem}
\begin{proof}
For grids $P_A,P_B$ write $P=P_A\times P_B$. For $x\in C\in P_A$,
\[
 \sum_{D\in P_B}\inf_{C\times D}f\,|D|
 \leq \ell(x)\leq u(x)\leq
 \sum_{D\in P_B}\sup_{C\times D}f\,|D|.
\]
Take the infimum or supremum over $C$ and sum after multiplying by $|C|$. It follows that
\[
 L(f,P)\leq L(\ell,P_A)\leq U(\ell,P_A)\leq U(u,P_A)\leq U(f,P)
\]
and also
$L(f,P)\leq L(\ell,P_A)\leq L(u,P_A)\leq U(u,P_A)\leq U(f,P)$.
Every grid of the product box is a product of coordinate grids. Choose such grids with $U(f,P)-L(f,P)$ tending to zero. These inequalities prove that the Darboux gaps for both outer functions tend to zero and trap their integrals at $\int f$. Interchanging blocks proves the other order; equality $\ell=u$ pointwise gives the ordinary-slice conclusion.
\end{proof}

\begin{example}[An exceptional slice]\label{d04:06:bad-slice}
On $[0,1]^2$ put $f(x,y)=\mathbf1_{\{0\}}(x)\mathbf1_{\Q}(y)$. Its support is contained in a line of content zero, so $\int f=0$, but its $x=0$ slice is not Riemann integrable. Here $\ell=0$ everywhere whereas $u=\mathbf1_{\{0\}}$. Both outer integrals are zero, exactly as the theorem states. Repeated ordinary integration at every parameter value is therefore a stronger hypothesis than product integrability.
\end{example}

\begin{proposition}[Size of the exceptional slice set]\label{d04:06:exceptional-slices}
In \cref{d04:06:fubini}, the parameters with a nonintegrable slice form a null set.
\end{proposition}
\begin{proof}
The nonnegative integrable function $w=u-\ell$ has integral zero. For $\eta>0$ and a grid whose upper sum for $w$ is less than $\eps\eta$, all cells meeting $\{w\geq\eta\}$ have supremum at least $\eta$, so their total content is less than $\eps$. Thus this level set has content zero. Their countable union for $\eta=1/k$ is precisely $\{w>0\}$. This argument proves nullity without a measure-theoretic Fubini theorem. It does not license arbitrary definitions on that null set: such definitions may destroy outer Riemann integrability.
\end{proof}

\family{K04.06.11}\family{K04.06.12}
\begin{proposition}[Slicing and geometric integrals]\label{d04:06:cavalieri}
Suppose $A\subset\R^{n-1}$ is Jordan and $a,b:\cl A\to\R$ are continuous, $a\leq b$. The bounded region
$E=\{(x,t):x\in A,\ a(x)\leq t\leq b(x)\}$ is Jordan, and for continuous $f$ on $\cl E$,
\[
 \int_E f(x,t)\,\dd(x,t)=\int_A\left(\int_{a(x)}^{b(x)}f(x,t)\,\dd t\right)\dd x.
\]
Endpoint inclusion choices have no effect. In particular $\vol(E)=\int_A(b-a)$. If $\rho\geq0$ is integrable on bounded Jordan $E$, its mass, first moments, and center of mass are
\[
 m=\int_E\rho,\qquad m_i=\int_E x_i\rho(x)\,\dd x,\qquad c_i=m_i/m\quad(m>0).
\]
The moment of inertia about a fixed affine axis $L$ is $\int_E\dist(x,L)^2\rho(x)\,\dd x$.
\end{proposition}
\begin{proof}
The graph of a continuous scalar function on a compact box has content zero: subdivide its base finely enough that its oscillation on each cell is less than $\eta$, and cover each graph piece by the cell times an interval of length at most $2\eta$; the total content is at most $2\eta$ times the base content. For a compact base $\cl A$ the same argument uses uniform continuity only on that base; cells meeting it have total content bounded by an enclosing box, and each such graph piece has arbitrarily small vertical oscillation. Also $\partial A$ times a bounded interval has content zero by multiplying a small cover by that interval. The boundary of $E$ is contained in this latter set together with the two graphs, proving Jordan measurability. The zero extension of $f$ is continuous except on that boundary. Its slices are continuous on their intervals and zero elsewhere, with at most endpoint jumps. Apply the ordinary-slice corollary of \cref{d04:06:fubini}. The same argument for $f=1$ proves the volume formula. The moment integrands are integrable because coordinate and squared-distance functions are continuous and bounded on an enclosing box, and products preserve integrability. These quantities are applications of this integral, not new definitions of integration; for instance $\int_E(x_i-c_i)\rho=0$ follows from linearity.
\end{proof}

\section{Improper integrals and a measure comparison}
\family{K04.06.13}
\begin{definition}\label{d04:06:improper-definition}
On an open $\Omega\subset\R^n$, call $f$ locally Riemann integrable if its restriction to each compact box in $\Omega$ is integrable. Let $\mathcal K$ be the compact Jordan subsets of $\Omega$. The restriction of such an $f$ to each $K\in\mathcal K$ is integrable: cover $K$ by finitely many boxes in $\Omega$, obtain boundedness and a null discontinuity set on their interiors, and apply the zero-extension criterion. Define
\[
 \int_\Omega |f|=\sup_{K\in\mathcal K}\int_K|f|\in[0,\infty].
\]
When this number is finite we call the integral absolutely convergent. A regular exhaustion is a nested sequence $K_j\in\mathcal K$ with $K_j\subset\Int K_{j+1}$ and $\bigcup_j\Int K_j=\Omega$. Conditional convergence along a specified exhaustion means convergence of $\int_{K_j}f$ without absolute convergence; the exhaustion is part of that assertion.
\end{definition}

\begin{theorem}[Absolute convergence and exhaustions]\label{d04:06:absolute-exhaustion}
Regular exhaustions exist. For $f\geq0$, the integrals along every regular exhaustion increase to the displayed supremum, possibly infinity. For absolutely convergent $f$, $\int_{K_j}f$ has the same finite limit for every regular exhaustion, denoted $\int_\Omega f$, and
\[
 \left|\int_\Omega f-\int_K f\right|\leq
 \int_\Omega|f|-\int_K|f|\qquad(K\in\mathcal K).
\]
\end{theorem}
\begin{proof}
For each $j$ the compact set $\{x:\norm x\leq j,\ \dist(x,\R^n\setminus\Omega)\geq1/j\}$ can be covered by finitely many closed cubes contained in $\Omega$. Inductively cover it together with the previous chosen union by the interiors of finitely many such cubes. Their union is compact Jordan and produces a regular exhaustion; an empty complement uses distance $+\infty$. Every compact subset of $\Omega$ is contained in one exhaustion member by a finite subcover and nestedness. Thus positive exhaustion integrals have the required supremum.
Write $A=\sup_K\int_K|f|<\infty$. Given $\eps>0$ choose $K$ with $\int_K|f|>A-\eps$. For any $L\supset K$ in $\mathcal K$, finite additivity gives $\int_{L\setminus K}|f|<\eps$. For any two $L,H\supset K$, compare their integrals through $L\cup H$ to obtain $|\int_Lf-\int_Hf|\leq\int_{(L\cup H)\setminus K}|f|<\eps$. This proves a common Cauchy limit along all exhaustions. Comparing $K$ with large exhaustion members containing it and taking limits proves the bound.
\end{proof}

\begin{example}[Conditional order dependence]\label{d04:06:conditional-order}
On $(0,1)^2$ let $f(x,y)=(x^2-y^2)/(x^2+y^2)^2$. Since
$\partial_y[y/(x^2+y^2)]=f$ and $\partial_x[x/(x^2+y^2)]=-f$,
the iterated improper integrals are respectively $\pi/4$ and $-\pi/4$. Indeed the inner integrals are $1/(1+x^2)$ and $-1/(1+y^2)$. Square truncations $[\delta,1-\delta]^2$, $0<\delta<1/2$, instead have integral zero by interchange of the two variables and antisymmetry; taking a strictly decreasing sequence $\delta\downarrow0$ gives a regular exhaustion inside $(0,1)^2$. Absolute convergence fails: on $0<y<x/2$, $f\geq 12/(25x^2)$, so integration there gives a divergent multiple of $\int_0^1\dd x/x$. This is a direct calculation on proper truncated triangles, using the slicing proposition. Thus ordinary proper Fubini alone does not justify improper interchange. A sufficient uniform-tail interchange theorem is proved in \cref{d04:08:improper-interchange}.
\end{example}

\family{K04.06.14}
\begin{proposition}[Explicit Riemann--Lebesgue bridge]\label{d04:06:measure-bridge}
Using only the admitted basic Lebesgue measure and dominated/monotone integration foundation, a bounded Riemann-integrable function on a compact box is Lebesgue measurable and integrable, and the two integrals agree. In particular Jordan content agrees with Lebesgue measure. The same agreement holds for absolutely convergent locally Riemann-integrable functions on open sets. This comparison supplies no multivariable change-of-variables theorem.
\end{proposition}
\begin{proof}
Choose nested grids $P_j$ with Darboux gaps tending to zero, by refining successive good grids. Let $s_j,t_j$ equal the lower and upper cell values off grid faces, assigning on faces any values between $-M$ and $M$, where $|f|\leq M$. They are measurable simple functions; their measure integrals are the Darboux sums, because grid faces have measure zero by their slab covers. Let $Z$ be the countable union of these faces. Off $Z$, $s_j$ increases and $t_j$ decreases. Define their limits $s,t$ there and extend both by zero on $Z$. These are measurable, $s\leq f\leq t$ almost everywhere, and $s_j\to s$, $t_j\to t$ almost everywhere; no convergence of the arbitrarily assigned face values is claimed. Dominated convergence gives $\int(t-s)=0$, so $s=t$ outside a measurable null set. Basic Lebesgue measure is complete; null-cover sets have measure zero since box volumes agree with measure and countable subadditivity holds. Therefore $f$, equal almost everywhere to $s$, is measurable, and its integral is squeezed between the equal limits of the two Darboux sums. For an absolutely convergent function on an open set, apply the box comparison to its restrictions on each compact Jordan exhaustion member (zero extended), and monotone convergence to $|f|\mathbf1_{K_j}$. Dominated convergence then gives the asserted agreement for $f$ itself.
\end{proof}

% END ORIGINAL INPUT: chapters/06.tex


% BEGIN ORIGINAL INPUT: chapters/07.tex
\chapter{Change of variables and geometric Jacobians}
\label{ch:d04-07}
\source{SPI, Chapter 3, Theorem 3-13 and its localization and coordinate-factorization proof; HUB, Sections 4.9--4.10 and Appendix A.19, printed pp. 765--772. The latter passage uses Lipschitz derivatives; the proof below instead proves the stated $C^1$ version by scalar substitution.}

An ordinary volume integral carries a positive density. The determinant enters with its absolute value, regardless of the sign of a coordinate system. We prove the substitution theorem before using geometric coordinate formulas. In particular no assertion about volumes of nonlinear images is a premise of that proof.

\section{Elementary maps and localization}
\family{K04.07.02}\family{K04.07.05}
\begin{lemma}[Finite continuous localization]\label{d04:07:localization}
Let $K$ be a compact subset of an open set $V$ covered by open sets $V_\alpha$. If continuous $f:V\to\R$ has support contained in $K$, it is a finite sum $f=\sum_i f_i$ of continuous functions with compact supports in individual $V_{\alpha_i}$. Consequently a substitution identity for continuous compact-supported functions localizes over an open cover of the target of a diffeomorphism.
\end{lemma}
\begin{proof}
If $K$ is empty then $f=0$ and the empty sum suffices. Otherwise, for each point of $K$ choose a ball $B_i$ whose closed ball of twice the radius is contained in one $V_\alpha$; finitely many $B_i$ cover $K$. Choose continuous functions $0\leq\eta_i\leq1$ equal to $1$ on $\cl B_i$ and zero outside that doubled ball, for example a piecewise linear function of the distance to its center. Their sum $s$ is positive on a neighborhood of $K$. Put $f_i=f\eta_i/s$ where $s>0$ and $f_i=0$ where $s=0$. This is continuous: the closed set $K$ has a neighborhood on which $s$ is bounded away from zero, and $f=0$ off $K$. Its support is compact and lies in $V_{\alpha_i}$, and the sum is $f$. If $\phi:U\to V$ is a diffeomorphism and each $V_{\alpha_i}=\phi(U_i)$, apply the local identity to $f_i$ and add. The functions $f_i\circ\phi$ have compact support in $U_i$ because $\phi^{-1}$ is continuous. Finite additivity of the integral completes the consequence. This construction needs no differentiable partition of unity.
\end{proof}

\begin{lemma}[Changing one coordinate]\label{d04:07:one-coordinate}
Suppose $T:U\to V$ is a $C^1$ diffeomorphism between open subsets of $\R^n$ and has the form
\[
 T(x_1,\ldots,x_n)=(x_1,\ldots,x_{j-1},\psi(x),x_{j+1},\ldots,x_n).
\]
For every continuous compact-supported $f$ in $V$,
\[
 \int_Vf(y)\,\dd y=\int_U f(T(x))|\partial_j\psi(x)|\,\dd x.
\]
Coordinate permutations also satisfy substitution, with absolute determinant $1$.
\end{lemma}
\begin{proof}
It is enough by \cref{d04:07:localization} to consider support compact in $T(B)$ for a bounded open box $B$ with $\cl B\subset U$ and $\partial_j\psi$ of constant nonzero sign there. Such boxes cover $U$: the determinant is $\partial_j\psi\ne0$. Write $x=(x',t)$, with coordinate $j$ last in the notation only, and $B=B'\times(a,b)$. For fixed $x'$, $t\mapsto\psi(x',t)$ is strictly monotone by the scalar mean-value theorem. The slice of $T(B)$ is exactly the interval with endpoints $\psi(x',a),\psi(x',b)$. Its boundary consists of two continuous graphs and a bounded set over $\partial B'$, hence is of content zero by the graph proof in \cref{d04:06:cavalieri}.

Extend $f$ by zero beyond $T(B)$. It is continuous on all of $\R^n$, since its support is compact in $T(B)$. Likewise extend $f(T(x))|\partial_j\psi(x)|$ by zero beyond $B$; it is continuous. Scalar substitution on $[a,b]$, with reversed endpoints when the derivative is negative, gives for every $x'\in B'$
\[
 \int_{\min(\psi(x',a),\psi(x',b))}^{\max(\psi(x',a),\psi(x',b))} f(x',s)\,\dd s
 =\int_a^b f(x',\psi(x',t))|\partial_j\psi(x',t)|\,\dd t.
\]
All functions are supported in fixed boxes and are integrable there; every scalar slice is continuous with a zero extension or at most endpoint jumps. Integrate over $x'$ and use \cref{d04:06:fubini}, obtaining the desired identity. A coordinate permutation simply reorders the finite sums on product grids; alternatively repeated ordinary-slice Fubini for continuous functions exchanges any adjacent coordinates. Its matrix determinant has absolute value $1$.
\end{proof}

\begin{lemma}[Composition and invertible affine maps]\label{d04:07:affine}
If continuous compact-supported substitution holds for two diffeomorphisms, it holds for their composition. It holds for every invertible affine map $x\mapsto Ax+b$, with Jacobian $|\det A|$.
\end{lemma}
\begin{proof}
For $S:U\to V$ and $T:V\to W$, first substitute $T$ and then $S$ into the continuous compact-supported function $(f\circ T)|\det DT|$. The chain rule and determinant multiplication identify the resulting factor as $|\det D(T\circ S)|$. Row elimination expresses every invertible matrix as a product of elementary matrices: at each column choose a nonzero pivot, interchange rows, rescale the pivot row, and eliminate the other entries; induction reduces the matrix to the identity. Each row addition $x_j\mapsto x_j+cx_i$, row scaling $x_j\mapsto cx_j$ ($c\ne0$), and coordinate translation is a one-coordinate map. Row interchanges are permutations. Their Jacobian determinants are exactly the algebraic elementary determinants. Thus \cref{d04:07:one-coordinate} and composition prove the affine assertion without assuming linear volume scaling.
\end{proof}

\begin{theorem}[Continuous substitution for a $C^1$ diffeomorphism]\label{d04:07:continuous-substitution}
Let $U,V\subset\R^n$ be open and let $\phi:U\to V$ be a bijective $C^1$ map with $\det D\phi\ne0$. Then $\phi^{-1}$ is $C^1$ and, for every continuous compact-supported $f:V\to\R$,
\[
 \int_Vf(y)\,\dd y=\int_U f(\phi(x))|\det D\phi(x)|\,\dd x.
\]
Both integrals are integrals of continuous zero extensions with compact support. The same conclusion holds coordinatewise for finite-dimensional vector-valued $f$.
\end{theorem}
\begin{proof}
The inverse function theorem \cref{d04:04:inverse-function} gives a $C^1$ local inverse at each point. Bijectivity makes these local inverses the restrictions of the global inverse, proving its asserted regularity.

We supply the local factorization that proves the identity. Fix $a\in U$, translate source and target, and postcompose with $D\phi(a)^{-1}$. These affine changes have already been handled. The resulting map $F$ satisfies $F(0)=0$ and $DF(0)=I$. For $k=0,\ldots,n$ set
\[
 H_k(x)=(F_1(x),\ldots,F_k(x),x_{k+1},\ldots,x_n).
\]
Here $H_0=\id$, $H_n=F$, and $DH_k(0)=I$ for every $k$. Apply the inverse function theorem to all finitely many $H_k$ and intersect their inverse neighborhoods, obtaining an open neighborhood $O$ of $0$ on which all restrictions are diffeomorphisms onto their open images. Define
\[
 E_k=H_k\circ H_{k-1}^{-1}:H_{k-1}(O)\longrightarrow H_k(O).
\]
For $i<k$ its $i$th coordinate is the $i$th input coordinate because both $H_k$ and $H_{k-1}$ have coordinate $F_i$; for $i>k$ both have coordinate $x_i$. Thus $E_k$ changes only coordinate $k$, and is a $C^1$ diffeomorphism. The one-coordinate lemma applies, and composition gives the identity for
$F|_O=E_n\circ\cdots\circ E_1$. Undoing the affine changes gives the identity on some neighborhood of $a$. The images of these neighborhoods cover $V$. Localize a given compact-supported continuous $f$ with \cref{d04:07:localization}, and sum the local identities. This proves the theorem in every positive dimension using only the scalar substitution theorem, the proved Riemann Fubini theorem, and the finite-dimensional inverse function theorem. No estimate of nonlinear image volume has been suppressed in the argument.
\end{proof}

\section{Null preservation and the full Riemann extension}
\family{K04.07.03}
\begin{lemma}[$C^1$ maps preserve null sets in equal dimensions]\label{d04:07:null-preservation}
If $F:U\to\R^n$ is $C^1$ on open $U$, the image of every null subset of $U$ is null. The image of every compact content-zero subset of $U$ has content zero. If $\phi:U\to V$ is a diffeomorphism and $\cl E\subset U$ is compact with $E$ Jordan, then $\phi(E)$ is Jordan and $\partial\phi(E)=\phi(\partial E)$.
\end{lemma}
\begin{proof}
First null covers may be made of cubes. For a box with positive side lengths, cover it by cubes of side $h\leq\min_i(b_i-a_i)$ from a grid; their total content is at most $2^n$ times the box content. Choose $h$ as small as desired. Enlarge degenerate boxes slightly first, with arbitrarily small total added content. Thus a null set has countable cube covers with arbitrarily small summed cube contents and arbitrarily small individual diameters.

On a ball whose closure lies in $U$, the segment derivative estimate bounds $\norm{F(x)-F(y)}\leq L\norm{x-y}$, because $DF$ is bounded on that closure. If cubes of side $s_i$ cover a null subset of this ball, choose in each nonempty intersection a point $z_i$. Its image lies in a cube centered at $F(z_i)$ of side $2L\sqrt n\,s_i$, whose content is $(2L\sqrt n)^n s_i^n$. The sum is arbitrarily small. A countable collection of such balls covers $U$ (use rational centers and radii with closure in $U$), proving the first assertion. A compact content-zero set has compact null image, hence content-zero image by \cref{d04:06:cover-facts}. Finally a homeomorphism preserves relative interiors and closures. Since $\cl E$ and its image are compact inside the respective open sets, their relative boundaries equal their ambient boundaries. This proves the boundary equality, and the boundary criterion proves the Jordan assertion.
\end{proof}

\begin{lemma}[Continuous upper and lower approximations]\label{d04:07:sandwich}
Let $f:\R^n\to\R$ be bounded, Riemann integrable with compact support $K$ in open $V$. For every $\eps>0$ there are continuous compact-supported $g_-,g_+$ in $V$ such that $g_-\leq f\leq g_+$ everywhere and $\int(g_+-g_-)<\eps$.
\end{lemma}
\begin{proof}
If $f=0$ take both approximations zero. Otherwise $K$ is nonempty. Choose a box $Q$ containing $K$ in its interior and a continuous cutoff $0\leq\eta\leq1$, equal to $1$ on $K$ and with compact support in $V\cap\Int Q$. Such a function follows from the ball construction in \cref{d04:07:localization}: if $s$ is a finite sum of cutoffs positive on $K$ and $c=\min_Ks>0$, use $\eta=\min(1,s/c)$.
Let $|f|\leq M$; $M=0$ is immediate. Choose a grid $P$ with gap less than $\eps/2$. Let $Z$ be its finite union of faces including $\partial Q$, and put $\lambda(x)=\min(1,\dist(x,Z)/\delta)$. In the interior of a cell $C$ define
\[
 h_+(x)=\lambda(x)M_C+(1-\lambda(x))M,
 \qquad h_-(x)=\lambda(x)m_C-(1-\lambda(x))M.
\]
On all faces set $h_+=M$, $h_-=-M$. These functions are continuous on $Q$ because $\lambda$ tends to zero on every face. They bound $f$. The strip $Z_\delta=\{\dist(x,Z)<\delta\}\cap Q$ has content tending to zero, by the slab estimate. Off it the difference is exactly $M_C-m_C$, and everywhere the additional error above this cell difference is at most $2M$. Consequently
\[
 \int_Q(h_+-h_-)\leq U(f,P)-L(f,P)+2M\vol(Z_\delta)<\eps
\]
for sufficiently small $\delta$. The strips are finite unions of rectangular slabs, so this integral estimate uses only finite additivity and comparison. Put $g_\pm=\eta h_\pm$ on $Q$ and zero off $Q$. They are continuous and compact-supported in $V$. On $K$ multiplication changes nothing; off $K$, $f=0$ and $h_-\leq0\leq h_+$, so the bounds persist. Also $0\leq g_+-g_-\leq h_+-h_-$ on $Q$, proving the estimate.
\end{proof}

\begin{theorem}[Full compact-supported Riemann substitution]\label{d04:07:substitution}
For a $C^1$ diffeomorphism $\phi:U\to V$ and bounded $f:V\to\R$ with compact support in $V$, its zero extension is Riemann integrable if and only if the zero extension of
\[
 g(x)=f(\phi(x))J(x),\qquad J(x)=|\det D\phi(x)|,
\]
is Riemann integrable. When either holds,
\[
 \int_Vf=\int_U(f\circ\phi)J.
\]
\end{theorem}
\begin{proof}
Suppose first that $f$ is integrable. Its discontinuity set is null. Discontinuities of $f\circ\phi$ are contained in its inverse image under $\phi$, which is null by applying \cref{d04:07:null-preservation} to $\phi^{-1}$. Multiplication by continuous $J$ adds no discontinuities. The resulting support is contained in the compact set $\phi^{-1}(\supp f)$, inside $U$, and $J$ is bounded there. The zero extension therefore has only the indicated null discontinuity set and is bounded; it is integrable by \cref{d04:06:riemann-criterion}.

Use \cref{d04:07:sandwich} to choose $g_-\leq f\leq g_+$ with integral gap less than $\eps$. Positivity of $J$, the continuous substitution theorem, and comparison give
\[
 \int_Vg_- =\int_U(g_-\circ\phi)J
 \leq \int_U(f\circ\phi)J
 \leq\int_U(g_+\circ\phi)J=\int_Vg_+.
\]
The integral of $f$ lies in the same interval of length less than $\eps$. Let $\eps\downarrow0$ to obtain equality. Conversely apply what was just proved to the integrable function $g$ and the diffeomorphism $\phi^{-1}$. The derivative inverse formula gives $|\det D\phi^{-1}(\phi(x))|=1/J(x)$, so its transformed function is $f$, proving both integrability and the reverse equality. This establishes the Riemann extension explicitly, rather than inferring it from a formula for continuous functions.
\end{proof}

\begin{corollary}[Bounded regions]\label{d04:07:regions}
Let $E\subset U$ be bounded Jordan with compact closure in $U$, and let $\phi:U\to V$ be a $C^1$ diffeomorphism. Then $\phi(E)$ is Jordan. For bounded $f:\phi(E)\to\R$, existence of either of the following Riemann integrals implies existence of the other, and
\[
 \int_{\phi(E)}f(y)\,\dd y=\int_Ef(\phi(x))|\det D\phi(x)|\,\dd x.
\]
In particular $\vol(\phi(E))=\int_E|\det D\phi|$.
\end{corollary}
\begin{proof}
Jordan measurability was proved in \cref{d04:07:null-preservation}. Extend $f$ by zero to $V$; its support is contained in compact $\phi(\cl E)\subset V$. Its pullback times $J$ is exactly the zero extension of the indicated integrand on $E$. Apply \cref{d04:07:substitution}. For $f=1$, both zero extensions are integrable by the Jordan boundary criterion and multiplication by a continuous function.
\end{proof}

\family{K04.07.01}\family{K04.07.08}
\begin{corollary}[Linear volume scaling]\label{d04:07:linear-volume}
For bounded Jordan $E\subset\R^n$ and every linear $A:\R^n\to\R^n$, $A(E)$ is Jordan and
\[
 \vol(A(E))=|\det A|\vol(E).
\]
For nonsingular $A$ this also gives substitution in any nonorthogonal linear coordinate system.
\end{corollary}
\begin{proof}
If $A$ is invertible, apply \cref{d04:07:regions} with constant derivative $A$. If $A$ is singular, its image lies in a proper linear hyperplane. In a coordinate whose normal component is nonzero this hyperplane is the graph of a linear function of the remaining coordinates. Its bounded parts have content zero by the graph proof in \cref{d04:06:cavalieri}. The bounded closure of $A(E)$ is contained in such a part, so $A(E)$ has content zero and is Jordan with volume zero. Its determinant is zero as well. This proves the singular case separately; no inverse substitution is asserted there.
\end{proof}

\family{K04.07.04}
\begin{proposition}[Signs and oriented top forms]\label{d04:07:orientation}
On every connected component of a diffeomorphism domain the continuous nonzero function $\det D\phi$ has a constant sign. The pullback of a normalized top form is
\[
 \phi^*(f(y)\,\dd y_1\wedge\cdots\wedge\dd y_n)
 =f(\phi(x))\det D\phi(x)\,\dd x_1\wedge\cdots\wedge\dd x_n.
\]
Thus an orientation-preserving coordinate change gives the ordinary integral formula, whereas reversal changes the oriented integral's sign.
\end{proposition}
\begin{proof}
If both signs occurred on a connected component, the inverse images of the positive and negative half-lines would disconnect it. Evaluation of the pullback on the ordered coordinate basis is the alternating multilinear determinant formula for the columns of $D\phi$, proving the displayed algebraic identity. Integrate its coefficient and compare $\det D\phi=\sgn(\det D\phi)|\det D\phi|$ with \cref{d04:07:substitution}. General form and manifold conventions are developed in \cref{ch:d04-13,ch:d04-15}; no oriented-volume convention is needed to prove the positive-density theorem above.
\end{proof}

\section{Polar, cylindrical, and spherical coordinates}
\family{K04.07.06}\family{K04.07.07}
\begin{proposition}[Geometric coordinate formulas]\label{d04:07:coordinates}
The following maps are diffeomorphisms on the indicated open coordinate ranges onto their images:
\begin{align*}
 P(r,\theta)&=(r\cos\theta,r\sin\theta),& J_P&=r,\\
 C(r,\theta,z)&=(r\cos\theta,r\sin\theta,z),&J_C&=r,\\
 S(r,\varphi,\theta)&=(r\sin\varphi\cos\theta,r\sin\varphi\sin\theta,r\cos\varphi),&J_S&=r^2\sin\varphi.
\end{align*}
For $P$ and $C$ take $r>0$, $\alpha<\theta<\alpha+2\pi$, and for $C$
allow every $z\in\R$. For $S$ take $r>0$, $0<\varphi<\pi$ and the same
open angular range for $\theta$.
On compact Jordan subregions of these ranges, substitution multiplies the integrand by the displayed positive factor. The determinants in the displayed coordinate orders are positive. Bounded Cartesian integrals may include their seams and axes by a limiting argument described below.
\end{proposition}
\begin{proof}
For $P$, the derivative columns are $(\cos\theta,\sin\theta)$ and $(-r\sin\theta,r\cos\theta)$, of determinant $r$. The radius determines $r$, and the chosen angular interval determines $\theta$ uniquely off its excluded ray; the inverse function theorem gives the local $C^1$ inverse and injectivity makes it global. The cylindrical derivative adds the unchanged $z$ direction, giving determinant $r$ and the same inverse argument. For $S$, put
\begin{align*}
 e_r&=(\sin\varphi\cos\theta,\sin\varphi\sin\theta,\cos\varphi),\\
 e_\varphi&=(\cos\varphi\cos\theta,\cos\varphi\sin\theta,-\sin\varphi),\\
 e_\theta&=(-\sin\theta,\cos\theta,0).
\end{align*}
These form an orthonormal positively oriented basis: direct cross multiplication gives $e_r\times e_\varphi=e_\theta$. The derivative columns are $e_r,re_\varphi,r\sin\varphi e_\theta$, giving determinant $r^2\sin\varphi$. The radius and $z/r=\cos\varphi$ determine $r,\varphi$, then the angular range determines $\theta$. Nonzero determinant supplies its $C^1$ inverse.

For clarity consider the full disk of radius $R$ and a bounded Riemann integrand $f$. Its transformed integrand on the full closed coordinate rectangle is Riemann integrable: away from $r=0$ and the angular faces, the discontinuity set is null by local inverse null preservation; those coordinate faces have content zero, so \cref{d04:06:riemann-criterion} applies. Remove $r<\delta$ and angular intervals of length $\delta$ at the two ends, apply the proper region theorem, and pass to the limit. On the parameter side the absolute contribution removed is at most
$\norm f_\infty(\pi\delta^2+R^2\delta)$, by integrating $r$. On the Cartesian side the origin is in a square of area $4\delta^2$, and the seam neighborhoods lie in a strip of width tending to zero inside a fixed bounding square; their upper integral contributions tend to zero. Thus the limiting identity is valid, with the isolated origin and ray assigned arbitrary bounded values whenever the integrability hypotheses hold. A bounded cylindrical region has the same estimates multiplied by a bounded $z$ length. For a bounded ball the removed $r<\delta$ part has parameter integral $4\pi\delta^3/3$; polar caps have parameter integral at most $(4\pi R^3/3)(1-\cos\delta)$; angular seam pieces have parameter integral at most $4R^3\delta/3$. Cartesian neighborhoods of the axis and seam lie in boxes or slabs of volume tending to zero. These estimates prove the analogous spherical limiting identity. They also apply to any subregion by comparison, provided the original zero extension is integrable and the truncated regions are Jordan. A bounded integrand on a compact coordinate rectangle may have seam values chosen in any way on those content-zero faces; integrability and the integral are unchanged by \cref{d04:06:exceptions}.
\end{proof}

\begin{example}[Volumes and nonorthogonal systems]\label{d04:07:volumes}
Putting $f=1$ in the disk and ball formulas gives $\pi R^2$ and $4\pi R^3/3$: integrate $r$ over $0<r<R$, $0<\theta<2\pi$, and $r^2\sin\varphi$ over $0<r<R$, $0<\varphi<\pi$, $0<\theta<2\pi$, respectively. Hence an ellipsoid $\{Ax:\norm x<R\}$ with nonsingular $A$ has volume $|\det A|4\pi R^3/3$. More generally coordinates $x=\sum_i u_i v_i$ in an arbitrary basis contribute the factor $|\det(v_1,\ldots,v_n)|$. These are all the same substitution theorem; orthogonality is unnecessary.
\end{example}

\family{K04.07.09}
\begin{proposition}[Finite sheet counting]\label{d04:07:multiplicity}
Suppose a parametrization has a finite decomposition into disjoint Jordan pieces $E_i$, up to a content-zero set, and on a neighborhood of each $E_i$ its restriction agrees with a diffeomorphism $\phi_i$. For bounded $f$ with all relevant zero extensions integrable,
\[
 \sum_i\int_{E_i}f(\phi_i(x))|\det D\phi_i(x)|\,\dd x
 =\int_{\R^n} f(y)\sum_i\mathbf1_{\phi_i(E_i)}(y)\,\dd y.
\]
Thus overlap in the images is counted, not erased.
\end{proposition}
\begin{proof}
Apply \cref{d04:07:regions} on each piece and add the finitely many integrable zero extensions. The sum of image indicators is exactly the integer number of pieces covering a point. Parameter-boundary sets of content zero contribute zero when integrands are bounded. This is a finite-sheet assertion, not a general area or coarea theorem.
\end{proof}

\begin{example}[A double angular traversal]\label{d04:07:double-angle}
The polar parametrization with $1<r<2$ and $0<\theta<4\pi$ traverses the annulus twice. Split it at angular multiples of $\pi$ into finitely many injective charts, and discard the radial boundaries (content zero). Its integral of $r\,\dd r\,\dd\theta$ is $6\pi$, while the annulus has area $3\pi$. The zero determinant at $r=0$ and the repeated seam are different issues: removing the origin does not repair the twofold covering.
\end{example}

\family{K04.07.10}
\begin{theorem}[Absolutely convergent improper substitution]\label{d04:07:improper}
Let $\phi:U\to V$ be a $C^1$ diffeomorphism and $f$ locally Riemann integrable on $V$. Then $f$ is absolutely integrable in the sense of \cref{d04:06:improper-definition} if and only if $g=(f\circ\phi)|\det D\phi|$ is, and in that case $\int_Vf=\int_Ug$. More generally a specified exhaustion and its image give equal conditional limits when those limits exist; an unrelated exhaustion requires a separate justification.
\end{theorem}
\begin{proof}
Local integrability of $g$ follows from the region theorem on compact boxes in $U$, because their images are compact Jordan in $V$. Choose a regular exhaustion $K_j$ of $U$. Its images $\phi(K_j)$ are compact Jordan, nested with each contained in the interior of the next, and exhaust $V$, because a homeomorphism preserves these properties. Proper substitution gives equality on each pair, also with absolute values. The positive exhaustion theorem proves equivalence of absolute convergence and equality of absolute integrals. The signed exhaustion theorem then proves the claimed equality. Without absolute convergence equality on the corresponding truncations still gives equality of their limits, but proves no exhaustion independence. In dimension zero substitution is evaluation and the empty determinant is $1$; all nonzero-dimensional geometric arguments above are then unnecessary.
\end{proof}

% END ORIGINAL INPUT: chapters/07.tex


% BEGIN ORIGINAL INPUT: chapters/08.tex
\chapter{Integrals depending on parameters}
\label{ch:d04-08}
\source{Z2, Sections 17.1.2--17.1.4 (printed pp. 406--410), the uniform-tail tests and limit theorems in 17.2.1--17.2.4 (checked pp. 417--425), and 17.5.1--17.5.2 (pp. 471--473). The moving-domain statement below is obtained from the proved substitution theorem and a determinant calculation.}

The useful hypothesis is control on an entire parameter neighborhood. Separate partial-derivative formulas, without that control, neither prove a total derivative nor justify passage to an improper limit. Our parameter space is always finite dimensional.

\section{A fixed bounded region}
\family{K04.08.01}
\begin{theorem}[Continuity and uniform limits]\label{d04:08:continuity}
Let $E\subset\R^n$ be bounded Jordan, $P\subset\R^m$ a parameter set, and $f:E\times P\to\R^k$. Suppose each $f(\cdot,p)$ is Riemann integrable on $E$. If
\[
 \sup_{x\in E}\norm{f(x,p)-f(x,p_0)}\longrightarrow0\quad(p\to p_0\text{ in }P),
\]
then $F(p)=\int_E f(x,p)\,\dd x$ is continuous at $p_0$. In particular joint continuity on $\cl E\times P$ implies continuity of $F$ when $P$ is open, and also on any compact parameter set on which that joint continuity is given. Uniform convergence of integrands on $E\times P$ implies uniform convergence of their parameter integrals.
\end{theorem}
\begin{proof}
The vector integral estimate in \cref{d04:06:algebra} gives
\[
 \norm{F(p)-F(p_0)}\leq\vol(E)\sup_{x\in E}\norm{f(x,p)-f(x,p_0)}.
\]
This proves the first and the uniform-convergence assertions. For joint continuity, restrict parameters to a compact neighborhood of $p_0$ in the open case; uniform continuity on its product with $\cl E$ gives exactly the required supremum convergence. On a compact parameter set use the product itself and relative convergence. The integrability in these continuous cases follows from \cref{d04:06:region-test}. A region of volume zero gives the zero function, consistent with the estimate.
\end{proof}

\family{K04.08.02}\family{K04.08.03}
\begin{theorem}[A total derivative under a proper integral]\label{d04:08:proper-derivative}
Let $P\subset\R^m$ be open and $E$ bounded Jordan. Suppose $f:E\times P\to\R^k$ is integrable in $x$ for each $p$, is differentiable in $p$ for each $x$, and $x\mapsto D_pf(x,p)$ is integrable as a function with values in $\Hom(\R^m,\R^k)$. At $p_0$ suppose
\[
 \sup_{x\in E}\norm{D_pf(x,p)-D_pf(x,p_0)}\longrightarrow0\qquad(p\to p_0).
\]
Then $F(p)=\int_E f(x,p)\,\dd x$ is totally differentiable at $p_0$, and
\[
 DF(p_0)h=\int_E D_pf(x,p_0)h\,\dd x\qquad(h\in\R^m).
\]
If $f$ and $D_pf$ are continuous on $\cl E\times P$, then $F\in C^1(P,\R^k)$ and this formula holds everywhere. No differentiability in $x$ is needed.
\end{theorem}
\begin{proof}
For small $r>0$ with $\cl B(p_0,r)\subset P$ put
\[
 \omega(r)=\sup_{\substack{x\in E\\\norm{p-p_0}\leq r}}
 \norm{D_pf(x,p)-D_pf(x,p_0)}.
\]
The hypothesis says $\omega(r)\to0$. For fixed $x$ subtract the linear function $p\mapsto D_pf(x,p_0)(p-p_0)$ from $f(x,p)$. Along the segment from $p_0$ to $p_0+h$, its derivative has norm at most $\omega(\norm h)\norm h$. The vector segment mean-value inequality \cref{d04:02:mean-value} therefore gives
\[
 \norm{f(x,p_0+h)-f(x,p_0)-D_pf(x,p_0)h}\leq\omega(\norm h)\norm h.
\]
Integrate and use the norm bound to obtain a remainder at most $\vol(E)\omega(\norm h)\norm h=o(\norm h)$. The proposed derivative is a linear map since finite-dimensional integrals are linear. This proves a total derivative, uniformly in all small directions, rather than just separate directional equalities. In the continuous case uniform continuity on compact parameter neighborhoods gives the hypotheses at every point; applying \cref{d04:08:continuity} to the operator-valued integrand $D_pf$ gives continuity of $DF$.
\end{proof}

\family{K04.08.04}
\begin{corollary}[Finite repeated differentiation]\label{d04:08:repeated}
If the parameter derivatives $D_p^j f$, $0\leq j\leq r$, exist and are continuous on $\cl E\times P$, then $F\in C^r(P,\R^k)$ and
\[
 D^jF(p)[h_1,\ldots,h_j]
 =\int_E D_p^jf(x,p)[h_1,\ldots,h_j]\,\dd x\quad(0\leq j\leq r).
\]
Equivalently $\partial_p^\alpha F=\int_E\partial_p^\alpha f$ for $|\alpha|\leq r$.
\end{corollary}
\begin{proof}
The spaces of $j$-linear maps between the fixed finite-dimensional spaces are finite dimensional. Apply \cref{d04:08:proper-derivative} first to $f$, then to the integrand $D_pf$ with this finite-dimensional target, and continue $r$ times. At each step its next parameter derivative is jointly continuous, so the uniform estimate and integrability requirements have already been verified. Evaluation on $h_1,\ldots,h_j$ is a linear map on that target and commutes with finite tagged-sum limits; this yields the displayed formula. Evaluation on coordinate vectors gives the multi-index version, with \cref{d04:03:symmetry,d04:03:factorial-bridge} identifying the coordinate and multilinear derivatives.
\end{proof}

\begin{example}[Exponential moments]\label{d04:08:exponential-moments}
For a nonnegative integrable density $\rho$ on bounded Jordan $E$, the function $Z(p)=\int_E e^{p\cdot x}\rho(x)\,\dd x$ is $C^r$ for every finite $r$, and
\[
 DZ(p)h=\int_E(h\cdot x)e^{p\cdot x}\rho(x)\,\dd x,
 \quad D^2Z(p)[h,k]=\int_E(h\cdot x)(k\cdot x)e^{p\cdot x}\rho(x)\,\dd x.
\]
Here $\rho$ need not be continuous. Indeed each derivative integrand is integrable as a product of $\rho$ with a continuous function; since $\rho$ and $E$ are bounded, all parameter remainders and derivative differences are uniformly controlled on compact parameter balls. The proof of \cref{d04:08:proper-derivative}, iterated with these bounds, proves the assertions. If $\int_E\rho>0$, then $Z(p)>0$ since $e^{p\cdot x}$ has a positive minimum on $\cl E$; $D\log Z$ is the center of mass of the reweighted density. Its Hessian is nonnegative because
$D^2\log Z(p)[h,h]=\int_E(h\cdot x-c_h)^2e^{p\cdot x}\rho/Z\geq0$, where $c_h=\int_E(h\cdot x)e^{p\cdot x}\rho/Z$. This follows by the quotient rule and expansion of the square.
\end{example}

\section{Moving endpoints and bounded moving domains}
\family{K04.08.05}
\begin{theorem}[Leibniz rule with vector parameters]\label{d04:08:endpoints}
Let $I\subset\R$ be an open interval and $P\subset\R^m$ be open, $f:I\times P\to\R^k$ and $D_pf$ continuous, and let $a,b:P\to I$ be $C^1$. Then, with oriented scalar integration when $b<a$,
\[
 F(p)=\int_{a(p)}^{b(p)}f(t,p)\,\dd t
\]
is $C^1$ and
\[
 DF(p)h=\int_{a(p)}^{b(p)}D_pf(t,p)h\,\dd t
 +f(b(p),p)Db(p)h-f(a(p),p)Da(p)h.
\]
This formula also holds at parameters where $a(p)=b(p)$.
\end{theorem}
\begin{proof}
Locally choose a compact interval $J\subset I$ containing the values of both endpoints, and fix $c\in J$. Define $H(s,p)=\int_c^s f(t,p)\,\dd t$. Scalar FTC gives $D_sH=f(s,p)$, and the proper parameter theorem on the fixed interval from $c$ to $s$ gives $D_pH(s,p)=\int_c^sD_pf(t,p)\,\dd t$. These partial derivatives are jointly continuous: the moving piece is bounded by the interval length times a uniform bound, and the fixed piece is controlled by uniform continuity on $J$ times a compact parameter neighborhood. Thus \cref{d04:02:continuous-partials} makes $H$ a $C^1$ map. Now $F(p)=H(b(p),p)-H(a(p),p)$; \cref{d04:02:chain-rule} gives the displayed formula. Oriented integration and subtraction are valid for every order of endpoints, including equality, so no division by $b-a$ or noncrossing assumption has entered the proof.
\end{proof}

\family{K04.08.06}
\begin{lemma}[Derivative of a positive Jacobian]\label{d04:08:jacobian-derivative}
For a differentiable family of invertible matrices $A(p)$, $J(p)=|\det A(p)|$ satisfies
\[
 DJ(p)h=J(p)\tr(A(p)^{-1}DA(p)h).
\]
\end{lemma}
\begin{proof}
Multilinearity of the determinant in columns shows
$\det(I+tB)=1+t\tr B+O(t^2)$: the linear terms choose one column from $tB$ and the rest from $I$, giving $tB_{ii}$, and every remaining term chooses at least two columns from $tB$. This is a finite polynomial estimate, locally uniform for bounded $B$. Factor
$A+H=A(I+A^{-1}H)$ to obtain
$D(\det)(A)H=\det A\tr(A^{-1}H)$. Since the determinant is nonzero, it has constant sign near $A$; differentiating its absolute value multiplies this formula by that sign, proving the claim.
\end{proof}

\begin{theorem}[A bounded moving-domain formula]\label{d04:08:moving-domain}
Let $E\subset\R^n$ be bounded Jordan with compact closure in open $W$, and let $P\subset\R^m$ be open. Suppose $\Phi:P\times W\to\R^n$ is $C^1$, every $\Phi_p=\Phi(p,\cdot)$ is a diffeomorphism onto its open image, and $D_x\Phi$ is differentiable in $p$ with continuous $D_pD_x\Phi$. For example, a jointly $C^2$ family meets these finite-regularity requirements. Let $f(p,y)$ be $C^1$ on an open neighborhood of $\{(p,\Phi(p,x)):p\in P,x\in\cl E\}$. Set $E_p=\Phi_p(E)$ and
\[
 F(p)=\int_{E_p}f(p,y)\,\dd y,\quad A(p,x)=D_x\Phi(p,x),\quad J(p,x)=|\det A(p,x)|.
\]
Then $F\in C^1(P)$ and
\begin{align*}
 DF(p)h=\int_E J(p,x)\big(&D_pf(p,\Phi(p,x))h
 +D_yf(p,\Phi(p,x))D_p\Phi(p,x)h\\
 &+f(p,\Phi(p,x))\tr(A(p,x)^{-1}D_pA(p,x)h)\big)\,\dd x.
\end{align*}
If, in addition, $D_xD_p\Phi=D_pD_x\Phi$ (in particular for jointly $C^2$ families), define
$v_h(y)=D_p\Phi(p,\Phi_p^{-1}(y))h$. The same identity becomes
\[
 DF(p)h=\int_{E_p}\big(D_pf(p,y)h+D_yf(p,y)v_h(y)+f(p,y)\operatorname{div}v_h(y)\big)\,\dd y.
\]
Here $p$ is fixed in the definition of $v_h$. This is a formula for a bounded regular chart family, not an assertion about arbitrary weakly moving domains.
\end{theorem}
\begin{proof}
The region theorem \cref{d04:07:regions} makes every $E_p$ Jordan and gives
$F(p)=\int_Ef(p,\Phi(p,x))J(p,x)\,\dd x$. On the product of $\cl E$ with a sufficiently small compact parameter ball, $A$ remains invertible and its inverse is continuous and bounded: inversion is given by cofactors divided by the nonzero determinant. The hypotheses, the product and chain rules, and \cref{d04:08:jacobian-derivative} show that this pulled-back integrand and its parameter derivative are continuous. Apply \cref{d04:08:proper-derivative} to obtain the first formula and its continuous dependence on $p$.

Under the additional mixed-derivative hypothesis the spatial derivative of the velocity is
\[
 D_yv_h(\Phi(p,x))=(D_xD_p\Phi(p,x)h)A(p,x)^{-1}.
\]
Its trace equals $\tr(A^{-1}D_pA[h])$, since the two mixed derivatives agree and $\tr(BC)=\tr(CB)$ by summing matrix entries. Define divergence as this trace, and substitute back from $E$ to $E_p$ with \cref{d04:07:regions}. This proves the second formula. If desired, for the regular boundaries and orientations established in \cref{ch:d04-12}, its last two terms can be expressed as the boundary flux of $fv_h$ by that chapter's divergence theorem. No boundary theorem is used to prove the present result.
\end{proof}

\begin{example}[Expanding a region]\label{d04:08:expansion}
For $t>0$ let $E_t=tE$ and let $f$ be $C^1$ near the relevant closures. Taking $\Phi(t,x)=tx$ gives $v(y)=y/t$, $\operatorname{div}v=n/t$, and
\[
 \frac{\dd}{\dd t}\int_{tE}f(y)\,\dd y
 =\frac1t\int_{tE}\big(Df(y)y+nf(y)\big)\,\dd y.
\]
In particular the derivative of volume is $nt^{n-1}\vol(E)$, agreeing with linear volume scaling. A moving-domain formula with only $\int_{E_t}\partial_tf$ would miss this volume term even for $f=1$.
\end{example}

\section{Improper families and uniform tails}
\family{K04.08.07}
\begin{definition}\label{d04:08:uniform-tails}
Fix a regular exhaustion $K_j$ of open $\Omega$ as in \cref{d04:06:improper-definition}. For locally Riemann-integrable $f(\cdot,p)$ write $F_j(p)=\int_{K_j}f(x,p)\,\dd x$. The improper family converges uniformly on $B\subset P$ if $F_j$ converges uniformly there. Equivalently, for every $\eps>0$ some $j_0$ satisfies
\[
 \sup_{p\in B}\norm{\int_{K_k\setminus K_j}f(x,p)\,\dd x}<\eps
 \qquad(k\geq j\geq j_0),
\]
For scalar integrands the norm is an absolute value. Uniform absolute tails mean the stronger assertion with $\int_{K_k\setminus K_j}\norm{f(x,p)}$ in place of that norm. For multiple singular endpoints, each endpoint tail must be controlled separately.
\end{definition}

\begin{theorem}[Uniform improper limits and continuity]\label{d04:08:improper-continuity}
Suppose every $F_j$ is continuous on a parameter neighborhood $B$ and the improper family converges uniformly on $B$. Then its limit $F$ is continuous on $B$. More generally, if $f(\cdot,p)\to g$ uniformly on every $K_j$ as $p\to p_0$ and the improper family has uniform Cauchy tails near $p_0$, then $g$ has an improper integral along this exhaustion and
\[
 \lim_{p\to p_0}\int_\Omega f(x,p)\,\dd x=\int_\Omega g(x)\,\dd x.
\]
A parameter-independent, nonnegative, absolutely integrable locally Riemann function $M$ with $\norm{f(x,p)}\leq M(x)$ implies uniform absolute tails.
\end{theorem}
\begin{proof}
For continuity choose $j$ so that $\sup_B\norm{F-F_j}<\eps/3$, then use continuity of $F_j$ at the parameter under consideration; the three-term triangle inequality proves continuity. For the second assertion, uniform convergence on $K_j$ implies integrability of $g$ there and convergence of its proper integrals. Fix $k\geq j$ and pass $p\to p_0$ in the uniform tail inequality to obtain the same Cauchy bound for $\int_{K_k\setminus K_j}g$. Thus its proper integrals have a Cauchy limit. Compare the two improper integrals through their common $K_j$ piece; the two tails are uniformly small and the proper piece tends to zero, proving the equality. Finally domination bounds every tail by $\int_{K_k\setminus K_j}M\leq\int_\Omega M-\int_{K_j}M$, which tends to zero by \cref{d04:06:absolute-exhaustion}.
\end{proof}

\begin{lemma}[A uniform $C^1$ limit with vector parameters]\label{d04:08:c1-limit}
Let $B=\cl B(p_*,r)\subset\R^m$, and let $F_j$ be $C^1$ on a neighborhood of $B$ with values in $\R^k$. If $F_j(p_*)$ converges and $DF_j$ converges uniformly on $B$ to $G$, then $F_j$ converges uniformly on $B$ to $F$, and $F$ is $C^1$ on $\Int B$ with $DF=G$.
\end{lemma}
\begin{proof}
Along the segment from $p_*$ to $p$, scalar FTC applied to each target coordinate gives
\[
 F_j(p)-F_k(p)=F_j(p_*)-F_k(p_*)+
 \int_0^1(DF_j-DF_k)(p_*+t(p-p_*))(p-p_*)\,\dd t.
\]
The norm is at most $\norm{F_j(p_*)-F_k(p_*)}+r\sup_B\norm{DF_j-DF_k}$, so the functions are uniformly Cauchy. Their uniform limit exists because the target is complete. The uniform limit $G$ of continuous maps is continuous. On any short segment from $p$ to $p+h$ inside $B$, take limits in the analogous FTC identity, using the integral norm bound to justify uniform passage, obtaining
$F(p+h)-F(p)=\int_0^1G(p+th)h\,\dd t$. Subtract $G(p)h$ and use continuity of $G$ to bound the remainder by $\norm h\sup_{0\leq t\leq1}\norm{G(p+th)-G(p)}=o(\norm h)$. This is the required total derivative and its continuity.
\end{proof}

\begin{theorem}[Improper differentiation from derivative tails]\label{d04:08:improper-derivative}
On a compact parameter ball $B=\cl B(p_*,r)\subset P$, suppose $f$ meets the proper differentiation hypotheses on each $K_j$, so that
\[
 F_j\in C^1,\qquad DF_j(p)=\int_{K_j}D_pf(x,p)\,\dd x.
\]
Assume the derivative integrals converge uniformly on $B$ and $F_j(p_*)$ converges at one parameter. Then the improper integrals $F(p)$ converge uniformly on $B$, $F\in C^1(\Int B)$, and
\[
 DF(p)h=\lim_{j\to\infty}\int_{K_j}D_pf(x,p)h\,\dd x.
\]
A sufficient derivative-tail hypothesis is $\norm{D_pf(x,p)}\leq M(x)$ on $B$, with $M$ absolutely integrable and locally Riemann integrable. Uniform convergence of the signed derivative integrals also suffices without absolute domination. Repetition gives finite higher parameter derivatives when the corresponding truncated derivatives and their uniform tails satisfy the same conditions at each order.
\end{theorem}
\begin{proof}
The assumptions are exactly those of \cref{d04:08:c1-limit}; apply that lemma and the proper derivative formula. For the sufficient hypothesis,
\[
 \sup_{p\in B}\norm{DF_k(p)-DF_j(p)}
 \leq\int_{K_k\setminus K_j}M\longrightarrow0
\]
uniformly for $k\geq j$, so derivative convergence follows by completeness in the finite-dimensional operator space. The assertion for higher derivatives follows by applying the same lemma in the finite-dimensional spaces of multilinear maps, order by order. Each application requires the stated uniform tails and base-point convergence for the order to which it is applied; it is not a purely formal repeated interchange. Absolute convergence, when present, additionally gives exhaustion independence by \cref{d04:06:absolute-exhaustion}; the proof itself only needs the chosen exhaustion.
\end{proof}

\begin{proposition}[An improper interchange justified by uniform tails]\label{d04:08:improper-interchange}
Let $A$ be a compact box, $K_j$ a regular exhaustion of open $\Omega$, and suppose $f$ is continuous on $A\times\Omega$. If $H_j(x)=\int_{K_j}f(x,y)\,\dd y$ converges uniformly for $x\in A$, then
\[
 \int_A\left(\lim_j\int_{K_j}f(x,y)\,\dd y\right)\dd x
 =\lim_j\int_{K_j}\left(\int_Af(x,y)\,\dd x\right)\dd y.
\]
Both sides exist; the right side is the indicated improper integral. The assertion also holds for vector-valued $f$.
\end{proposition}
\begin{proof}
Each $H_j$ is continuous by the proper continuity theorem. On the bounded product $A\times K_j$, zero extend in $y$; its discontinuities are confined to $A\times\partial K_j$, a content-zero set by a product cover. Every inner slice is integrable. Proper Fubini gives equality of the two iterated integrals for fixed $j$. Uniform convergence gives an integrable continuous limit $H$ and $\int_AH_j\to\int_AH$ by the integral bound. Passing to the limit in the proper equalities proves both existence and equality as stated. Thus the interchange is justified by an explicit uniform-tail condition, not by reusing proper Fubini at an infinite boundary.
\end{proof}

\begin{example}[Exponential tails]\label{d04:08:exponential-tail}
On $x\geq0$ and parameters $a\geq a_0>0$, the derivatives of $e^{-ax}$ of order $j$ are $(-x)^je^{-ax}$ and are dominated by $x^je^{-a_0x}$. These are absolutely integrable: the scalar exponential series gives $x^j\leq j!(2/a_0)^j e^{a_0x/2}$, leaving an exponentially integrable bound. Thus every finite derivative order has uniform tails on compact subsets of $a>0$. Since $\int_0^\infty e^{-ax}\,\dd x=1/a$, differentiation gives
\[
 \int_0^\infty x^je^{-ax}\,\dd x=\frac{j!}{a^{j+1}}.
\]
The interval $a>0$ is the parameter domain; the tail bounds do not assert uniformity as $a\downarrow0$.
\end{example}

\family{K04.08.08}
\section{Failures of insufficient hypotheses}
\begin{example}[Pointwise convergence on a proper interval]\label{d04:08:proper-counterexample}
For $0<t<1$ set $f(x,t)=t^{-1}\mathbf1_{(0,t)}(x)$ on $[0,1]$ and $f(x,0)=0$. For every fixed $x$, $f(x,t)\to0$ as $t\downarrow0$, but $\int_0^1f(x,t)\,\dd x=1$. Each proper slice is integrable, so pointwise convergence alone cannot replace the uniform hypothesis. Nor does altering a null set preserve arbitrary outer Riemann integrals, as \cref{d04:06:exceptions} already demonstrates.
\end{example}

\begin{example}[Escape of mass at an improper boundary]\label{d04:08:escape}
Choose a nonnegative $C^1$ function $\psi$ supported on $[1,2]$ with $\int\psi=1$; for example normalize $(u-1)^2(2-u)^2$ on that interval and set it to zero outside. Put $f(x,t)=t\psi(tx)$ for $t>0$, $x\geq0$, and $f(x,0)=0$. On every bounded $x$ interval, $f(\cdot,t)$ is identically zero for small enough $t$, and hence converges uniformly to zero there. Yet scalar substitution gives $\int_0^\infty f(x,t)\,\dd x=1$ for every $t>0$. The support moves to infinity, preventing uniform tails. For each fixed $x$, the parameter derivative at $t=0$ is zero as well; the parameter integral is not even continuous at $0$. Thus good behavior on every fixed truncation does not justify improper parameter differentiation.
\end{example}

\begin{example}[A translated center and its tails]\label{d04:08:moving-singularity}
The scalar integrals $\int_{\R}(1+(x-p)^2)^{-1}\,\dd x=\pi$ have their mass centered at the moving point $p$. On each compact parameter set, comparison with $C/(1+x^2)$ supplies uniform tails. On all of $\R$ as parameter set there are no uniform tails relative to $[-R,R]$: choose $p$ far beyond $R$, and almost all the mass lies outside that interval. Uniformity is therefore always attached to a specified parameter set, even when the value of the integral happens to be constant. These integrands have no finite singularity. For actual finite moving singularities a fixed exhaustion may not avoid the singularities uniformly at all; one must supply a change of variables with proved bounds, or a different explicit improper-limit argument.
\end{example}

\begin{remark}[The measure alternative and the finite-dimensional scope]\label{d04:08:measure-alternative}
An integrable dominating function can replace uniform convergence in a Lebesgue dominated-convergence argument when the functions involved are measurable. That is a distinct measure-theoretic theorem, available only under the comparison foundation \cref{base:measure}; it does not manufacture Riemann integrability of a pointwise limit. For Riemann functions already known to satisfy \cref{d04:06:measure-bridge}, one can compare the resulting values, but a measurable non-Riemann limit remains a measure integral. Full general parameter theorems in Banach spaces belong to functional analysis; weak transport, rough boundaries, and flow dependence require their own measure/PDE/ODE foundations. The proper and improper finite-dimensional results proved here use only the hypotheses displayed above.
\end{remark}

% END ORIGINAL INPUT: chapters/08.tex


% BEGIN ORIGINAL INPUT: chapters/09.tex
\chapter{Parametrized curves and line integrals}
\label{ch:d04-09}
\source{FLE Sections 6.2--6.4, especially printed pp. 247--255; SPI Chapter 4, printed pp. 100--104.}

\section{Paths, length and traversal}
\family{K04.09.01}\family{K04.09.02}\family{K04.09.03}
\begin{definition}[Paths and polygonal length]\label{d04:09:paths}
A path is a continuous map $\gamma:[a,b]\to\R^n$. Its image is its
trace, which forgets the order and number of visits. The path is piecewise
$C^1$ if a finite partition $a=t_0<\cdots<t_m=b$ makes each restricted
path $C^1$ up to its endpoints. Derivatives at the joins are understood
one-sided. It is piecewise regular if the derivative on each such closed
piece is nonzero. A simple arc is injective; a simple closed path is
injective except that its two endpoints agree. Constant paths are allowed; on a nondegenerate interval $a<b$, they are not piecewise regular.
For $a=b$, use the zero-piece partition: the piecewise $C^1$ and piecewise regular conditions hold vacuously, and all integrals and lengths below are zero.
Arc-length parametrization conclusions requiring a nondegenerate interval explicitly assume $a<b$.

For a partition $P$ put
$\ell_P(\gamma)=\sum_j|\gamma(t_j)-\gamma(t_{j-1})|$.
The length is $L(\gamma)=\sup_P\ell_P(\gamma)$, possibly infinite.
Refining a partition increases its polygonal length by the triangle inequality.
\end{definition}

\begin{theorem}[Length equals the speed integral]\label{d04:09:length}
Every piecewise $C^1$ path has finite length and
\[
 L(\gamma)=\sum_{j=1}^m\int_{t_{j-1}}^{t_j}|\gamma'(t)|\,dt.
\]
For partitions containing the joins, polygonal lengths converge to this
number as mesh tends to zero.
\end{theorem}
\begin{proof}
For one $C^1$ piece, the componentwise scalar fundamental theorem gives
$\gamma(v)-\gamma(u)=\int_u^v\gamma'(t)dt$. The norm of this integral is
at most the integral of its norm: the inequality holds for Riemann sums
by the triangle inequality and passes to their limit. Thus every polygonal
sum is at most $\int_a^b|\gamma'|$.
Let $\omega(\delta)$ be the modulus of continuity of $\gamma'$ on the
compact interval. For any tag $s_j$ in a subinterval of width $\Delta t_j$,
\[
 |\gamma(t_j)-\gamma(t_{j-1})-\gamma'(s_j)\Delta t_j|
 \le\omega(\operatorname{mesh}P)\Delta t_j.
\]
The reverse triangle inequality bounds the difference between the polygonal
sum and the tagged Riemann sum for $|\gamma'|$ by
$(b-a)\omega(\operatorname{mesh}P)$. This tends to zero, proving the
claim for a piece. In the piecewise case refine by all joins and apply the
argument on each of the finitely many pieces. Refinement shows that no
unrefined polygon has a larger supremum. This also proves length additivity
under restriction at any parameter value.
\end{proof}

\begin{proposition}[Arc-length parametrization]\label{d04:09:arclength}
For a regular $C^1$ path put $s(t)=\int_a^t|\gamma'(u)|du$ and $L=s(b)$.
If $a<b$, then $s$ is a strictly increasing $C^1$ bijection onto $[0,L]$
with $C^1$ inverse, and $\widetilde\gamma=\gamma\circ s^{-1}$ satisfies
$|\widetilde\gamma'(s)|=1$. Piecewise regular paths have the corresponding
piecewise unit-speed parametrization, with possible tangent jumps at joins.
\end{proposition}
\begin{proof}
The fundamental theorem gives $s'=|\gamma'|>0$. Strict monotonicity and
the intermediate-value theorem give its continuous inverse. The scalar
inverse derivative follows from the quotient
$(s^{-1}(v)-s^{-1}(u))/(v-u)$ and the definition of $s'$; it equals
$1/s'(s^{-1}(u))$ and is continuous. The chain rule now gives unit speed.
Each piecewise inverse joins continuously at the cumulative lengths, proving
the last statement. A stationary interval would make $s$ constant and is
the reason regularity is required for the inverse assertion.
\end{proof}

\section{Densities, work and reparametrization}
\family{K04.09.04}\family{K04.09.05}\family{K04.09.06}
\begin{definition}[Two different line integrals]\label{d04:09:integrals}
For a piecewise $C^1$ path and a scalar function $f$ on its trace define
\[
 \int_\gamma f\,ds=\sum_j\int_{t_{j-1}}^{t_j}f(\gamma(t))|\gamma'(t)|dt,
\]
provided the displayed bounded Riemann integrands are integrable. This
holds, in particular, if $f$ is continuous on the trace. For a continuous
one-covector field $\alpha_x$, define
\[
 \int_\gamma\alpha=\sum_j\int_{t_{j-1}}^{t_j}
       \alpha_{\gamma(t)}(\gamma'(t))dt.
\]
In Euclidean coordinates $\alpha=\sum_iF_i dx_i$ corresponds to the vector
field $F$ by $\alpha_x(v)=F(x)\cdot v$; the second integral is also denoted
$\int_\gamma F\cdot d x$. The first uses positive length and the second
directed displacement. Common refinements and additivity of ordinary
integrals prove that the definitions do not depend on the chosen join partition.
\end{definition}

\begin{theorem}[Reparametrization]\label{d04:09:reparametrize}
Let $h:[c,d]\to[a,b]$ be a piecewise $C^1$ bijection whose inverse is
piecewise $C^1$, with derivative nonzero on its regular pieces. It is either
strictly increasing or strictly decreasing. For continuous data,
\[
 L(\gamma h)=L(\gamma),\qquad
 \int_{\gamma h}f\,ds=\int_\gamma f\,ds,\qquad
 \int_{\gamma h}\alpha=\begin{cases}
 \int_\gamma\alpha,&h\text{ increasing},\\
 -\int_\gamma\alpha,&h\text{ decreasing}.
 \end{cases}
\]
The same formulas hold for bounded Riemann data whenever the substitution
hypotheses guarantee the indicated integrability on both parameter intervals.
\end{theorem}
\begin{proof}
A continuous injective map of an interval is strictly monotone: if three
successive arguments violate monotonicity, the intermediate-value theorem
produces a repeated value. Subdivide at all joins and their preimages.
The chain rule gives $(\gamma h)'=\gamma'(h)h'$ on every piece.
For scalar densities its norm contributes $|h'|$; the ordinary scalar
substitution formula gives exactly the same nonnegative length variable in
both monotonicity cases. For the one-form integral the factor is $h'$, so
the decreasing case instead reverses the integration endpoints and gives
a minus sign. Apply the scalar formula on each piece and sum. Taking
$f=1$ and using \cref{d04:09:length} proves the length assertion as well.
\end{proof}

\family{K04.09.07}
\begin{proposition}[Reversal, concatenation and estimates]\label{d04:09:path-algebra}
Let $\gamma^{-}(t)=\gamma(a+b-t)$. Scalar integrals are unchanged under
reversal and one-form integrals are negated. If the terminal point of
$\gamma_1$ is the initial point of $\gamma_2$, their concatenation (after
affine rescaling to two adjacent intervals) satisfies
\[
 L(\gamma_1*\gamma_2)=L(\gamma_1)+L(\gamma_2),\quad
 \int_{\gamma_1*\gamma_2}f\,ds=\int_{\gamma_1}f\,ds+\int_{\gamma_2}f\,ds,
 \quad\int_{\gamma_1*\gamma_2}\alpha=\int_{\gamma_1}\alpha+\int_{\gamma_2}\alpha.
\]
Moreover $|\int_\gamma f\,ds|\le\sup_\gamma|f|\,L(\gamma)$ and
$|\int_\gamma\alpha|\le\sup_\gamma\|\alpha\|\,L(\gamma)$.
\end{proposition}
\begin{proof}
Reversal is the decreasing case of the preceding theorem. Split the
concatenation integral at its join and undo its two increasing affine
parameter changes. Length additivity follows either from the same argument
with $f=1$ or directly from refined polygonal sums. The inequalities follow
from $|\alpha_x(v)|\le\|\alpha_x\||v|$ and the scalar integral inequality.
\end{proof}

\section{The fundamental theorem and a circulation test}
\family{K04.09.08}\family{K04.09.09}
\begin{theorem}[Gradient integral]\label{d04:09:gradient-integral}
If $f\in C^1(U)$ and the piecewise $C^1$ path $\gamma$ lies in the open
set $U$, then
\[
 \int_\gamma df=\int_\gamma\nabla f\cdot dx=f(\gamma(b))-f(\gamma(a)).
\]
This is the path-endpoint part of the canonical potential equivalence in
\cref{d04:11:potential}; the metric correspondence is
$df_x(v)=\ip{\nabla f(x)}{v}$.
\end{theorem}
\begin{proof}
On a smooth piece the chain rule gives
$(f\circ\gamma)'(t)=Df(\gamma(t))\gamma'(t)
=\nabla f(\gamma(t))\cdot\gamma'(t)$.
The scalar fundamental theorem evaluates its integral. The endpoint values
at internal joins cancel in the finite sum.
\end{proof}

\begin{example}[Multiplicity changes both kinds of integral]\label{d04:09:circle}
For $\gamma_m(t)=(\cos mt,\sin mt)$, $0\le t\le2\pi$, and a nonzero
integer $m$, all traces are the unit circle but $L(\gamma_m)=2\pi|m|$.
For the angular field $F=(-y,x)/(x^2+y^2)$,
$F(\gamma_m)\cdot\gamma_m'=m$, so the circulation is $2\pi m$.
Direct differentiation gives $\partial_xF_2=\partial_yF_1$ on punctured
space, yet the field has no global scalar potential there: a gradient has
zero integral on every closed path by \cref{d04:09:gradient-integral}.
The scalar integral $\int_{\gamma_m}1\,ds$ records absolute traversal
multiplicity, and the work integral records signed traversal. Neither is
determined solely by the trace without a single-traversal convention.
\end{example}

\begin{example}[A helical arc]\label{d04:09:helix}
For $\gamma(t)=(a\cos t,a\sin t,bt)$ on $[0,T]$, with $a>0$ and $T\ge0$, the speed
is $\sqrt{a^2+b^2}$, so $L=T\sqrt{a^2+b^2}$ and arc length is
$s=t\sqrt{a^2+b^2}$. For $f(x,y,z)=z$, the work of $\nabla f=(0,0,1)$
is $bT$, whereas $\int_\gamma z\,ds=bT^2\sqrt{a^2+b^2}/2$.
These values exhibit the two different geometric factors explicitly.
\end{example}

% END ORIGINAL INPUT: chapters/09.tex


% BEGIN ORIGINAL INPUT: chapters/10.tex
\chapter{Surface area and oriented flux}
\label{ch:d04-10}
\source{SPI Chapter 5, printed pp. 126--130; FLE Sections 8.1--8.4 and 8.8; MUN Chapters 5--6, with the chart transition passage in Section 24 checked visually.}

\section{Patches and their geometric density}
\family{K04.10.01}\family{K04.10.02}\family{K04.10.03}
\begin{definition}[Regular patches]\label{d04:10:patch}
A regular $C^1$ surface parametrization is a map
$\phi:U\subset\R^2\to\R^n$ with rank two everywhere. It is an embedded
patch if it is a homeomorphism onto its image; full rank alone only gives
this after local restriction, by \cref{d04:14:charts}. At $p=\phi(u,v)$ its
tangent plane is $\operatorname{im}D\phi(u,v)$, which is independent of an
embedded chart by \cref{d04:14:tangent}. For $n=3$ the two choices of unit
normal are
\[
 \pm N_\phi=\pm\frac{\phi_u\times\phi_v}{|\phi_u\times\phi_v|}.
\]
The cross product is nonzero precisely when the two tangent columns are
independent, so this formula states exactly the rank condition.
\end{definition}

\begin{definition}[Area integral]\label{d04:10:area}
Write $G_\phi=D\phi^TD\phi$ and $J_2\phi=\sqrt{\det G_\phi}$.
For compact Jordan measurable $D\subset U$ and continuous $f$ on
$\phi(D)$ put
\[
 \int_{\phi|D} f\,dA=\int_D f(\phi(u,v))J_2\phi(u,v)\,du\,dv.
\]
More general bounded data are allowed only when the indicated Riemann
integrand is integrable. The value with $f=1$ is the parametrized area.
For a noninjective parametrization it counts its traversal, and is not
automatically the area of the image once. This definition uses the metric
area density, rather than a supremum over arbitrary inscribed triangles.
\end{definition}

\begin{proposition}[Gram and cross-product formulas]\label{d04:10:gram}
For a regular patch in $\R^3$,
\[
 J_2\phi=\sqrt{|\phi_u|^2|\phi_v|^2-(\phi_u\cdot\phi_v)^2}
          =|\phi_u\times\phi_v|.
\]
For an affine patch it is the area of the parallelogram spanned by its
derivative columns. More generally for a $k$-patch in $\R^n$ the factor
$J_k\phi=\sqrt{\det(D\phi^TD\phi)}$ is the $k$-dimensional
parallelepiped volume factor.
\end{proposition}
\begin{proof}
The two-by-two Gram determinant gives the displayed scalar expression.
Expanding the three cross-product components and summing their squares
gives the same expression. In general choose an orthonormal basis of the
column image and let $A$ be the square matrix of the columns in this basis.
Then the Gram matrix is $A^TA$, and its determinant's square root is
$|\det A|$. The linear-volume theorem \cref{d04:07:linear-volume}, applied
inside that orthonormal image space, identifies this number with its
Euclidean parallelepiped volume. Thus the definition has the expected flat normalization.
It agrees with speed for $k=1$ and absolute volume Jacobian for $k=n$.
\end{proof}

\family{K04.10.04}\family{K04.10.05}\family{K04.10.06}
\begin{theorem}[Actual reparametrization invariance]\label{d04:10:reparametrize}
Let $h:V\to U$ be a $C^1$ diffeomorphism and $\psi=\phi h$.
For compact Jordan $E\subset V$, $D=h(E)$ is Jordan measurable and
\[
 \int_{\psi|E}f\,dA=\int_{\phi|D}f\,dA.
\]
For patches in $\R^3$,
\[
 \psi_s\times\psi_t=(\det Dh)(\phi_u\times\phi_v)\circ h.
\]
Consequently oriented flux
\[
 \int_{\phi|D}F\cdot N_\phi\,dA
 =\int_D F(\phi)\cdot(\phi_u\times\phi_v)\,du\,dv
\]
is unchanged under positive reparametrization and changes sign under
negative reparametrization when the parameter orientations are used.
For a fixed geometric normal, the chart's sign must be adjusted so that
both computations use that same normal; then flux is invariant.
\end{theorem}
\begin{proof}
The chain rule gives $D\psi=(D\phi\circ h)Dh$, whence
$G_\psi=Dh^T(G_\phi\circ h)Dh$ and
$J_2\psi=|\det Dh|(J_2\phi\circ h)$.
The ordinary region substitution theorem \cref{d04:07:regions} proves both
Jordan measurability and equality of the scalar integrals; it also supplies
the integrability equivalence for bounded Riemann data when these are used.
For the cross product expand each of the two new tangent columns as a
linear combination of the old columns. Alternation removes the repeated
column terms and leaves $\det Dh$ times the old cross product. Inserting
this into flux and applying substitution proves the signed assertion.
A many-to-one parameter map is not a diffeomorphism, so this proof does
not remove its multiplicity.
\end{proof}

\section{Assembly, orientation and exceptional parameters}
\family{K04.10.07}\family{K04.10.08}\family{K04.10.09}
\begin{proposition}[Overlapping surface patches]\label{d04:10:assembly}
A compact embedded $C^1$ surface, with or without boundary, has a
well-defined scalar integral of every continuous function, obtained by
integrating the Gram density in finitely many charts with a subordinate
partition of unity. It has a global oriented flux integral if a continuous
unit normal has been chosen. Overlap weights cannot be discarded unless
the patches instead form a decomposition with only area-zero intersections.
\end{proposition}
\begin{proof}
Use the finite partition of \cref{d04:14:partition}. On a chart overlap
the Gram density transformation is precisely the absolute-Jacobian identity
proved in \cref{d04:10:reparametrize}. The double-partition argument in
\cref{d04:15:independence} therefore proves independence of weights and
charts for scalar densities. For a chosen normal, use the signed chart
adjustment in the same reparametrization theorem; it is equivalently
integration of the two-form $\beta_F$ of \cref{d04:13:dictionary}.
Adding whole overlapping charts without weights would add the integral
over their common open portion twice, so the partition is essential.
\end{proof}

\begin{lemma}[Negligible seams]\label{d04:10:seams}
A finite union of compact $C^1$ curves and points on an embedded surface
has area zero. It can be omitted when computing an integral of a bounded
continuous function by regular charts that cover the complement, provided
the limiting integrals are taken over regular compact truncations.
\end{lemma}
\begin{proof}
Cover the compact exceptional set by finitely many restricted surface
charts. On each compact parameter subset the inverse chart is $C^1$ and
locally Lipschitz by \cref{d04:14:charts}; thus a curve lifts, after finite
subdivision, to a Lipschitz curve in $\R^2$. Subdivide its parameter interval
into $O(1/\delta)$ intervals of width $\delta$. Their images fit in squares
of side $C\delta$, with total area $O(\delta)\to0$.
Points have the same property using a single small square. Gram densities
and finitely many partition coefficients are bounded on the compact chart
subsets, so these square covers give surface neighborhoods of arbitrarily
small density integral. The integral of a bounded function over them is
bounded by its supremum times their area. A compact truncation exhausting
the complement eventually contains every closed part outside such a
neighborhood; this estimate proves convergence and equality. No arbitrary
change on a dense null set is being licensed.
\end{proof}

\begin{example}[The Mobius strip]\label{d04:10:mobius}
For $R>a>0$, consider
\[
 \phi(u,v)=((R+v\cos(u/2))\cos u,
 (R+v\cos(u/2))\sin u,v\sin(u/2)),\quad |v|\le a.
\]
Use $0\le u\le2\pi$ and identify $(0,v)$ with $(2\pi,-v)$.
The radius $R+v\cos(u/2)>0$, so polar angle determines $u$ modulo $2\pi$,
and the remaining coordinates determine $v$ with exactly that seam
identification. This gives an embedded strip covered by regular local
patches, including seam patches obtained by extending the formula in $u$.
In the orthonormal radial, angular and vertical frame, $\phi_v$ has norm
one, is orthogonal to $\phi_u$, and
$|\phi_u|^2=(R+v\cos(u/2))^2+v^2/4>0$. Thus the rank is two.
Along $v=0$, the cross-product normal at $u=2\pi$ is the negative of its
value at $u=0$. A hypothetical continuous unit normal must on this entire
connected center curve equal either $N_\phi$ or $-N_\phi$, with a constant
choice of sign. It would return with the wrong value at the same point,
a contradiction. The strip is nonorientable. Its positive Gram area
density nevertheless exists and can be integrated by the preceding
assembly theorem. A globally normal-oriented flux integral is unavailable
until the geometry is replaced by a suitable oriented parametrized chain.
\end{example}

\section{Calculations and calibration}
\family{K04.10.10}
\begin{example}[Graphs]\label{d04:10:graph}
For $\phi(u,v)=(u,v,g(u,v))$,
$\phi_u\times\phi_v=(-g_u,-g_v,1)$ and
$dA=\sqrt{1+g_u^2+g_v^2}\,du\,dv$.
The upward flux of $F$ is therefore
$\int_D[-F_1(\phi)g_u-F_2(\phi)g_v+F_3(\phi)]du\,dv$.
For the paraboloid $g=u^2+v^2$ above $u^2+v^2\le a^2$,
polar substitution on regular sectors gives
\[
 A=2\pi\int_0^a r\sqrt{1+4r^2}\,dr
   =\frac\pi6[(1+4a^2)^{3/2}-1].
\]
The radial origin and a ray in the polar parameter representation are
removed by truncation; their planar and surface seams are negligible by
\cref{d04:10:seams}.
\end{example}

\begin{example}[Sphere and flux]\label{d04:10:sphere}
The sphere of radius $R>0$ is parametrized off a seam and its poles by
\[
 \phi(\theta,\varphi)
 =(R\sin\theta\cos\varphi,R\sin\theta\sin\varphi,R\cos\theta),
 \quad 0<\theta<\pi,\quad0<\varphi<2\pi.
\]
Its cross product in this order is
$\phi_\theta\times\phi_\varphi=R^2\sin\theta\,n$, with $n=\phi/R$
outward. This follows either by direct expansion or by orthogonality of
the tangent vectors, of lengths $R$ and $R\sin\theta$, with sign checked
at $(\theta,\varphi)=(\pi/2,0)$ by continuity. Consequently
\[
 A=\int_0^{2\pi}\int_0^\pi R^2\sin\theta\,d\theta\,d\varphi=4\pi R^2,
 \qquad\int_{S_R} x\cdot n\,dA=4\pi R^3.
\]
Use compact truncations before the scalar integrations. The omitted
meridian and poles form a compact curve and points, so
\cref{d04:10:seams} proves that the limits are the whole-sphere integrals.
The displayed parametrization is not a regular global chart at the poles
and is not injective across its closed seam.
\end{example}

\begin{example}[Ruled surfaces]\label{d04:10:ruled}
For $\phi(s,t)=a(s)+t b(s)$, the tangent columns are $a'(s)+t b'(s)$ and
$b(s)$, so the area factor is
$|(a'(s)+t b'(s))\times b(s)|$ wherever it is nonzero. A cylinder
$(R\cos s,R\sin s,t)$ has factor $R$ and lateral area $2\pi RH$ for
$0\le t\le H$, after removing its seam. A cone parametrized by
$((1-t)R\cos s,(1-t)R\sin s,tH)$ has factor
$R(1-t)\sqrt{R^2+H^2}$ and lateral area $\pi R\sqrt{R^2+H^2}$.
At its apex $t=1$ rank fails; the calculation is the limit of regular
truncations, with the single apex point contributing zero area. This is an
area computation on regular pieces, not a claim that the apex has a tangent plane.
\end{example}

% END ORIGINAL INPUT: chapters/10.tex


% BEGIN ORIGINAL INPUT: chapters/11.tex
\chapter{Vector operators, potentials and topology}
\label{ch:d04-11}
\source{FLE Sections 6.3--6.4, 8.4 and 8.8; SPI Chapter 4, printed pp. 88--96; HUB Section 6.13, printed pp. 688--695.}

\section{Euclidean operators and identities}
\family{K04.11.01}\family{K04.11.02}\family{K04.11.03}
\begin{definition}[Gradient, divergence and curl]\label{d04:11:operators}
For $f\in C^1(U)$ the gradient is the unique vector satisfying
$Df(x)v=\ip{\nabla f(x)}v$ for every $v$; hence its column coordinates
are $(\partial_1 f,\ldots,\partial_n f)^T$.
For a $C^1$ vector field $F:U\to\R^n$ define
$\operatorname{div}F=\tr DF=\sum_i\partial_iF_i$.
In oriented $\R^3$ define
\[
 \operatorname{curl}F=(\partial_yF_3-\partial_zF_2,
 \partial_zF_1-\partial_xF_3,\partial_xF_2-\partial_yF_1).
\]
The curl convention uses the standard orientation: it is the unique vector
with $d\alpha_F=\beta_{\operatorname{curl}F}$ in
\cref{d04:13:dictionary}. A change of orientation changes this vector
identification. For $f\in C^2$, the Laplacian is
$\Delta f=\operatorname{div}\nabla f=\sum_i\partial_i^2f$.
\end{definition}

\family{K04.11.04}\family{K04.11.05}
\begin{proposition}[Product and differential identities]\label{d04:11:identities}
For $C^1$ data of the indicated types,
\[
 \nabla(fg)=f\nabla g+g\nabla f,\qquad
 \operatorname{div}(fF)=\nabla f\cdot F+f\operatorname{div}F.
\]
In $\R^3$,
\[
 \operatorname{curl}(fF)=\nabla f\times F+f\operatorname{curl}F,
 \qquad\operatorname{div}(F\times G)=G\cdot\operatorname{curl}F-F\cdot\operatorname{curl}G.
\]
For $C^2$ data,
\[
 \operatorname{curl}\nabla f=0,\quad
 \operatorname{div}\operatorname{curl}F=0,\quad
 \operatorname{curl}\operatorname{curl}F
 =\nabla\operatorname{div}F-\Delta F.
\]
Here $\Delta F$ means applying $\Delta$ to each component.
\end{proposition}
\begin{proof}
The first two formulas are the scalar product rule applied to each
coordinate and then summed. For the first curl component,
\[
 \partial_y(fF_3)-\partial_z(fF_2)
 =f(\partial_yF_3-\partial_zF_2)+(\partial_yf)F_3-(\partial_zf)F_2;
\]
the other two are its cyclic permutations. Expanding
$\sum_i\partial_i(F\times G)_i$ gives the six terms differentiated in $F$
as $G\cdot\operatorname{curl}F$ and the six differentiated in $G$ as
$-F\cdot\operatorname{curl}G$.
The curl of a gradient has components of the form
$\partial_i\partial_jf-\partial_j\partial_if$, hence zero by continuous
mixed-partial symmetry. The divergence of a curl cancels the same three
pairs. For the iterated curl its first component is
\[
 \partial_y(\partial_xF_2-\partial_yF_1)
 -\partial_z(\partial_zF_1-\partial_xF_3)
 =\partial_x(\operatorname{div}F)-\Delta F_1.
\]
Cyclic permutation gives the others. These are also the vector shadows
of $d^2=0$ for the first two identities, with precisely the regularity in
\cref{d04:13:d-rules}.
\end{proof}

\section{The canonical potential theorem}
\family{K04.11.06}\family{K04.11.07}\family{K04.11.08}\family{K04.11.09}
\begin{theorem}[Potential, path independence and periods]\label{d04:11:potential}
Let $U\subset\R^n$ be open and connected, and let $F:U\to\R^n$ be
continuous. The following statements are equivalent:
\begin{enumerate}[label=(\roman*)]
\item $F=\nabla f$ for some $f\in C^1(U)$;
\item work integrals along piecewise $C^1$ paths in $U$ depend only on
their endpoints;
\item the work integral on every closed piecewise $C^1$ path is zero.
\end{enumerate}
The potential is unique up to one additive constant. If $F\in C^r$, then
the potential is $C^{r+1}$ for $r\ge0$ (where $C^0$ means continuous).
For a disconnected open set the statements apply on each component, with
one independent constant per component.
\end{theorem}
\begin{proof}
An open connected Euclidean set is polygonally connected: the points
reachable from a fixed point by finite segments in the set form an open
subset, since a small ball allows one additional segment; its complement
is also open by the same argument. Connectedness forces reachability of all
points. The implication (i)$\Rightarrow$(ii) is exactly the chain-rule
calculation of \cref{d04:09:gradient-integral}. Statement (ii) gives (iii)
by comparison with the constant path. If (iii) holds, concatenate one path
with the reverse of another and use \cref{d04:09:path-algebra} to obtain
(ii).

Fix $a\in U$. Define $f(x)$ as the work on any polygonal path from $a$ to
$x$; (ii) makes this well-defined. If a small ball about $x$ lies in $U$,
the segment from $x$ to $x+h$ gives
\[
 f(x+h)-f(x)=\int_0^1 F(x+th)\cdot h\,dt.
\]
Subtract $F(x)\cdot h$. Its absolute value is at most
$|h|\sup_{0\le t\le1}|F(x+th)-F(x)|=o(|h|)$ by continuity. Thus $Df(x)$
is the covector $v\mapsto F(x)\cdot v$, continuously in $x$, giving (i).
Two potentials have zero derivative difference and hence are constant on
every segment, then on every polygonal path. Higher regularity follows
because the first partials of $f$ are the $C^r$ functions $F_i$.
\end{proof}

\begin{corollary}[Local and star-shaped potentials]\label{d04:11:local-potential}
For $F\in C^1(U)$, the symmetry conditions
$\partial_iF_j=\partial_jF_i$ are necessary for a potential. They are
sufficient locally on every ball in $U$, and globally if $U$ is star-shaped
about $a$. In that case
\[
 f(x)=\int_0^1F(a+t(x-a))\cdot(x-a)\,dt
\]
is a potential, up to a constant. In $\R^3$ the symmetry conditions are
exactly $\operatorname{curl}F=0$.
\end{corollary}
\begin{proof}
A potential for $C^1$ data is $C^2$ by \cref{d04:11:potential}, so the
necessary condition follows from equality of mixed partials. In form
notation the condition is $d\alpha_F=0$. Apply the explicit Poincare
primitive \cref{d04:13:poincare} in degree one. It gives the displayed
formula and $df=\alpha_F$; a ball is star-shaped. The coordinate curl
formula identifies the three independent symmetry conditions in dimension
three.
\end{proof}

\begin{theorem}[The simply connected one-form case]\label{d04:11:simply-connected}
If $U$ is open, connected and simply connected, a $C^1$ field with symmetric
Jacobian has a global potential. Here simply connected means that every
continuous loop admits a continuous homotopy to a constant loop inside
$U$; the homotopy's two endpoint tracks agree, so its intermediate curves
are loops. This agrees with the based-contraction convention in
\cref{ch:d04-01}: given a free contraction, put $c(s)=H(s,0)$ and precede
the loop $H(s,\cdot)$ by the path $c|_{[0,s]}$ from the original basepoint,
then follow it by the reverse of that path. These are continuously based
loops. At the final stage the loop is a path followed by its reverse,
which contracts while keeping the basepoint fixed by shortening that path.
The initial harmless stationary pieces are removed by a continuous
reparametrization. Conversely a based contraction is a free contraction.
No differentiability of a homotopy is required.
\end{theorem}
\begin{proof}
By \cref{d04:11:local-potential} there are local potentials on balls.
Take a closed piecewise $C^1$ path $\gamma$ and a continuous contraction
$H:[0,1]^2\to U$, with $H(0,t)=\gamma(t)$, $H(1,t)$ constant and
$H(s,0)=H(s,1)$. Cover the compact image by balls lying in $U$ with local
potentials. The Lebesgue-number property and uniform continuity allow a
finite rectangular grid in the square such that the image of every grid
cell lies in one such ball. Include all breakpoints of $\gamma$ in the
horizontal grid. Replace each grid edge by the straight segment joining
the images of its endpoints. This segment lies in each adjacent cell's
ball: both endpoints lie there and balls are convex. In any grid cell the
sum of the four directed segment integrals is zero, because a local
potential evaluates each by endpoint differences. Sum over all cells;
interior edges cancel. Along $s=0$, each original path subarc and its
replacement chord have the same integral in their common ball. Along
$s=1$ all chords are constant, and the two endpoint tracks give identical
chords with opposite directions. Thus the only remaining value is
$\int_\gamma F\cdot dx=0$. The period criterion
\cref{d04:11:potential} supplies the global potential.
\end{proof}

\begin{example}[A domain matters]\label{d04:11:angular-obstruction}
The field $(-y,x)/(x^2+y^2)$ on punctured $\R^2$ has symmetric Jacobian,
but its unit-circle period is $2\pi$, by \cref{d04:09:circle}. It therefore
has local potentials and no global one. On a slit plane it has an angular
potential; any continuously chosen branch of polar angle differentiates to
that field by differentiating the polar-coordinate inverse. This is a
statement about a single-valued branch on that slit, not a globally defined
angle on punctured space.
\end{example}

\section{Vector potentials and higher-degree obstructions}
\family{K04.11.10}
\begin{theorem}[Vector potential on a star-shaped domain]\label{d04:11:vector-potential}
If $F\in C^1(U;\R^3)$, $U$ is star-shaped about $a$, and
$\operatorname{div}F=0$, then there is a $C^1$ vector potential $A$ with
$\operatorname{curl}A=F$. One choice is
\[
 A(x)=\int_0^1t\,F(a+t(x-a))\times(x-a)\,dt.
\]
A $C^2$ vector potential necessarily has divergence-free curl. A global
potential of either regularity need not exist on an arbitrary domain.
\end{theorem}
\begin{proof}
Use the closed two-form $\beta_F$ from \cref{d04:13:dictionary}.
The Poincare primitive $H\beta_F$ is a $C^1$ one-form with derivative
$\beta_F$. On a vector $v$ its value is
\[
 \int_0^1t F(a+t(x-a))\cdot((x-a)\times v)dt
 =\int_0^1t[F(a+t(x-a))\times(x-a)]\cdot v\,dt,
\]
so it is exactly $\alpha_A$ for the displayed field. The dictionary gives
$\operatorname{curl}A=F$. The necessary $C^2$ statement is
\cref{d04:11:identities}. The radial field of
\cref{d04:15:radial-obstruction} is divergence free but has nonzero sphere
flux, which excludes any global $C^1$ vector potential by Stokes.
\end{proof}

\begin{example}[Simply connected is insufficient in degree two]\label{d04:11:punctured-three}
The domain $\R^3\setminus\{0\}$ is simply connected, yet the closed
two-form in \cref{d04:15:radial-obstruction} is not exact.
Here is an elementary verification of the first assertion. Any continuous
loop has positive distance from zero. By uniform continuity it can be
approximated uniformly by a closed polygon within less than half this
distance; the straight homotopy to that polygon avoids zero. The finitely
many polygon edges each lie in some plane through zero (enlarge their
linear span to a plane if necessary). Choose a point $p$ outside their
finite union. Such a point exists: a product of their nonzero defining
linear polynomials is a nonzero polynomial and cannot vanish everywhere,
as follows by induction from the scalar polynomial fact. No edge point
is a negative multiple of $p$, so the straight homotopy from the polygon
to the constant point $p$ also avoids zero. This contracts the original
loop. Thus the one-form theorem cannot be applied indiscriminately to
other form degrees. These examples exhaust the topology needed here;
there is no general de Rham or homology construction in this chapter.
\end{example}

\section{A bounded curvilinear-coordinate calculus}
\family{K04.11.11}
\begin{theorem}[Orthogonal coordinate formulas]\label{d04:11:orthogonal}
Let $x=\phi(u_1,u_2,u_3)$ be an orientation-preserving $C^2$
diffeomorphism on an open parameter domain, with mutually orthogonal
nonzero columns. Write
$h_i=|\partial_i\phi|$, $e_i=\partial_i\phi/h_i$, and $J=h_1h_2h_3>0$.
For a $C^1$ scalar field let $\widetilde f=f\circ\phi$; for a $C^1$ vector field use
physical components $F_i=F(\phi(u))\cdot e_i(u)$. Then
\[
 (\nabla f)\circ\phi=\sum_i\frac{\partial_i\widetilde f}{h_i}e_i,
 \qquad (\operatorname{div}F)\circ\phi
 =\frac1J\sum_i\partial_i\left(\frac{J}{h_i}F_i\right),
\]
and the physical components of curl are
\[
 ((\operatorname{curl}F)\circ\phi)\cdot e_i
 =\frac{\partial_j(h_\ell F_\ell)-\partial_\ell(h_jF_j)}{h_jh_\ell}
 \quad ((i,j,\ell)=(1,2,3),(2,3,1),(3,1,2)).
\]
For $f\in C^2$,
\[
 (\Delta f)\circ\phi=\frac1J\sum_i
 \partial_i\left(\frac{J}{h_i^2}\partial_i\widetilde f\right).
\]
\end{theorem}
\begin{proof}
The chain rule says $\partial_i\widetilde f=h_i\nabla f\cdot e_i$,
which gives the gradient in the orthonormal basis. The pullback of the
work form is $\phi^*\alpha_F=\sum_i h_iF_i\,du_i$.
The pullback of Euclidean volume is $J\,du_1\wedge du_2\wedge du_3$.
Insertion of $F$ corresponds to insertion of its parameter vector
$\sum_i(F_i/h_i)\partial_{u_i}$, so
\[
 \phi^*\beta_F=\sum_i(-1)^{i-1}\frac{J}{h_i}F_i\,du_{\widehat i}.
\]
Applying $d$ gives the divergence formula by naturality. Applying $d$ to
the pulled-back work form and comparing the $du_j\wedge du_\ell$
coefficient with the preceding flux expression for $\operatorname{curl}F$
gives the curl formula, since $J/h_i=h_jh_\ell$. Finally insert the
gradient components in the divergence formula to obtain the Laplacian.
Only $C^2$ coordinate maps and the stated field derivatives were used.
\end{proof}

\begin{example}[Cylindrical and spherical factors]\label{d04:11:coordinates}
For $(x,y,z)=(r\cos\theta,r\sin\theta,z)$ with $r>0$ and a single
angle interval of length less than $2\pi$, the factors are $(1,r,1)$.
Thus
\[
 \operatorname{div}F=\frac1r\partial_r(rF_r)+\frac1r\partial_\theta F_\theta+\partial_zF_z,
 \qquad\Delta f=\frac1r\partial_r(r\partial_rf)+\frac1{r^2}\partial_\theta^2f+\partial_z^2f.
\]
In spherical parameters $(r,\theta,\varphi)$, $r>0$, $0<\theta<\pi$,
with a single azimuth branch, the factors are $(1,r,r\sin\theta)$.
Consequently
\[
 \operatorname{div}F=\frac1{r^2}\partial_r(r^2F_r)
 +\frac1{r\sin\theta}\partial_\theta(\sin\theta F_\theta)
 +\frac1{r\sin\theta}\partial_\varphi F_\varphi,
\]
\[
 \Delta f=\frac1{r^2}\partial_r(r^2\partial_rf)
 +\frac1{r^2\sin\theta}\partial_\theta(\sin\theta\partial_\theta f)
 +\frac1{r^2\sin^2\theta}\partial_\varphi^2 f.
\]
These are identities on regular coordinate domains. The axis, polar points,
origin and angle seams require another chart or a limit; a denominator
vanishing there does not extend the chart theorem to a singular chart.
Full Riemannian operator theory is outside this selected Euclidean bridge.
\end{example}

% END ORIGINAL INPUT: chapters/11.tex


% BEGIN ORIGINAL INPUT: chapters/12.tex
\chapter{Green, Gauss and classical Stokes}
\label{ch:d04-12}
\source{FLE Section 8.4, printed pp. 340--343, and Section 8.8, pp. 362--365; SPI Chapter 5, classical formulas after the volume-element construction; MUN Section 37, checked printed pp. 303--306.}

The general-dimensional proof is in \cref{d04:15:stokes}. This chapter
supplies the domain, metric and sign identifications that turn that theorem
into the classical formulas. In particular a vector formula here is not
just a differently named theorem: its normal, tangent and differential-form
integrands will be compared explicitly.

\section{Domains and the volume interface}
\begin{definition}[Classical regular domains]\label{d04:12:domains}
A bounded $C^1$ domain in $\R^n$, $n\ge1$, is a bounded open set $D$ such that near
each point of its ambient boundary there is a $C^1$ function $\rho$ with
$D=\{\rho<0\}$ locally, $\partial D=\{\rho=0\}$ locally, and
$D\rho\ne0$ on that level set. We permit finitely or multiply connected
domains and do not require $D$ itself to be connected. Its closure
$K=\cl D$ carries the standard ambient orientation. The outward unit normal
is $\nu=\nabla\rho/|\nabla\rho|$, independently of the local defining
function. A function $C^r$ on a closure means locally extendible $C^r$
to an open neighborhood. For the compact closures used here finitely many
local extensions and Euclidean cutoffs combine to a single extension near
the closure, as in the proof of \cref{d04:14:partition}.
\end{definition}

\begin{lemma}[Regular-domain and Riemann interfaces]\label{d04:12:region-interface}
For such a domain, $K$ is a compact $C^1$ embedded $n$-manifold with
boundary $\partial K=\partial D$. The region $K$ is Jordan measurable;
for continuous $f$,
\[
 \int_K f\,dx_1\wedge\cdots\wedge dx_n=\int_D f\,dx
 =\int_K f\,dx,
\]
where the first integral is the oriented manifold integral and the other
two are the bounded-region Riemann integrals. Its boundary normal agrees
with the outward-first manifold convention.
\end{lemma}
\begin{proof}
At a boundary point choose a coordinate with $\partial_i\rho\ne0$.
The map consisting of the other $n-1$ coordinates and $-\rho$ has invertible
derivative; the inverse theorem supplies a local $C^1$ coordinate map
carrying $K$ to a half-space. The points with $\rho>0$ are locally outside
$K$ and those with $\rho<0$ are interior, so intrinsic and ambient
boundaries coincide. Boundedness and closedness give compactness.
The boundary has a finite cover by compactly restricted $C^1$ graph charts.
Each such graph has content zero: divide its $(n-1)$-dimensional parameter
box into cubes of side $\delta$; a uniform derivative bound places each
image in an $n$-box of side $C\delta$, giving total volume $O(\delta)$.
For $n=1$ this is the finite set of boundary points. Thus the boundary is
content zero and $K$ is Jordan measurable by \cref{d04:06:region-test}. Its inclusion
or removal leaves the Riemann integral unchanged.

Use the compact chart partition of \cref{d04:14:partition}. In any
full-dimensional chart its orientation sign relative to the ambient
orientation is $\sigma_\phi=\sgn\det D\phi$. The signed pulled-back volume
coefficient is therefore
$\sigma_\phi(f\circ\phi)\det D\phi=(f\circ\phi)|\det D\phi|$.
The bounded compact-support substitution \cref{d04:07:substitution},
localized near half-space charts as in
\cref{d04:15:directed-substitution}, identifies each weighted chart integral
with the ordinary integral of $\lambda_i f$ on $K$. The zero extension is
Riemann integrable because it is continuous off the content-zero boundary.
Sum the finite weights to obtain the equality. Finally
$D\rho(\nu)=|\nabla\rho|>0$, so $\nu$ points to the exterior in the
half-space model and is exactly the outward vector used by the manifold
definition.
\end{proof}

\section{The planar formulas}
\family{K04.12.01}\family{K04.12.02}\family{K04.12.05}\family{K04.12.06}
\begin{theorem}[Green circulation and flux]\label{d04:12:green}
Let $D\subset\R^2$ be a bounded $C^1$ domain and let $P,Q$ be $C^1$ on
a neighborhood of $K=\cl D$. Give $\partial D$ the positive orientation
induced by $dx\wedge dy$. Then
\[
 \int_{\partial D}P\,dx+Q\,dy
 =\iint_D(\partial_xQ-\partial_yP)\,dx\,dy,
\]
\[
 \int_{\partial D}(P,Q)\cdot\nu\,ds
 =\int_{\partial D}P\,dy-Q\,dx
 =\iint_D(\partial_xP+\partial_yQ)\,dx\,dy.
\]
The same formulas hold for finite piecewise regular regions with the
decomposition and cancellation hypotheses of \cref{d04:15:piecewise}.
\end{theorem}
\begin{proof}
For $\alpha=Pdx+Qdy$, direct exterior differentiation gives
$d\alpha=(\partial_xQ-\partial_yP)dx\wedge dy$.
Apply the $C^1$ ambient Stokes theorem \cref{d04:15:stokes} to $K$ and
use \cref{d04:12:region-interface} and the line-integral definition.
If $\tau$ is the positively oriented unit boundary tangent, outward-first
means $\det[\nu,\tau]=1$, so
$\tau=(-\nu_y,\nu_x)$ and $\nu=(\tau_y,-\tau_x)$.
Thus $(P,Q)\cdot\nu\,ds=Pdy-Qdx$, with no extra minus sign.
Apply the circulation formula to the one-form $-Qdx+Pdy$ to obtain the
flux formula. For finite pieces apply the same computations piece by piece;
the internal directed terms cancel by \cref{d04:15:piecewise}.
\end{proof}

\begin{example}[Orientation around holes]\label{d04:12:annulus}
For the annulus $a<\sqrt{x^2+y^2}<b$ with $0<a<b$, the outer unit normal points radially outward
on the outer circle and radially inward on the inner circle. The tangent
$\tau=(-\nu_y,\nu_x)$ is therefore counterclockwise on the outer boundary
and clockwise on the inner boundary. In both cases the region lies to the
left of positive traversal. The angular field
$(-y,x)/(x^2+y^2)$ has zero curl near the closed annulus. Its two boundary
circulations are $2\pi$ and $-2\pi$, summing to zero, exactly as Green's
formula requires. A disk containing the origin would violate the field's
domain hypothesis, so the same formula cannot be applied there.
\end{example}

\section{Divergence and surface Stokes}
\family{K04.12.03}\family{K04.12.04}
\begin{theorem}[Gauss divergence theorem]\label{d04:12:gauss}
For a bounded $C^1$ domain $D\subset\R^n$, $n\ge1$, and a $C^1$ vector
field on a neighborhood of its closure,
\[
 \int_D\operatorname{div}F\,dx=\int_{\partial D}F\cdot\nu\,dA.
\]
Here $dA$ is the positive $(n-1)$-dimensional Gram density; in dimension
one the right side is the sum of endpoint values times their outward
normal signs. The finite decomposable piecewise version follows under
\cref{d04:15:piecewise}.
\end{theorem}
\begin{proof}
Put $\Omega=dx_1\wedge\cdots\wedge dx_n$ and
$\beta_F=\sum_i(-1)^{i-1}F_i dx_{\widehat i}=\iota_F\Omega$.
The only surviving derivative of its $i$th term is the $i$th one; moving
$dx_i$ back into place cancels $(-1)^{i-1}$, hence
$d\beta_F=(\operatorname{div}F)\Omega$.
For $n\ge2$ and a positively oriented boundary basis $v_1,\ldots,v_{n-1}$,
decompose $F=(F\cdot\nu)\nu+F_T$. The tangential term in
$\Omega(F,v_1,\ldots,v_{n-1})$ vanishes, since all its vectors belong
to the $(n-1)$-dimensional tangent space. The remaining term is
\[
 (F\cdot\nu)\Omega(\nu,v_1,\ldots,v_{n-1})
 =(F\cdot\nu)\sqrt{\det(\ip{v_i}{v_j})},
\]
because $\nu$ is a unit normal and the outward-first basis is positive.
Thus the boundary restriction of $\beta_F$ is exactly the flux density
with its chosen orientation. Apply \cref{d04:15:stokes} and
\cref{d04:12:region-interface}. In $n=1$ the same formula is the ordinary
fundamental theorem with signed endpoints. Finite piecewise cancellation
is the theorem already proved in \cref{d04:15:piecewise}.
\end{proof}

\begin{theorem}[Classical Stokes for a surface]\label{d04:12:surface-stokes}
Let $S\subset\R^3$ be a compact oriented embedded $C^1$ surface with
boundary, and choose the continuous unit normal $N$ corresponding to its
orientation, so that $v_1\times v_2$ points along $N$ for a positive basis.
For $F\in C^1$ on a neighborhood of $S$,
\[
 \int_S(\operatorname{curl}F)\cdot N\,dA
 =\int_{\partial S}F\cdot\tau\,ds.
\]
The boundary tangent is $\tau=N\times\nu$, where $\nu$ is the outward
unit conormal, tangent to $S$ and perpendicular to $\partial S$.
The assertion also holds for a finite oriented union of regular surface
pieces with matching seams and the cancellation conditions of
\cref{d04:15:piecewise}; the remaining boundary is then piecewise regular.
\end{theorem}
\begin{proof}
The work form is $\alpha_F=F_1dx+F_2dy+F_3dz$. By
\cref{d04:13:dictionary}, $d\alpha_F=\beta_{\operatorname{curl}F}$.
In a positive surface chart,
\[
 \phi^*\beta_{\operatorname{curl}F}
 =(\operatorname{curl}F)(\phi)\cdot(\phi_u\times\phi_v)\,du\wedge dv
 =(\operatorname{curl}F)(\phi)\cdot N\,J_2\phi\,du\wedge dv.
\]
This is exactly the claimed scalar flux. The boundary restriction of
$\alpha_F$ is its work integral. Outward-first requires $(\nu,\tau)$
positive on $S$, hence $\nu\times\tau=N$; for perpendicular unit vectors
this is equivalent to $\tau=N\times\nu$. Applying the $C^1$ ambient
general Stokes theorem proves the formula. In the finite-piece case apply
the cell or manifold result to each piece and cancel the seam integrals.
\end{proof}

\begin{example}[Disk, hemisphere and box signs]\label{d04:12:signs}
On the horizontal disk oriented by $N=(0,0,1)$ the conormal is radial,
so $\tau=N\times\nu$ is counterclockwise as viewed from above. For the
upper hemisphere with outward normal, the conormal at its equator points
downward within the sphere, yielding the same counterclockwise tangent.
The lower hemisphere with outward normal yields the opposite equatorial
traversal, so their seam integrals cancel to give a closed sphere.
For the standard oriented box, the upper and lower $i$-faces have coefficients
$(-1)^{i-1}$ and $-(-1)^{i-1}$ in the remaining coordinate order. These
are precisely the cell signs of \cref{d04:15:boundary-square}; a box uses
the finite-cell theorem at corners, rather than a nonexistent smooth
boundary chart there.
\end{example}

\section{Integration by parts and Green identities}
\family{K04.12.07}\family{K04.12.08}
\begin{corollary}[Integration by parts]\label{d04:12:parts}
For a bounded $C^1$ domain, $f,g\in C^1$ near its closure and each coordinate
$i$,
\[
 \int_Df\partial_i g\,dx
 =\int_{\partial D}fg\nu_i\,dA-\int_Dg\partial_i f\,dx.
\]
For $u,v\in C^2$ near the closure, with
$\partial_\nu v=\nabla v\cdot\nu$,
\[
 \int_D(u\Delta v+\nabla u\cdot\nabla v)\,dx
 =\int_{\partial D}u\partial_\nu v\,dA,
\]
\[
 \int_D(u\Delta v-v\Delta u)\,dx
 =\int_{\partial D}(u\partial_\nu v-v\partial_\nu u)\,dA.
\]
If the relevant product or function is supported compactly in $D$, the
corresponding boundary term is zero. In particular
$\int_Du\Delta u=-\int_D|\nabla u|^2$ for $u\in C_c^2(D)$.
\end{corollary}
\begin{proof}
Apply Gauss to $fg e_i$, using
$\operatorname{div}(fg e_i)=f\partial_i g+g\partial_i f$.
Apply it to $u\nabla v$, whose divergence is
$\nabla u\cdot\nabla v+u\Delta v$, to obtain the first Green identity.
Subtract the identity with $u,v$ exchanged to obtain the second. Compact
support in $D$ means vanishing on a neighborhood of the boundary, so the
claimed boundary terms vanish. For $u=v$ in the first identity the final
energy equality follows. All differentiated vector fields are $C^1$ under
the stated orders; no PDE existence assertion is required.
\end{proof}

\section{Calibration and infinitesimal meaning}
\family{K04.12.09}
\begin{example}[Simple exact values]\label{d04:12:calibration}
For $F=x$ on a radius-$R$ sphere, \cref{d04:10:sphere} gives outward flux
$4\pi R^3$. Since $\operatorname{div}F=3$, Gauss gives the ball volume
$4\pi R^3/3$. On the box $[0,a]\times[0,b]\times[0,c]$ with $a,b,c>0$, the three lower
faces have zero flux and each upper face contributes $abc$, giving total
$3abc$, matching the integral of divergence. For
$F=(-y/2,x/2,0)$ one has $\operatorname{curl}F=(0,0,1)$.
On the radius-$R$ horizontal disk with upward normal, Stokes gives
circulation $\pi R^2$; direct substitution of
$(R\cos t,R\sin t,0)$ gives the same integrand $R^2/2$ on $[0,2\pi]$.
Reversing the surface normal reverses the traversal and both sides' signs.
\end{example}

\begin{corollary}[Local flux and circulation densities]\label{d04:12:local-meaning}
For a $C^1$ field near $x$, with balls contained in its domain,
\[
 \operatorname{div}F(x)=\lim_{r\downarrow0}
 \frac1{\vol B_r(x)}\int_{\partial B_r(x)}F\cdot\nu\,dA.
\]
For $F$ in $\R^3$ and a unit vector $N$, let $D_r$ be the oriented planar
disk of radius $r$ centered at $x$ with normal $N$. Then
\[
 \operatorname{curl}F(x)\cdot N
 =\lim_{r\downarrow0}\frac1{\pi r^2}
       \int_{\partial D_r}F\cdot\tau\,ds.
\]
\end{corollary}
\begin{proof}
Gauss and Stokes identify the respective quotients with averages of the
continuous functions $\operatorname{div}F$ and
$\operatorname{curl}F\cdot N$ over the shrinking ball or disk. The
difference from the value at $x$ is bounded by the supremum of the
corresponding difference on that ball or disk, which tends to zero.
This proves the geometric meanings of trace and curl under the chosen
normalizations, rather than assuming them as definitions.
\end{proof}

% END ORIGINAL INPUT: chapters/12.tex


% BEGIN ORIGINAL INPUT: chapters/13.tex
\chapter{Differential forms in Euclidean space}
\label{ch:d04-13}
\source{SPI Chapter 4, printed pp. 75--96, and Addenda p. 146; HUB Section 6.13, especially 6.13.9--6.13.12.}

The algebra in this chapter concerns covectors: the space in degree $k$ is
$\Alt^k(V;\R)$, naturally the exterior power of $V^*$, rather than of $V$.
All vector spaces are finite dimensional. We keep regularity in statements,
instead of giving the word ``differentiable'' the meaning $C^\infty$.

\section{Alternating algebra and pullback}
\family{K04.13.01}\family{K04.13.02}\family{K04.13.03}
\begin{definition}[Normalized covectors and forms]\label{d04:13:forms}
An alternating $k$-covector on $V$ is a multilinear map $V^k\to\R$ that
changes sign when two arguments are interchanged. In degree zero it is a
scalar. Relative to a basis $e_1,\ldots,e_n$ with dual $dx_1,\ldots,dx_n$,
write $dx_I(v_1,\ldots,v_k)=\det((v_b)_{i_a})_{a,b=1}^k$ for
$I=(i_1<\cdots<i_k)$. A $C^r$ differential $k$-form on an open
$U\subset\R^n$ is a field $x\mapsto\omega_x\in\Alt^k(\R^n;\R)$
whose coordinate coefficients are $C^r$. Degrees $k>n$ contain only zero;
degree zero consists of functions. We do not introduce negative degrees.
\end{definition}

\begin{proposition}[Basis and wedge normalization]\label{d04:13:wedge}
Every alternating covector is uniquely $\sum_I a_I dx_I$, with
$a_I=\omega(e_{i_1},\ldots,e_{i_k})$. For degrees $p,q$, define
\[
 (\alpha\wedge\beta)(v_1,\ldots,v_{p+q})
 =\sum_{\sigma\in\operatorname{Sh}(p,q)}\sgn(\sigma)
 \alpha(v_{\sigma(1)},\ldots,v_{\sigma(p)})
 \beta(v_{\sigma(p+1)},\ldots,v_{\sigma(p+q)}),
\]
where a shuffle preserves the order in each of its two blocks. This product
is bilinear, associative, and satisfies
$\alpha\wedge\beta=(-1)^{pq}\beta\wedge\alpha$; a scalar factor is ordinary
scalar multiplication. In particular $dx_{i_1}\wedge\cdots\wedge dx_{i_k}=dx_I$
for increasing indices. If $\operatorname{Alt}$ denotes the average over all
$(p+q)!$ permutations, then
$\alpha\wedge\beta=((p+q)!/(p!q!))\operatorname{Alt}(\alpha\otimes\beta)$.
\end{proposition}
\begin{proof}
Expand each argument in the basis. Terms with repeated basis vectors vanish,
and the remaining terms group by their increasing index set, giving the
determinant formula and uniqueness. In the shuffle sum each permutation of
all arguments occurs $p!q!$ times in the unnormalized alternation, proving
the last identity and alternation. For associativity, both bracketings sum
over ordered partitions of the argument positions into blocks of sizes
$p,q,r$, with the sign of the same permutation in each term. Exchanging the
first two blocks requires $pq$ transpositions, proving graded commutativity.
For one-covectors the shuffle expansion is the determinant expansion, proving
the coordinate normalization. These arguments also include a block of size
zero; the empty determinant and empty product equal one.
\end{proof}

\family{K04.13.04}
\begin{definition}[Pullback]\label{d04:13:pullback}
For $f:U\subset\R^m\to V\subset\R^n$ of class $C^1$, define
\[
 (f^*\omega)_x(v_1,\ldots,v_k)
 =\omega_{f(x)}(Df(x)v_1,\ldots,Df(x)v_k).
\]
In degree zero this is composition. A continuous form has a continuous
pullback. If $f$ is $C^{r+1}$ and $\omega$ is $C^r$, the pullback is $C^r$.
\end{definition}
\begin{proposition}[Pullback calculus]\label{d04:13:pullback-rules}
Pullback is linear, $f^*(a\omega)=(a\circ f)f^*\omega$,
$f^*(\alpha\wedge\beta)=f^*\alpha\wedge f^*\beta$, and
$(g\circ f)^*=f^*g^*$ wherever the composition is defined. Moreover
\[
 f^*dy_i=\sum_j\partial_j f_i\,dx_j,\qquad
 f^*(a\,dy_1\wedge\cdots\wedge dy_n)
 =(a\circ f)\det Df\,dx_1\wedge\cdots\wedge dx_n
\]
in the second formula when $m=n$.
\end{proposition}
\begin{proof}
Insert $Df(x)v_j$ into each defining sum. The shuffle sums agree term by
term, and the chain rule proves the composition order. Evaluation on $v$
gives $dy_i(Df(x)v)=\sum_j\partial_jf_i(x)v_j$. Evaluating the top form on
the standard basis gives the determinant, and the top alternating space is
one dimensional. This proves both coordinate formulas.
\end{proof}

\section{Exterior differentiation with finite regularity}
\family{K04.13.05}\family{K04.13.06}\family{K04.13.07}
\begin{definition}[Exterior derivative]\label{d04:13:exterior-derivative}
For $\omega=\sum_I a_I dx_I$ of class $C^1$, set
\[
 d\omega=\sum_{I,j}\partial_j a_I\,dx_j\wedge dx_I.
\]
It is continuous, and is $C^{r-1}$ for $C^r$ coefficients, $r\ge1$.
For $k=n$ it is zero. For a function this definition is precisely its total
derivative represented as a one-form.
\end{definition}

\begin{theorem}[Differential identities]\label{d04:13:d-rules}
For $C^1$ forms of degrees $p,q$,
\[
 d(\alpha\wedge\beta)=d\alpha\wedge\beta+(-1)^p\alpha\wedge d\beta.
\]
For a $C^2$ form, $d(d\omega)=0$. If $f$ is $C^2$ and $\omega$ is $C^1$,
then $f^*\omega$ is $C^1$ and
\[
 d(f^*\omega)=f^*(d\omega).
\]
\end{theorem}
\begin{proof}
Linearity reduces the product identity to $\alpha=a\,dx_I$,
$\beta=b\,dx_J$. Differentiate $ab$, and move $db$ past the $p$ coordinate
one-forms in $dx_I$; the sign is $(-1)^p$. The expansion of $d^2\omega$ is
$\sum_{I,i,j}\partial_i\partial_j a_I\,dx_i\wedge dx_j\wedge dx_I$.
Diagonal terms vanish, and the $(i,j)$ and $(j,i)$ terms cancel by equality
of continuous mixed partials from \cref{d04:03:symmetry}.
For naturality it suffices to treat $a\,dy_{i_1}\wedge\cdots\wedge dy_{i_k}$.
Its pullback is $(a\circ f)df_{i_1}\wedge\cdots\wedge df_{i_k}$.
The product identity and $d(df_i)=0$ leave only
$d(a\circ f)\wedge df_{i_1}\wedge\cdots\wedge df_{i_k}$.
The scalar chain rule identifies this with $f^*(da\wedge dy_I)$.
\end{proof}

\begin{remark}\label{d04:13:regularity-warning}
A $C^1$ map does not in general pull a $C^1$ form to a $C^1$ form:
the coefficients contain first derivatives of the map. Thus the displayed
pointwise naturality theorem is deliberately stated for $C^2$ maps. The
integrated Stokes formula for $C^1$ parametrizations will nevertheless be
proved by approximation in \cref{d04:15:cell-stokes-c1}; no nonexistent second
derivative will be used there.
\end{remark}

\family{K04.13.08}\family{K04.13.10}
\begin{definition}[Closed and exact]\label{d04:13:closed-exact}
A $C^1$ form is closed if $d\omega=0$. A positive-degree form is exact if
$\omega=d\eta$ for a $C^1$ form $\eta$. The coordinate $d^2$ proof above
requires $\eta\in C^2$. More generally a $C^1$ exact form is closed even
when its primitive is only $C^1$; \cref{d04:15:exact-closed-c1} proves this
using the independently proved cell formula and shrinking cubes.
For a $C^1$ zero-form, closed means locally constant: on a ball, the scalar
fundamental theorem along a segment gives $f(y)-f(x)=0$; overlapping balls
then give constancy on each connected component. No primitive of degree
minus one is asserted.
\end{definition}

\section{A primitive with an explicit homotopy}
\family{K04.13.09}
\begin{theorem}[Euclidean Poincare lemma]\label{d04:13:poincare}
Let $U\subset\R^n$ be open and star-shaped about $a\in U$.
For a continuous $k$-form, $k\ge1$, define the continuous $(k-1)$-form
\[
 (H\omega)_x(v_1,\ldots,v_{k-1})=
 \int_0^1 t^{k-1}\omega_{a+t(x-a)}(x-a,v_1,\ldots,v_{k-1})\,dt.
\]
If $\omega$ is $C^1$, then $H\omega$ is $C^1$ and
$dH\omega+H d\omega=\omega$. In particular every closed $C^1$ form of
positive degree on $U$ has the displayed $C^1$ primitive. If its coefficients
are $C^r$, the primitive is $C^r$, $1\le r\le\infty$.
\end{theorem}
\begin{proof}
Translate to $a=0$. For any $x_0$, the compact segment $[0,x_0]$ has positive
distance from the complement of $U$. Consequently all segments $[0,x]$
for $x$ in a sufficiently small ball about $x_0$ stay in one compact subset
of $U$. The coefficient integrands and their first $x$ derivatives are
continuous and bounded there. Their difference quotients converge uniformly,
by the one-variable integral mean-value estimate, so differentiation under
the proper integral is justified. Iteration proves the $C^r$ claim.

Here is the full homotopy computation. Let $\iota_X$ mean insertion of
$X(x)=x$ in the first argument. For a single coefficient term $b\,dx_I$ of
degree $k$, expansion of the product rule gives
\[
 d\iota_X(b\,dx_I)+\iota_X d(b\,dx_I)
 =\bigl(\sum_jx_j\partial_jb+kb\bigr)dx_I.
\]
Indeed $\iota_Xdx_I=\sum_{r=1}^k(-1)^{r-1}x_{i_r}dx_{I\setminus i_r}$.
Differentiating the $x_{i_r}$ produces $k b\,dx_I$; differentiation of $b$
produces $db\wedge\iota_Xdx_I$, which cancels the negative of that term
in $\iota_X(db\wedge dx_I)=(\sum_jx_j\partial_jb)dx_I-db\wedge\iota_Xdx_I$.
Apply the same expansion to the coefficients at $tx$, with the factors
$t^{k-1}$ in $H$. It gives
\[
 (dH\omega+Hd\omega)_x
 =\sum_I\int_0^1\bigl(kt^{k-1}b_I(tx)
       +t^k\sum_jx_j\partial_jb_I(tx)\bigr)\,dt\,dx_I
 =\sum_I[t^kb_I(tx)]_0^1dx_I.
\]
The lower endpoint is zero because $k>0$ and coefficients are bounded near
zero. This is $\omega$. This proof uses only proper one-variable parameter
integration and does not use Stokes. In degree zero the corresponding
identity is $Hdf=f-f(0)$, by the scalar fundamental theorem.
\end{proof}

\begin{example}\label{d04:13:homotopy-example}
For $\omega=z\,dx\wedge dz$ on $\R^3$, insertion and integration give
$H\omega=(-z^2dx+xz\,dz)/3$. Direct differentiation gives
$dH\omega=z\,dx\wedge dz$, since $d\omega=0$.
\end{example}

\section{Metric identifications and obstructions}
\family{K04.13.11}\family{K04.13.12}
\begin{proposition}[The oriented Euclidean dictionary]\label{d04:13:dictionary}
In standard oriented $\R^3$ set
\[
 \alpha_F=F_1dx+F_2dy+F_3dz,\quad
 \beta_F=F_1dy\wedge dz+F_2dz\wedge dx+F_3dx\wedge dy,
 \quad \Omega=dx\wedge dy\wedge dz.
\]
Then $df=\alpha_{\nabla f}$, $d\alpha_F=\beta_{\operatorname{curl}F}$,
and $d\beta_F=(\operatorname{div}F)\Omega$ whenever the first derivatives
are continuous. On vectors $v,w$, $\beta_F(v,w)=F\cdot(v\times w)$.
More generally for $k\ge1$, an inner product and orientation determine the unique top
covector equal to one on positive orthonormal bases. Its absolute value on
$(v_1,\ldots,v_k)$ is $\sqrt{\det(\ip{v_i}{v_j})}$ in a $k$-space.
\end{proposition}
\begin{proof}
Expand $d$ and collect the coefficients of $dy\wedge dz,dz\wedge dx,
dx\wedge dy$; they are respectively
$\partial_yF_3-\partial_zF_2,\partial_zF_1-\partial_xF_3,
\partial_xF_2-\partial_yF_1$. For $d\beta_F$, the three surviving terms
all have positive sign in $dx\wedge dy\wedge dz$. The cross-product
identity follows by the determinant expansion. For the final statement,
write the vectors as columns of $A$ in a positive orthonormal basis. The
two quantities are $|\det A|$ and $\sqrt{\det(A^TA)}$, hence equal.
Orientation reversal changes the top covector's sign, and leaves this
positive Gram density unchanged.
\end{proof}

\begin{example}[A closed form that has no global primitive]\label{d04:13:angular}
On $\R^2\setminus\{0\}$ let
$\theta=(-y\,dx+x\,dy)/(x^2+y^2)$. Differentiating shows
$\partial_x(x/(x^2+y^2))=\partial_y(-y/(x^2+y^2))$, hence $d\theta=0$.
Its integral on $(\cos t,\sin t)$, $0\le t\le2\pi$, is $2\pi$.
Were $\theta=df$, the scalar chain rule and fundamental theorem would give
zero on this closed path. Thus local exactness does not imply global exactness.
In \cref{d04:15:radial-obstruction} a closed two-form on punctured three-space
will supply the higher-degree counterpart. This finite collection of
obstructions requires no general cohomology theory.
\end{example}

% END ORIGINAL INPUT: chapters/13.tex


% BEGIN ORIGINAL INPUT: chapters/14.tex
\chapter{Embedded manifolds and finite localization}
\label{ch:d04-14}
\source{SPI Chapter 5, printed pp. 109--128; MUN Sections 23--24 and 37, checked printed pp. 199, 202--203, 303--305.}

Only subsets of Euclidean spaces are needed here. Coordinate changes,
tangent spaces, orientation and compact localization will be proved in that
setting. Abstract manifold topology, curvature and general cohomology belong
to subsequent subjects.

\section{Charts, graphs and tangent spaces}
\family{K04.14.01}\family{K04.14.02}\family{K04.14.04}
\begin{definition}[Embedded charts]\label{d04:14:embedded}
Let $1\le r\le\infty$ and $0\le k\le n$. Put
$H^k=\{u\in\R^k:u_k\ge0\}$ when $k>0$. An embedded $C^r$
$k$-manifold with boundary in $\R^n$ is a subset $M$ locally parametrized
by homeomorphisms $\phi:V\to M\cap O$, where $O\subset\R^n$ is open,
$V$ is open in $\R^k$ or relatively open in $H^k$, and $\phi$ has rank
$k$. A map on a half-space is $C^r$ if near every point it extends $C^r$
to an open Euclidean neighborhood, with the stated full rank. A boundary
chart is centered at $u_k=0$. Charts may always be restricted. In dimension
zero $M$ is discrete and has empty intrinsic boundary. A compact zero-manifold
is finite: its isolated point neighborhoods have a finite subcover.
\end{definition}

\begin{theorem}[Local graph and transition theorem]\label{d04:14:charts}
Every full-rank $C^r$ parametrization is, after shrinking near a parameter
point, a graph over a $k$-coordinate plane, with parameter domain a local
$C^r$ image of its original domain. Its inverse on that image is the
restriction of a $C^r$ ambient map. Two embedded charts have $C^r$
transition maps with invertible derivatives on their overlap. An interior
chart point can never be a boundary chart point; boundary charts map their
flat boundary onto the same subset $\partial M$. This subset is a $C^r$
$(k-1)$-manifold without boundary, and $M\setminus\partial M$ is a
$k$-manifold without boundary.
\end{theorem}
\begin{proof}
Choose $k$ independent rows of $D\phi(u)$ and let $P:\R^n\to\R^k$
be the associated coordinate projection. For a local extension of $\phi$,
$D(P\phi)$ is invertible. The inverse function theorem of
\cref{d04:04:inverse-function} gives a $C^r$ inverse $q$ on an open neighborhood.
Then $\phi q(v)$ has selected coordinates $v$; the remaining coordinates
are $C^r$ functions of $v$. On the patch, its inverse is $qP$. This proves
the graph and inverse assertions even at a half-space point.

Thus $\psi^{-1}\phi$ on an overlap is a restriction of an ambient $C^r$
composition. The identities with its inverse, differentiated in interior
parameter points, give inverse derivatives. At a flat boundary the same
identities follow by limits from the interior; equivalently first derivatives
are determined by the inward half-space, which spans $\R^k$. In particular
the transition derivative is invertible at every overlap point.

For completeness the half-space cannot disappear under a transition. A
one-sided differentiable curve starting at a parameter point has derivative
in all of $\R^k$ at an interior point, and in $H^k$ at a point of the flat
boundary (after translation). Conversely every vector in that cone is
realized by a short straight ray. A transition and its inverse therefore
carry these cones onto one another by their invertible derivatives. A
linear isomorphism cannot carry a proper half-space cone onto an entire
space. This proves the intrinsic distinction. Between boundary charts it
carries the largest vector subspace of the cone, $\{v_k=0\}$, onto itself.
Restriction to flat boundaries gives full-rank charts and $C^r$ transitions
in dimension $k-1$. Their domains are open in $\R^{k-1}$; hence they have
no boundary. Interior restrictions give the remaining assertion.
\end{proof}

\begin{proposition}[Regular level sets and graphs]\label{d04:14:level}
If $g:O\subset\R^n\to\R^{n-k}$ is $C^r$ and $Dg$ has rank $n-k$
on $g^{-1}(0)$, that level set is an embedded $C^r$ $k$-manifold without
boundary. Conversely every embedded chart near an interior point has a
local regular-level description. Graphs of $C^r$ maps are examples.
\end{proposition}
\begin{proof}
The implicit function theorem \cref{d04:04:implicit-function} solves $n-k$ variables
as a $C^r$ function of the other $k$. This is a graph chart, injective with
continuous inverse the coordinate projection. Conversely the preceding
graph construction writes the set locally as $(v,h(v))$; the function
$g(v,w)=w-h(v)$ has derivative $[-Dh\ I]$ and zero set exactly the graph.
\end{proof}

\family{K04.14.03}
\begin{proposition}[Tangent spaces]\label{d04:14:tangent}
For $x=\phi(u)$ define $T_xM=\operatorname{im}D\phi(u)$. This definition
is chart independent. At an interior point it is precisely the set of
velocities of two-sided $C^1$ curves in $M$ through $x$; at a boundary
point it is the span of the one-sided velocities. For a regular level set,
$T_xM=\ker Dg(x)$, and at a boundary point $T_x\partial M$ is the
image of $\{v_k=0\}$ under a boundary chart derivative.
\end{proposition}
\begin{proof}
The transition chain rule and its invertible derivative show equality of
the images. A curve in a chart has derivative $D\phi(u)c'(0)$, and every
parameter velocity is realized by a straight line at an interior point.
At the boundary the straight rays with inward velocities span the parameter
space, giving the span statement. Differentiating $g\phi=0$ gives inclusion
in the kernel, and both spaces have dimension $k$. The last assertion is
the flat-boundary chart description.
\end{proof}

\begin{remark}[Immersions and embeddings]\label{d04:14:immersion}
A $C^r$ immersion is a full-rank parametrization and is locally an embedded
patch by \cref{d04:14:charts}. It need not be globally injective: the circle
traversed twice has rank one and multiplicity two. An embedding additionally
is a homeomorphism onto its image. An injective immersion need not be an
embedding on a noncompact domain: $t\mapsto(\cos t,\sin t)$ on $[0,2\pi)$
is injective and regular, but its inverse is discontinuous at $(1,0)$.
On a compact domain an injective continuous map into Euclidean space has
continuous inverse, because images of closed subsets are compact and closed.
The sphere has empty intrinsic boundary although its ambient topological
boundary is the entire sphere.
\end{remark}

\section{Orientation and boundary signs}
\family{K04.14.05}\family{K04.14.06}\family{K04.14.09}
\begin{definition}[Orientation]\label{d04:14:orientation}
For $k>0$ an orientation of a $k$-space is one of its two classes of ordered
bases, where two bases are equivalent if their change-of-basis determinant
is positive. An orientation of $M$ is a continuous choice of these classes.
Precisely, the ordered derivative columns of any chart $\phi$ agree or
disagree with the choice by a locally constant sign $\sigma_\phi\in\{1,-1\}$.
If $\psi=\phi h$, the chain rule gives
$\sigma_\psi=(\sigma_\phi\circ h)\sgn\det Dh$.
A positive chart means one with sign $+1$. Interior charts can always be
made positive by a coordinate reflection. Boundary charts use the fixed
half-space $u_k\ge0$; in dimension one a boundary chart may necessarily
have sign $-1$, since a reflection would change the half-space.
We therefore retain chart signs throughout integration and do not assume
an all-positive boundary atlas. In dimension zero an orientation is an explicitly chosen
sign $+1$ or $-1$ at each point, as needed by boundary integration.
\end{definition}

\begin{proposition}[Outward-first boundary orientation]\label{d04:14:boundary-orientation}
Let $M$ be oriented. At $x\in\partial M$ there is a unique unit vector
$\nu\in T_xM$ perpendicular to $T_x\partial M$ and pointing out of $M$.
For $k\ge2$, declare $(v_1,\ldots,v_{k-1})$ positive on $\partial M$ precisely when
$(\nu,v_1,\ldots,v_{k-1})$ is positive on $M$. This is chart independent
and continuous. For $k=1$, give a boundary point sign $+1$ if its outward
vector agrees with the chosen tangent-line orientation, and sign $-1$
otherwise; this is a signed-point convention, not an empty-basis class.
For $H^k$ in its standard orientation the boundary
$u_k=0$ has $(-1)^k$ times the standard orientation of $\R^{k-1}$.
For an oriented interval the terminal endpoint is $+1$ and initial endpoint
is $-1$.
\end{proposition}
\begin{proof}
The quotient $T_xM/T_x\partial M$ is a line. The inward tangent cone,
whose invariance was proved in \cref{d04:14:charts}, singles out its positive
half. Orthogonal projection of a transverse inward chart vector onto the
orthogonal complement of $T_x\partial M$ gives a nonzero inward vector;
normalize and negate to obtain $\nu$. The construction proves uniqueness,
independence and continuity. Replacing $\nu$ by any outward vector adds a
tangential component and multiplies by a positive scalar, so the resulting
orientation is unchanged. In the standard half-space $\nu=-e_k$; moving
$-e_k$ from first to last position has sign $(-1)^k$. The case $k=1$
states exactly the endpoint signs, with orientation on a zero-space
interpreted as the stated integration sign.
\end{proof}

\family{K04.14.08}
\begin{proposition}[Orientations and densities]\label{d04:14:metric-density}
An oriented Euclidean $k$-manifold, $k\ge1$, has the metric top form $dV_M$ with
\[
 \phi^*dV_M=\sqrt{\det(D\phi^TD\phi)}\,du_1\wedge\cdots\wedge du_k
\]
in a positive chart; in an arbitrary chart the right side is multiplied
by $\sigma_\phi$. Its unsigned density exists without orientability.
For $k=1,2,n$ the displayed factor is respectively speed, the magnitude of
the cross product in $\R^3$, and $|\det D\phi|$ in $\R^n$.
For $k=0$, the unsigned density is $1$, and the signed metric zero-form
has value equal to the chosen orientation sign at the point.
\end{proposition}
\begin{proof}
For $k\ge1$, at each tangent space choose a positive orthonormal basis and use its
normalized top covector. Orthogonal positive basis changes have determinant
one, so it is independent of this choice. The Gram identity in
\cref{d04:13:dictionary} gives the pullback. Under a transition $h$ the
Gram matrix becomes $Dh^T G Dh$ and its positive square-root determinant
becomes $|\det Dh|\sqrt{\det G}$. Thus the absolute density glues even when
some transitions reverse orientation. The special cases follow by direct
Gram determinants; for two vectors this is
$|v|^2|w|^2-(v\cdot w)^2=|v\times w|^2$.
In degree zero the empty Gram determinant is one, and the stated signed
value is the definition consistent with a signed point's integral.
\end{proof}

\section{Bumps and finite compact localization}
\family{K04.14.07}\family{K04.14.10}
\begin{lemma}[Explicit Euclidean cutoffs]\label{d04:14:bumps}
For compact $K\subset O\subset\R^n$ open there is a $C^\infty$ function
$\chi:\R^n\to[0,1]$ equal to one near $K$, with compact support in $O$.
\end{lemma}
\begin{proof}
Put $b(t)=e^{-1/t}$ for $t>0$ and $b(t)=0$ for $t\le0$.
Every right derivative is a polynomial in $1/t$ times $e^{-1/t}$ and tends
to zero at zero (the exponential series dominates every fixed power).
Induction using the scalar derivative definition proves $b\in C^\infty$.
Thus
$s(t)=b(t)/(b(t)+b(1-t))$ is smooth, equals zero for $t\le0$, equals one
for $t\ge1$, and lies in $[0,1]$.
For $0<r<R$, the function
$s((R^2-|x-p|^2)/(R^2-r^2))$ equals one on $\cl B(p,r)$ and vanishes
outside $B(p,R)$. Choose finitely many such smaller balls covering $K$,
with larger closed balls in $O$, and associated functions $\rho_i$.
The function $\chi=1-\prod_i(1-\rho_i)$ has all the required properties.
Only scalar differentiation, composition and compactness were used.
\end{proof}

\begin{theorem}[Finite partition near a compact set]\label{d04:14:partition}
Let $K\subset M$ be compact and let relative open sets $U_\alpha$ cover it.
There are finitely many nonnegative functions $\lambda_i$ on $M$, each the
restriction of an ambient smooth function near $M$, with compact support
in some $U_{\alpha_i}$, such that $\sum_i\lambda_i=1$ near $K$ in $M$.
Their sum can be chosen between zero and one everywhere. The supports can
be placed in arbitrarily small chart neighborhoods. This assertion includes
boundary points and uses no integration or change of variables.
\end{theorem}
\begin{proof}
First shrink the covering neighborhoods. The graph construction makes $M$
locally closed in an ambient neighborhood, including at a boundary point:
the graph over the local closed half-space is closed there. Hence for each
$p\in K$ one can choose ambient balls $B(p,r_p)$ and $B(p,R_p)$ whose
closed larger ball meets $M$ in a compact subset of one chosen $U_\alpha$
and one chosen chart. Take finitely many smaller balls covering $K$ and
smooth cutoffs $\rho_i$ supported in their larger balls, equal to one on
the smaller balls. Set $S=\sum_i\rho_i$. It is positive near $K$.
Choose an ambient cutoff $\chi$ equal to one near $K$, with compact support
in the open set $\{S>0\}$, using \cref{d04:14:bumps}.
Define $\lambda_i=\chi\rho_i/S$ on that open set and zero off it.
These extensions are smooth because $\supp\chi$ is a compact subset of
$\{S>0\}$. Restriction to $M$ gives compact subordinate supports, and
$\sum_i\lambda_i=\chi$ is one near $K$. This also proves the claimed
support control. No appeal to substitution, a manifold integral, or Stokes
has occurred.
\end{proof}

\begin{proposition}[An approximation tool]\label{d04:14:mollify}
If $f$ is $C^1$ on an open neighborhood of a compact Euclidean rectangle
$Q$, there are $C^\infty$ maps $f_j$ on a neighborhood of $Q$ converging
to $f$ together with their first derivatives uniformly on $Q$.
\end{proposition}
\begin{proof}
Use \cref{d04:14:bumps} to replace $f$, component by component, by a $C^1$
map $g$ of compact support agreeing with it near $Q$; extension by zero is
$C^1$ because the cutoff vanishes near the edge of its defining open set.
Let $\eta(z)=c\,b(1-|z|^2)$, with $c$ chosen so that
$\int\eta(z)\,dz=1$. This ordinary Riemann integral exists since $\eta$
is continuous with compact support; it is positive since $\eta>0$ on a
small box about zero. Set
$g_\eps(x)=\int\eta(z)g(x-\eps z)\,dz$ on a fixed box containing the unit
ball. Differentiating using the continuous first derivatives on a compact
set gives $Dg_\eps(x)=\int\eta(z)Dg(x-\eps z)\,dz$.
Uniform continuity yields
$\sup_Q(|g_\eps-g|+|Dg_\eps-Dg|)\to0$.
To see smoothness rather than merely $C^1$, change the affine integration
variable to write
$g_\eps(x)=\int\eps^{-n}\eta((x-y)/\eps)g(y)\,dy$.
This affine substitution follows already by scaling and translating
rectangular Riemann sums (or \cref{ch:d04-07}); it does not use any manifold
result. On a bounded neighborhood of $Q$ a fixed box contains every support
in $y$. All derivatives of the kernel are continuous and bounded there;
differentiation under the proper integral to every order proves smoothness.
Take any sequence $\eps_j\downarrow0$.
\end{proof}

\begin{remark}\label{d04:14:scope}
The construction supplies exactly the compact localization needed below.
For a noncompact manifold and compactly supported form, only a finite
partition near its support is required. A general locally finite partition
on an arbitrary abstract manifold is not an additional assumption here.
\end{remark}

% END ORIGINAL INPUT: chapters/14.tex


% BEGIN ORIGINAL INPUT: chapters/15.tex
\chapter{Integration of forms and generalized Stokes}
\label{ch:d04-15}
\source{SPI Chapter 4, printed pp. 97--104, and Chapter 5, pp. 122--130; MUN Section 37, checked printed pp. 303--306; HUB Section 6.10 and Appendix A.23 (the latter's trimming route is not needed for the finite-chart proof here).}

\section{Cells and finite chains}
\family{K04.15.01}\family{K04.15.02}\family{K04.15.05}
\begin{definition}[Cells and their integrals]\label{d04:15:cells}
A $C^r$ parametrized $k$-cell in an open $O\subset\R^n$ is a map
$c:[0,1]^k\to O$ having a $C^r$ extension to a neighborhood of the cube.
Neither rank nor injectivity is required. For a continuous $k$-form, write
$c^*\omega=a(u)\,du_1\wedge\cdots\wedge du_k$ and define
$\int_c\omega=\int_{[0,1]^k}a(u)\,du$. This is a usual proper Riemann
integral. A zero-cell is a point and $\int_c f=f(c)$.
A finite $k$-chain is a finite formal integer linear combination of such
cells; its integral is the corresponding linear combination. No quotient
by reparametrizations or homology is implicit.

Let $F_i^\eps$ insert $\eps\in\{0,1\}$ as coordinate $i$, in the order
of the remaining coordinates. Set
\[
 \partial c=\sum_{i=1}^k(-1)^{i-1}(cF_i^1-cF_i^0),
\]
and extend linearly to chains. The boundary of a zero-chain is zero.
\end{definition}

\begin{proposition}[Boundary cancellation]\label{d04:15:boundary-square}
For every finite chain, $\partial^2c=0$. The face signs are exactly the
outward-first signs for the standard oriented cube.
\end{proposition}
\begin{proof}
Fix two distinct coordinates $i<j$ and fixed values on both. The resulting
codimension-two face occurs by removing $i$ and then position $j-1$, or
by removing $j$ and then position $i$. The exponents of the two orientation
factors differ by one: $(i-1)+(j-2)$ and $(j-1)+(i-1)$.
The factors contributed by choosing lower or upper faces agree in both
orders, so the two occurrences cancel. Every term is paired this way.
On an upper $i$-face the outward vector is $e_i$; placing it first before
the remaining ordered coordinate basis has sign $(-1)^{i-1}$.
The lower face uses $-e_i$ and therefore the negative sign.
\end{proof}

\begin{theorem}[Stokes on a $C^2$ cell]\label{d04:15:cell-stokes}
For a $C^2$ $k$-cell, $k\ge1$, and a $C^1$ $(k-1)$-form defined on a
neighborhood of its image,
\[
 \int_c d\omega=\int_{\partial c}\omega.
\]
The same holds for finite chains by linearity.
\end{theorem}
\begin{proof}
First let $c$ be the identity cube. Write
$\omega=\sum_i b_i\,du_1\wedge\cdots\wedge\widehat{du_i}
\wedge\cdots\wedge du_k$. Then
\[
 d\omega=\sum_i(-1)^{i-1}\partial_i b_i\,du_1\wedge\cdots\wedge du_k.
\]
All coefficients and their derivatives are continuous. The continuous-slice
Fubini theorem \cref{d04:06:fubini}, followed by the scalar fundamental
theorem in coordinate $i$, integrates this to
\[
 \sum_i(-1)^{i-1}\int_{[0,1]^{k-1}}
 [b_i(u_1,\ldots,1,\ldots,u_k)-b_i(u_1,\ldots,0,\ldots,u_k)]\,du_{\widehat i}.
\]
Pullback to the $i$-faces kills every summand except $b_i du_{\widehat i}$,
so this is precisely the boundary formula with its signs.
For arbitrary $c$, use \cref{d04:13:d-rules} to replace $c^*d\omega$
by $d(c^*\omega)$ and functoriality on every face. When $k=1$ the
calculation is the scalar endpoint fundamental theorem; it includes the
zero-dimensional boundary convention.
\end{proof}

\begin{corollary}[Closedness with a $C^1$ primitive]\label{d04:15:exact-closed-c1}
If a $C^1$ $k$-form on an open subset of $\R^n$, $k\ge1$, satisfies
$\omega=d\eta$ for a $C^1$ $(k-1)$-form, then $d\omega=0$.
\end{corollary}
\begin{proof}
For $k\ge n$ the exterior derivative is zero by degree. Otherwise take
any sufficiently small affine $(k+1)$-cube $c$ in the open set. The cell
formula, applied first to $\omega$ and then to $\eta$ on its $k$-faces,
gives
\[
 \int_c d\omega=\int_{\partial c}\omega
 =\int_{\partial c}d\eta=\int_{\partial^2c}\eta=0.
\]
Every integrand has the required regularity and every cell is affine, so
\cref{d04:15:cell-stokes} and \cref{d04:15:boundary-square} apply without
assuming the present conclusion. For a fixed increasing $(k+1)$-index
set, choose cubes parallel to those coordinate directions, based at any
given point and with side length $h\downarrow0$. Divide the zero integral
by $h^{k+1}$. Continuity of $d\omega$ makes this average converge to its
corresponding coefficient at the point, which is therefore zero. All
coefficients vanish.
\end{proof}

\begin{theorem}[Stokes on a $C^1$ cell]\label{d04:15:cell-stokes-c1}
The conclusion of \cref{d04:15:cell-stokes} holds for a $C^1$ cell and a
$C^1$ ambient form. In particular no second derivative of the cell is needed.
\end{theorem}
\begin{proof}
Extend the cell $C^1$ to an open neighborhood of the cube and apply
\cref{d04:14:mollify} componentwise. The resulting smooth $c_j$ converge
uniformly with first derivatives on the cube. Their images lie in the
form's domain for large $j$, since the original image is compact there.
Every coefficient of $c_j^*d\omega$ is a finite sum of a continuous
coefficient of $d\omega$ evaluated at $c_j$ times a minor of $Dc_j$.
Uniform continuity on a compact neighborhood and convergence of the minors
give uniform convergence to $c^*d\omega$. The same holds for $c_j^*\omega$
on each face: the tangent derivative columns in each face's pullback are
restrictions of $Dc_j$, so the pullback coefficients converge uniformly.
No derivatives of those pullback coefficients are asserted to converge.
Hence the integrals converge, since the absolute difference is
bounded by cube volume times the uniform difference. Apply the $C^2$ cell
theorem to $c_j$ and take limits. This proof does not assert that
$d(c^*\omega)$ exists when $c$ is merely $C^1$.
\end{proof}

\section{Chart independence and manifold integration}
\family{K04.15.03}\family{K04.15.04}\family{K04.15.10}
\begin{lemma}[Directed substitution]\label{d04:15:directed-substitution}
If $h:V\to W$ is a $C^1$ diffeomorphism between connected open subsets
of $\R^k$, and $\eta$ is a continuous compactly supported top form on $W$,
then
\[
 \int_Vh^*\eta=\sgn(\det Dh)\int_W\eta.
\]
It also holds for relative open half-space domains and compactly supported
coefficients continuous up to their flat boundary, when $h$ is a chart
transition; signs are taken separately on overlap components.
\end{lemma}
\begin{proof}
Write $\eta=a(w)dw_1\wedge\cdots\wedge dw_k$. Its pullback coefficient
is $(a\circ h)\det Dh$, by \cref{d04:13:pullback-rules}.
The nonzero continuous determinant has constant sign on a connected domain.
Apply \cref{d04:07:continuous-substitution} to its absolute value.
For a half-space transition, \cref{d04:14:charts} supplies, near each point
of the compact coefficient support, an extension with invertible derivative;
the inverse theorem gives a local Euclidean diffeomorphism. It maps the
local half-space onto the local half-space because the transition and its
inverse do. A finite compact partition in these neighborhoods reduces the
assertion to each such extension. Zero extension of a compactly supported
coefficient across the flat boundary is Riemann integrable: it is continuous
off a compact portion of that hyperplane, which is content zero by covering
it with boxes of arbitrarily small thickness. The Riemann substitution
\cref{d04:07:substitution} therefore applies to the extension, with indicator
of the half-space. This proves the relative version, including existence of
both integrals. The same localization deals with disconnected overlaps.
\end{proof}

\begin{definition}[Integration on an embedded manifold]\label{d04:15:manifold-integral}
Let $M$ be an oriented $C^1$ embedded $k$-manifold with boundary, $k\ge1$, and let
$\omega$ be a continuous alternating top-covector field on its tangent
spaces, with compact support. Choose a finite partition $\lambda_i$ equal
to one near that support as in \cref{d04:14:partition}, subordinate to
charts $\phi_i:V_i\to U_i$ with constant orientation signs $\sigma_i$, and define
\[
 \int_M\omega=\sum_i\sigma_i\int_{V_i}\phi_i^*(\lambda_i\omega).
\]
The coefficients have compact support relative to $V_i$ and are extended
by zero inside a containing rectangle, restricting to $H^k$ when necessary.
The sign is $+1$ precisely when the chart's ordered derivative basis is
positive; neighborhoods can be restricted to make the sign constant.
For a compact zero-manifold this is the sum of the values times their
chosen orientation signs. For an unsigned density replace the oriented
coefficient by its density coefficient and use absolute determinants.
\end{definition}

\begin{theorem}[Independence of charts and weights]\label{d04:15:independence}
The preceding definition is independent of every chart and partition
choice. It is linear, additive under finite localization, reverses sign
when the orientation reverses, and agrees with a single-chart integral
when the support lies in that chart. Unsigned density integration is
orientation independent, even on a nonorientable manifold.
\end{theorem}
\begin{proof}
For a form supported compactly inside an overlap, the two coordinate
expressions are related by pullback through the transition, by the chain
rule. For $\psi=\phi h$ their signs satisfy
$\sigma_\psi=(\sigma_\phi\circ h)\sgn\det Dh$, by
\cref{d04:14:orientation}. Multiplying directed substitution by these signs
makes the two signed chart integrals equal. If the
overlap is not a single coordinate neighborhood, a finite compact partition
there reduces to the local statement in that lemma; each term is supported
in one common neighborhood. Now compare partitions $\lambda_i$ and $\mu_j$.
Both give, by finite linearity, the double sum of the integrals of
$\lambda_i\mu_j\omega$, because each sum of weights is one near the form's
support. Each double-sum term is chart independent by the overlap result.
Interchanging the finite sums proves equality. The same proof with absolute
Jacobians treats densities, with no orientation choice. Linearity and the
sign assertion follow directly. In dimension zero the assertions are
finite-sum identities. Thus there is neither an unproved overlap rule nor
an unnoticed extra counting of overlapping charts.
\end{proof}

\begin{proposition}[Intrinsic forms at finite regularity]\label{d04:15:intrinsic-d}
On a $C^2$ embedded manifold a $C^1$ intrinsic $p$-form is a tangent
covector field whose pullbacks in charts are $C^1$. There is a unique
continuous $(p+1)$-form $d_M\omega$ satisfying
$\phi^*d_M\omega=d(\phi^*\omega)$. For a $C^1$ ambient form its tangent
restriction has exterior derivative equal to the tangent restriction of
the ambient exterior derivative.
\end{proposition}
\begin{proof}
Transitions are $C^2$. The pointwise naturality theorem
\cref{d04:13:d-rules} shows that the local definitions agree on overlaps.
Chart derivatives are isomorphisms onto tangent spaces, which gives
existence and uniqueness of the covector field. The same naturality with
the chart map proves the ambient restriction assertion. The product rule
for $d_M$ follows in each chart from the Euclidean product rule.
\end{proof}

\section{Generalized Stokes, including the local proof}
\family{K04.15.06}\family{K04.15.07}
\begin{theorem}[Generalized Stokes]\label{d04:15:stokes}
Let $M\subset\R^n$ be an oriented $C^2$ embedded $k$-manifold with
boundary, $k\ge1$. Let $\omega$ be a compactly supported $C^1$ intrinsic
$(k-1)$-form. Give $\partial M$ the outward-first orientation. Then
\[
 \int_Md_M\omega=\int_{\partial M}\omega.
\]
In particular this holds for every such form on compact $M$.
There is also the following $C^1$ ambient version: $M$ may be $C^1$ if
$\omega$ is a $C^1$ form on an ambient neighborhood and
$M\cap\supp\omega$ is compact (in particular if $M$ is compact). In this
version the left side means the tangent restriction of ambient $d\omega$.
\end{theorem}
\begin{proof}
We first refine the chart cover before choosing any weights. Around each
interior parameter point choose an open rectangle whose closure lies in
the chart domain, and a smaller rectangle with closure inside it. Around
a boundary point choose a relatively open flat-bottom box in $H^k$, with
closure in the original chart domain, and a smaller such box whose closure
avoids all the larger box's side and upper faces. Their smaller images
cover the relevant compact set. Apply \cref{d04:14:partition} to these
smaller images; thus every localized support lies inside a larger parameter
rectangle and strictly away from all its artificial faces. This explicit
refinement is needed because an arbitrary compact subset of an irregular
chart domain need not fit in one rectangle contained in that domain.

We now prove the local calculation, including its artificial edges.
In an interior chart of sign $\sigma_\phi$, take a compactly supported coordinate
$(k-1)$-form $\eta$ on the open parameter domain. Extend its coefficients
by zero and choose a large rectangle containing their support in its
interior. The coordinate cube proof of \cref{d04:15:cell-stokes}, rescaled
affinely, gives $\int d\eta=0$, since every outer face has zero integrand.
In a boundary chart of sign $\sigma_\phi$ with parameters $u_k\ge0$, take instead a
rectangle $[-A,A]^{k-1}\times[0,A]$ containing the support away from its
side and upper faces. Coefficients extend by zero across the artificial
edges; they have $C^1$ extensions across the genuine flat boundary by the
definition of boundary regularity. The same calculation leaves precisely
the lower face, with coordinate factor $(-1)^k$. The bulk chart integral
has factor $\sigma_\phi$. By \cref{d04:14:boundary-orientation}, the
induced orientation of the flat-boundary parametrization has factor
$\sigma_\phi(-1)^k$. Thus the same signed equality holds on $M$.
This includes a necessarily negative right-endpoint chart when $k=1$.
Thus the local equality is proved, including the $k=1$
endpoint calculation.

Choose the finite partition $\lambda_i$ of
\cref{d04:14:partition} near $K=\supp\omega$, with supports inside such
small interior or boundary charts. It suffices to apply the local calculation
to $\lambda_i\omega$. Indeed $\sum_i\lambda_i\omega=\omega$, and the
product rule gives
\[
 \sum_i d_M(\lambda_i\omega)
 =\sum_i d_M\lambda_i\wedge\omega+\sum_i\lambda_i d_M\omega
 =d_M\omega.
\]
The first term is zero where $\omega$ can be nonzero, since the sum of
weights is one on a neighborhood of $K$; outside $K$ the form and its
derivative vanish locally. The restriction of the same partition sums to
one near the boundary support. Finite linearity and
\cref{d04:15:independence} now give the formula globally.

For the ambient version choose $K=M\cap\supp\omega$ and ambient weights
as in \cref{d04:14:partition}. On $M$ near $K$ their sum is one, so its
derivative on tangent vectors is zero, including at boundary points by
continuity from the interior. Outside $K$ the ambient form and derivative
are zero. Thus the same finite product identity holds after tangential
restriction. To justify the local calculation for a $C^1$ chart $\phi$,
choose a closed parameter rectangle containing the support with the same
edge conditions. The ambient form $\lambda_i\omega$ need not have compact
support off $M$, which causes no difficulty: its restriction to the chart
vanishes on all artificial faces. Apply
\cref{d04:15:cell-stokes-c1} to the $C^1$ cell obtained by restricting
$\phi$ to that rectangle. It gives the required local equality directly,
without differentiating $\phi^*\omega$. The surviving boundary-chart sign is
the same $\sigma_\phi(-1)^k$. This completes the $C^1$ version.
\end{proof}

\begin{remark}[The classical interfaces]\label{d04:15:classical-bridge}
For $k=2$ in the plane use $Pdx+Qdy$, whose derivative is
$(\partial_xQ-\partial_yP)dx\wedge dy$. For $k=n$ use
$\sum_i(-1)^{i-1}F_i dx_1\wedge\cdots\wedge\widehat{dx_i}
\wedge\cdots\wedge dx_n$, whose derivative is
$(\sum_i\partial_iF_i)dx_1\wedge\cdots\wedge dx_n$.
For $k=2$ in $\R^3$ use $\alpha_F$ and $d\alpha_F=\beta_{\operatorname{curl}F}$.
These are the actual coefficient identifications, with the induced signs;
\cref{ch:d04-12} proves the corresponding normal and tangent identifications
and states the resulting Green, Gauss and classical Stokes formulas.
\end{remark}

\section{Finite pieces and consequences}
\family{K04.15.08}\family{K04.15.09}
\begin{theorem}[Finite piecewise Stokes]\label{d04:15:piecewise}
Let $C=\sum_{j=1}^N a_jc_j$ be a finite chain of $C^1$ $k$-cells, or a
finite sum of oriented compact $C^1$ embedded $k$-manifolds with boundary.
Suppose the internal boundary pieces can be paired with identical
parametrizations up to an orientation-reversing $C^1$ diffeomorphism,
so their directed integrals cancel. Suppose the unpaired boundary pieces
give the prescribed outer boundary, with matching orientations. Then for
an ambient $C^1$ $(k-1)$-form,
\[
 \sum_j a_j\int_{c_j}d\omega=\int_{\text{outer boundary}}\omega.
\]
If pieces represent one geometric region or surface, require their interiors
to be disjoint and their remaining overlaps to lie in the displayed faces
or finitely many lower-dimensional $C^1$ pieces. Under those hypotheses the
sum has multiplicity one and is its usual piecewise integral.
\end{theorem}
\begin{proof}
Apply the cell theorem or ambient manifold theorem to each summand and add.
Directed substitution makes each paired internal contribution equal and
opposite. Faces intersecting along edges do not create extra integrals:
in face coordinates those edge sets have content zero. More generally a
compact Lipschitz image of a bounded $d$-rectangle in $\R^k$, $d<k$, has
content zero: divide its domain into $O(\delta^{-d})$ boxes of side
$\delta$; the images fit into boxes of side at most a constant times
$\delta$, with total $k$-volume $O(\delta^{k-d})\to0$.
A $C^1$ map on a compact set is locally Lipschitz and a finite subdivision
gives that estimate. Thus the stipulated lower-dimensional overlaps have
zero contribution to the relevant continuous integrals. The disjointness
condition precludes hidden multiple counting.
\end{proof}

\begin{remark}\label{d04:15:corners-scope}
A box and any explicitly assembled finite collection of regular cells fall
under \cref{d04:15:piecewise}, although a box is not a manifold with smooth
boundary at its corners. The theorem also covers surfaces made of finitely
many regular pieces with cancelling seams. It makes no assertion for an
arbitrary set merely described as having a ``piecewise smooth'' or rough
boundary without the stated decomposition and matching conditions.
\end{remark}

\begin{corollary}[Exact integrals on cycles]\label{d04:15:exact-cycles}
An exact form integrates to zero on every finite cell cycle on which its
primitive has the regularity of the cell Stokes theorem. A compact oriented
boundaryless manifold has $\int_Md\eta=0$ under the intrinsic or ambient
hypotheses of \cref{d04:15:stokes}.
\end{corollary}
\begin{proof}
Stokes gives the integral of $\eta$ over the zero boundary in each case.
\end{proof}

\begin{example}[A higher-degree obstruction]\label{d04:15:radial-obstruction}
On $\R^3\setminus\{0\}$ take $F(x)=x/|x|^3$ and $\omega=\beta_F$.
Direct differentiation gives
$\operatorname{div}F=3|x|^{-3}-3|x|^{-5}\sum_i x_i^2=0$; hence
$d\omega=0$. On the unit sphere with outward orientation,
$F\cdot n=1$, so $\int_{S^2}\omega=4\pi$, using the sphere-area
calculation in \cref{d04:10:sphere}. If $\omega=d\eta$ for a $C^2$
one-form on punctured space, the boundaryless Stokes theorem on the sphere
would give zero, a contradiction. Even a $C^1$ ambient primitive is
excluded by the ambient version of Stokes. Thus a divergence-free field
need not have a global vector potential. The obstruction is a flux through
a closed surface, distinct from the planar circulation obstruction.
\end{example}

% END ORIGINAL INPUT: chapters/15.tex

\backmatter

% BEGIN ORIGINAL INPUT: source-register.tex
\chapter*{Source register and editorial corrections}

\addcontentsline{toc}{chapter}{Source register and editorial corrections}

The abbreviations below identify the exact local comparison editions. Original

book bytes are preserved; their presence is not a redistribution licence.

The prose and proofs in this manuscript are newly written. References to a

chapter or appendix are structural comparison pointers; a claim of checking

a specific passage is recorded separately in the author and review records.

\section*{HUB: Vector Calculus, Linear Algebra, and Differential Forms: A Unified Approach}

John H. Hubbard; Barbara Burke Hubbard.

Fifth edition, 2015; ISBN 9780971576681. The retained Library of Congress text describing the fourth edition is original frontmatter, not evidence that this is the fourth edition.

Local PDF: 837 physical pages. Identifier: A04.

\par\noindent\texttt{SHA-256: 303eaa45ff2202c521b1bf2c79596042\allowbreak 3756e99baed53eb42c13207e12cf3b97}.

\section*{MUN: Analysis on Manifolds}

James R. Munkres.

1991 edition / reprint

Local PDF: 380 physical pages. Identifier: D04-S01.

\par\noindent\texttt{SHA-256: 22a29728b0ccf2a5afe5eb9af494b801\allowbreak 00a9f7fcfa067f844f30385a4a02454e}.

\section*{FLE: Functions of Several Variables}

Wendell Fleming.

Second edition, 1977

Local PDF: 420 physical pages. Identifier: D04-S02.

\par\noindent\texttt{SHA-256: 124faefa1caa1c3f514ad7320193287f\allowbreak 8b42b3fb143e9cd1b05d3e438961f0b2}.

\section*{SPI: Calculus on Manifolds}

Michael Spivak.

1965 edition; 24th printing, January 1995

Local PDF: 159 physical pages. Identifier: D04-S03.

\par\noindent\texttt{SHA-256: fa1b46456524ed708bf73ad0c6762166\allowbreak 3ea0bf9f168010634e7e8d2dcf2f217e}.

\section*{MAR: Vector Calculus}

Jerrold E. Marsden; Anthony Tromba.

Sixth edition; copyright page 2012

Local PDF: 578 physical pages. Identifier: D04-S04.

\par\noindent\texttt{SHA-256: 0cbe934f6ca025b0891a5b8a75498e45\allowbreak c7a206c040bfff5b9e64e1850a9ee1df}.

\section*{Z1: Mathematical Analysis I}

Vladimir A. Zorich.

Second English edition, copyright 2015.

Local PDF: 630 physical pages. Identifier: A01.

\par\noindent\texttt{SHA-256: 6e108f59f382204d892c42792dff1c9c\allowbreak 75ae77a90fd89a08d2af48698bf612b1}.

\section*{Z2: Mathematical Analysis II}

Vladimir A. Zorich.

Second English edition, copyright 2016.

Local PDF: 729 physical pages. Identifier: A02.

\par\noindent\texttt{SHA-256: ecb238c575e5c99bbcb875a1ec97bf4e\allowbreak 95e2e3e3f91c78880eb86881e7ddd61c}.

\section*{RUD: Principles of Mathematical Analysis}

Walter Rudin.

Third edition, 1976, international printing. Readable scan with OCR; check formulas against page images.

Local PDF: 353 physical pages. Identifier: A03.

\par\noindent\texttt{SHA-256: 76db95990342837ed9dbc436fd7035d0\allowbreak ae8df8dbc7614ef65de50d5f3e4cdcc6}.

\section*{APO: Mathematical Analysis}

Tom M. Apostol.

Second edition. Readable scan with OCR; first PDF page is the cover, followed by the preface and contents. Separate title/copyright leaves are not present in this copy. Body starts at printed p. 1 (PDF p. 14), and final index reaches printed p. 492 (PDF p. 505).

Local PDF: 505 physical pages. Identifier: D03-S01.

\par\noindent\texttt{SHA-256: 001553e8bd34d7e2e375a72064d2349e\allowbreak 489f84850505630277c10fcea168882c}.

\section*{PUG: Real Mathematical Analysis}

Charles Chapman Pugh.

Second edition, copyright 2015; copyright page states corrected at second printing 2017. eBook ISBN 978-3-319-17771-7. Body begins at printed p. 1 (PDF p. 12); final index reaches printed p. 478 (PDF p. 486). Printed-to-PDF offsets change later; use verified per-location pages, not one global offset.

Local PDF: 486 physical pages. Identifier: D03-S02.

\par\noindent\texttt{SHA-256: e66b19c61567435f19161ee81f149e70\allowbreak d271279ca606a575261b4fa26eee71cf}.

\section*{Corrections that affect mathematical content}

\begin{enumerate}[leftmargin=*]

\item \textbf{SPI, Addenda to Theorem 3-8.} The Riemann integrability criterion proof must account for boundaries of the partition rectangles, covering them with arbitrarily small total volume. A point outside a bad-oscillation set may still lie on a partition boundary. The proof here treats those boundaries explicitly.

\item \textbf{SPI, Theorem 3-10 and HUB, Theorem 4.5.10.} The upper/lower-integral formulation of Riemann Fubini is different from an iterated ordinary-integral statement with every slice integrable. The latter needs its slice hypothesis. A null exceptional set does not make all slices integrable.

\item \textbf{SPI, Addenda on alternating tensors.} Alternating covectors belong to the exterior algebra of the dual space. Coordinate wedge forms use the normalization fixed in this manuscript; factorial conventions are translated, not silently mixed.

\item \textbf{PUG, Chapter 5 Appendix D.} The heading ``Polar Form concerns polar decomposition of a linear isomorphism. It is not a source for polar-coordinate substitution. Polar, cylindrical and spherical formulas below are derived from the Jacobian theorem with their seams and exceptional sets treated.

\item \textbf{FLE, Chapter 5.} Measure-theoretic parameter-integration formulations are not silently identified with Riemann integrability. The classical statements below have their own uniform estimates and comparison bridges.

\item \textbf{MUN, bounded-set terminology.} A ``rectifiable set in a classical multiple-integration discussion is read against its stated boundary/content definition. It is not an admission of the modern geometric-measure theory of countably rectifiable sets.

\item \textbf{Finite differentiability.} The Morse supplement states the finite regularity supplied by its proof. Exterior differentiation, pullback and Stokes retain separate regularity hypotheses. None is justified by a blanket replacement of finite differentiability with smoothness.

\end{enumerate}

\section*{Coverage beyond the selected core}

Abstract manifolds and full de Rham cohomology are routed to the topology and

differential-geometry directions. Construction of Lebesgue measure belongs to

measure theory. General Banach calculus, PDE and numerical methods remain

separate. The bounded Euclidean examples included here do not claim those courses.

% END ORIGINAL INPUT: source-register.tex


% BEGIN ORIGINAL INPUT: verification-register.tex
\chapter*{Verification and future formalization}
\addcontentsline{toc}{chapter}{Verification and future formalization}
The files accompanying this manuscript contain the mathematical result
inventory, knowledge-family correspondence, inspected foundation interfaces,
source comparison records and independent review findings. A page count or a
count of theorem environments is not a proof of coverage. The coverage check
starts from the 158 prepared families and the source scope, then identifies
the concrete definitions, statements, counterexamples and bridges in the text.

The proof check distinguishes missing hypotheses, invalid arguments, unsupported
regularity, circular references and incorrect signs from harmless changes in
proof route. An independently stated assertion must be proved or explicitly
declared as a shared foundation. The bounded advanced supplements must state
their actual scope rather than imply a stronger general theory.

The accompanying audit files preserve author checks and independent review
separately. Repairs identify the statement and argument changed. PDF compilation
and visual inspection check typesetting and cross-references; they do not
replace mathematical review. No new Lean code or new formalization run is part
of this preparation delivery. Once authorized, each required result still
needs its actual Lean declaration, kernel acceptance, specification and axiom
checks and independent result-fidelity review in the normal agent workflow.

% END ORIGINAL INPUT: verification-register.tex

\end{document}
